Logic

It is a year and a half now since Kant commissioned me to prepare his
Logic for publication, as he expounded it to his listeners in public lectures,
and to transmit it to the public in the form of a compendious manual. For
this purpose I received from him his own manuscript, which he had used
in his lectures, with the expression of special, honorable confidence in me,
that, being acquainted with the principles of his system in general, I would
easily enter into the course of his ideas, that I would not distort or falsify
his thoughts, but rather would present them with the required clarity and
distinctness and at the same time in the appropriate order. Now since, as I
thus accepted the commission and have sought to carry it out as well as I
was able, in conformity with the wish and the expectation of that praisewor-
thy wise man, my most honored teacher and friend, everything that con-
cerns the exposition, the clothing and the execution, the presentation and
the ordering of the thoughts, is in part to be reckoned to my account, I am
naturally obliged to provide an account to the readers of this new Kantian
work. On this matter, then, a few closer explanations here.
Since the year 1765, Professor Kant has based his lectures on logic,
without interruption, on Meier's textbook as guiding thread (Georg Frie-
drich Meier's Excerpts from the Doctrine of Reason, Halle: Gebauer, 1752),
for reasons he explained in the program he published by way of announce-
ment of his lectures in the year 1765.' The copy of the mentioned compen-
dium that he himself used in his lectures, like all the other textbooks he
used for the same purpose, is interleaved with paper; his general remarks
and elucidations, as well as the more special ones that relate in the first
instance to the text of the compendium in its individual sections, are
found partly on the interleaved paper, partly on the empty margin of the
textbook itself. And what has been written by hand here and there in
scattered remarks and elucidations, taken together, constitutes now the
storehouse of materials which Kant built up in his lectures here, which in
part he expanded from time to time through new ideas, and which in part
he again and again revised anew and improved in regard to various individ-
ual materials. Hence it contains at least the essentials of what the famous
commentator on Meier's textbook was accustomed to communicate to his
listeners concerning logic in lectures that were given in a free manner, and
that which he esteemed worthy of writing down.
Now as for what concerns the presentation and ordering of things in
this work, I believed that I would put forth the great man's ideas and
principles most fittingly if, with respect to the economy and the division of
the whole as such, I held myself to his express explanation, according to
which nothing more may be taken up in the proper treatment of logic, and
in particular in its Doärine of Elements, than the theory of the three essen-
tial principal functions of thought: concepts, judgments, and inferences.
Hence everything that deals with cognition in general and with its logical
perfections, and which in Meier's textbook precedes the doctrine of con-
cepts and takes in almost half of the whole, must accordingly be reckoned
to the introduction. "Previously," Kant says right at the outset of the
eighth section, 2 in which his author expounds the doctrine of concepts -
"Previously cognition in general was treated, as propaedeutic to logic, now
logic itself follows."
In consequence of this explicit pointer, I have placed everything that
occurs before the mentioned section in the Introduction, which for this
reason contains a much greater extension than it customarily includes in
other logic manuals. The consequence of this was then also that the Doc-
trine of Method, as the other principal part of the treatise, had to turn out that
much shorter, the more the materials that had already been treated in the
Introduction, as for example the doctrine of proofs, etc. - which, by the
way, are now rightly included in the sphere of the doctrine of method by our
modern logicians. It would have been a repetition as unnecessary as im-
proper to give mention to this material yet again in its proper place, only in
order to make the incomplete complete and to put everything in its proper
place. I have done the latter with respect to the doctrine of definitions,
however, and of the logical division of concepts, which in Meier's compendium
belongs to the eighth section, namely, to the doctrine of elements of con-
cepts, an order which even Kant left unaltered in his exposition.
It is surely obvious that the great reformer of philosophy and, as concerns
the economy and external form of logic, of this part of theoretical philoso-
phy in particular, would have worked on logic according to his architectonic
plan, whose essential outlines are sketched in the Critique of Pure Reason, if
it had pleased him to do so, and if his occupation with a scientific grounding
of the whole system of philosophy proper - philosophy of the really true
and certain, which is an occupation immeasurably more important and
more difficult, and which in the first instance only he and he alone could
carry out with his originality - had permitted him to think of preparing a
logic himself. But this work he was quite able to leave to others, who, with
insight and with unbiased judgment, could use his architectonic ideas for a
truly purposeful and well-ordered arrangement and treatment of this sci-
ence. This was to be expected from several thorough and unbiased thinkers
among our German philosophers. And in this expectation Kant and the
friends of his philosophy were not disappointed. Several recent textbooks
on logic are to be regarded, in respect of the economy and the disposition of
the whole, more or less as fruit of those Kantian ideas on logic. That this
science has really gained through this; that it has become, though not richer
or really more solid as to its content or in itself better grounded, yet more
purified, partly from all its foreign components, partly from so many useless
subtleties and merely dialectical tricks; that it has become more systematic
and yet at the same time, with all scientific strictness of method, simpler - of
this everyone must be convinced, even by the most fleeting comparison of
older textbooks of logic with modern ones worked out in accordance with
Kantian principles, if only he has correct and clear concepts of the proper
character and the legitimate limits of logic. For however much many of the
older manuals on this science may stand out for scientific strictness of
method; for clarity, determinateness, and precision of explanations; and for
conciseness and evidence in proofs; still there is scarcely a one of them in
which the limits of the various spheres that belong to universal logic in its
broader extension - the merely propaedeutic, the dogmatic und technical, the
pure and the empirical - do not run into and through each other, so that the
one cannot be determinately distinguished from the other.
Herr Jakob, in the Preface to the first edition of his logic,3 does observe:
"Wolff grasped the idea of a universal logic exceptionally well, and if this
great man had thought of expounding pure logic in complete separation,
he would certainly, by means of his systematic mind, have presented us
with a masterpiece, which would have made all further work of this kind
unnecessary." But he did not execute this idea, and none of his successors
developed it either, however great and well-grounded in general the merit
may be that the Wolffian school has attained concerning the properly
logical, the formal perfection in our philosophical cognition.
But irrespective of what was able to happen and had to happen in
regard to external form for the perfection of logic by the necessary separa-
tion of pure and merely formal propositions from empirical and real or
metaphysical ones, when it comes to passing judgment on and determin-
ing the inner content of this science as science, there is no doubt about
Kant's judgment on this point. He frequently explained, determinately
and expressly, that logic is to be regarded as a separate science, existing
for itself and grounded in itself, and hence that from its origin and first
development with Aristotle, right down to our times, it could not really gain
anything in scientific grounding. In conformity with this claim, Kant did
not think either about grounding the logical principles of identity and
contradiction on a higher principle, or about deducing the logical forms of
judgment. He recognized and treated the principle of contradiction as a
proposition that has its evidence in itself and requires no derivation from a
higher principle. He only restricted the use, the validity, of this principle
by banishing it from the sphere of metaphysics, where dogmatism sought
to make it valid, and restricting it to the merely logical use of reason, as
valid only for this use alone.
But now whether the logical principle of identity and of contradiction is
really incapable of or does not need any further deduction, in itself and
without qualification, that is of course a different question, which leads to
the highly significant question of whether there is in general an absolutely
first principle of all cognition and science, whether such a thing is possible
and can be found.
The doctrine of science'' believes that it has discovered such a principle in
the pure, absolute I, and hence that it has grounded all philosophical
knowledge perfectly, not merely as to form but also as to content. And
having presupposed the possibility and the apodeictic validity of this abso-
lutely one and unconditioned principle, it then proceeds completely con-
sistently when it does not allow the logical principles of identity and of
contradiction, the propositions^ = A and -A = -A, to hold uncondition-
ally, but instead declares them to be subaltern principles, which can and
must be established and determined only through it and its highest propo-
sition: / am. (See the Foundation of the Doctrine of Science, p. ißf.) In an
equally consistent way Schellingf in his system of transcendental idealism,
declares himself against the presupposition of logical principles as uncondi-
tioned, i.e., as not derivable from any higher principle, since logic can arise
in general only through abstraction from determinate propositions and,
insofar as it arises in a scientific way, only through abstraction from the
highest principles of knowledge, and since it consequently presupposes
these highest principles of knowledge and with them the doctrine of
science itself. Since, however, from the other side, these highest princi-
ples of knowledge, considered as principles, just as necessarily presuppose
logical form, there arises here that circle which cannot be broken by
science, but which can at least be explained, explained by recognition of a
principle of philosophy that is first both as to form and as to content
(formally and materially), in which these two, form and content, recipro-
cally condition and ground one another. In this principle is supposed to lie
the point in which the subjective and the objective, identical knowledge
and synthetic knowledge, are one and the same.
Given the presupposition of such a dignity, which must undoubtedly
belong to such a principle, logic would accordingly have to be subordinated,
like every other science, to the doctrine of science and its principles.
But whatever the situation may be here, this much is settled: In the
inner part of its sphere, as concerns the essential, logic remains in every
case unaltered; and the transcendental question as to whether logical
propositions are still capable of and require a derivation from a higher,
absolute principle has as little influence on logic and on the validity and
evidence of its laws as the transcendental problem, How are synthetic
judgments a priori possible in mathematics? does on pure mathematics in
regard to its scientific content. Like the mathematician as mathematician,
the logician as logician, within the sphere of his science, can also continue
confidently and certainly to explain and to prove, without permitting him-
self to worry about the transcendental question, which lies outside his
sphere, as to horn pure mathematics or pure logic is possible as a science.
Given this universal recognition of the correctness of universal logic,
the battle between the skeptics and the dogmatists concerning the ulti-
mate grounds of philosophical knowledge has never been conducted in
the domain of logic, whose rules were recognized as valid by every rational
skeptic as well as by the dogmatist, but rather has always been conducted
in the sphere of metaphysics. And how could it be otherwise? The highest
task of philosophy proper concerns not subjective but objective, not identi-
cal but synthetic, knowledge. In this, logic as such remains completely on
the sidelines; it could not occur either to critique or to the doctrine of
science — nor will it be able to occur at all to a philosophy that knows how
to distinguish determinately the transcendental standpoint from the
merely logical - to seek die ultimate grounds of real philosophical knowl-
edge inside the sphere of mere logic, and to wish to cull a real object from a
proposition of logic, considered merely as such.
He who has grasped die enormous difference between logic proper
(universal logic), as a merely formal science, the science of mere diought
as diought, and transcendental philosophy, this sole material or real pure
science of reason, the science of knowledge proper, and who never again
fails to heed tiiis difference, will thus easily be able to pass judgment on
the question of what is to be said concerning die recent attempt lately
undertaken by Herr Bardili (in his Outline of Primary Logic* 1 ) to make out a
prius" for logic itself, in die expectation of finding on the path of this
investigation "a real object, eidier posited by it (mere logic) or nodiing can
be posited at all; the key to die essence of nature, either given by it, or no
logic and no philosophy is possible." In trudi, however, one cannot see in
what way Herr Bardili could discover a real object based on die prius of
logic that he advances, the principle of the absolute possibility of thought,
according to which we can repeat one as one and the same in the many (not
die manifold) infinitely many times. This prius of logic, supposedly newly
discovered, is in fact obviously nothing more and nothing less dian die
old, long recognized principle of identity, established within the sphere of
logic and placed at the head of this science: What I think, I think; and it is
only this and nothing else dial I can think repeated to infinity. — Who, after
all, in the case of the well understood logical principle of identity/ will
diink of a manifold and not of a mere many, which of course arises and can
arise through nothing other than the mere repetition of one and the same
thought, the mere repeated positing of an A = A = A, and so on to
infinity. It would be difficult, then, on the path that Herr Bardili has taken,
and by the heuristic method that he has used, to find what philosophical
reason seeks: the beginning and ending point, from which it may set out in
its investigations and to which it may again return. The principal and most
significan: objections that Herr Bardili has lodged against Kant and his
method of philosophizing could not so much touch Kant the logician, then,
but rather Kant the transcendental philosopher and metaphysician.
Here, then, we can let them be deferred entirely to their proper place.
Finally, I wish to mention here that as soon as leisure permits I will
work up the Kantian Metaphysics in the same manner and publish it, for
which I already have the manuscript in hand.

Introduction
i.
Concept of logic
Everything in nature, both in the lifeless and in the living world, takes
place according to rules, although we are not always acquainted with these
rules. - Water falls according to laws of gravity, and with animals locomo-
tion also takes place according to rules. The fish in water, the bird in the
air, move according to rules. The whole of nature in general is really
nothing but a connection of appearances according to rules; and there is
no absence of rules anywhere. If we believe we have found such a thing, then
in this case we can only say that we are not acquainted with the rules.
The exercise of our powers also takes place according to certain rules
that we follow, unconscious of them at first, until we gradually arrive at
cognition of them through experiments and lengthy use of our powers,
indeed, until we finally become so familiar with them that it costs us much
effort to think them in abstracto. Thus universal grammar is the form of a
language in general, for example. One speaks even without being ac-
quainted with grammar, however; and he who speaks without being ac-
quainted with it does actually have a grammar and speaks according to
rules, but ones of which he is not himself conscious.
Like all our powers, the understanding in particular is bound in its
actions to rules, which we can investigate. Indeed, the understanding is
to be regarded in general as the source and the faculty for thinking rules
in general. For as sensibility is the faculty of intuitions, so the under-
standing is the faculty for thinking, i.e., for bringing the representations
of the senses under rules. Hence it is desirous of seeking for rules and is
satisfied when it has found them. Since the understanding is the source
of rules, the question is thus, according to what rules does it itself
proceed?
For there can be no doubt at all: we cannot think, we cannot use our
understanding, except according to certain rules. But now we can in turn
think these rules for themselves, i.e., we can think them apart from their
application or in abstracto. Now what are these rules?
527
13
All rules according to which the understanding operates are either neces-
sary or contingent. The former are those without which no use of the
understanding would be possible at all, the latter those without which a
certain determinate use of the understanding would not occur. The contin-
gent rules, which depend upon a determinate object of cognition, are as
manifold" as these objects themselves. Thus there is, for example, a use of
the understanding in mathematics, in metaphysics, morals, etc. The rules
of this particular, determinate use of the understanding in the sciences
mentioned are contingent, because it is contingent whether I think of this
or that object, to which these particular rules relate.
If now we put aside all cognition that we have to borrow from objects and
merely reflect on the use just of the understanding, we discover those of
its rules which are necessary without qualification, for every purpose and
without regard to any particular objects of thought, because without them
we would not think at all. Thus we can have insight into these rules a
priori, i.e., independent of all experience, because they contain merely the
conditions for the use of the understanding in general, without distinction
among its objects, be that use pure or empirical. And from this it follows at
the same time that the universal and necessary rules of thought in general
can concern merely its form and not in any way its matter. Accordingly, the
science that contains these universal and necessary rules is merely a
science of the form of our cognition through the understanding, or of
thought. And thus we can form for ourselves an idea of the possibility of
such a science, just as we can of a universal grammar, which contains
nothing more than the mere form of language in general, without words,
which belong to the matter of language.
Now this science of the necessary laws of the understanding and of
reason in general, or what is one and the same, of the mere form of
thought as such, we call logic.
As a science that deals with all thought in general, without regard to
objects as the matter of thought, logic
1. is to be regarded as foundation for all the other sciences and as the
propaedeutic to all use of the understanding. Just because it does abstract
wholly from all objects, however, it also
2. cannot be an organon of the sciences.
By an organon we understand, namely, a directive as to how a certain
cognition is to be brought about. This requires, however, that I already be
acquainted with the object of the cognition that is to be produced accord-
ing to certain rules. An organon of the sciences is thus not mere logic,
°
"vielfältig."
because an organon presupposes exact acquaintance with the sciences,
their objects and sources. Thus mathematics, for example, as a science
that contains the ground for the extension of our cognition in regard to a
certain use of reason, is an excellent organon. Logic, on the other hand, as
universal propaedeutic to all use of the understanding and of reason in
general, may not go into the sciences and anticipate their matter. It is only
a universal an of reason (canonica Epicuri) for making cognitions in general
conform to the form of die understanding in general, and hence is only to
this extent to be called an organon, which serves of course merely for
passing judgment and for correcting our cognition, but not for expanding it.
3. As a science of the necessary laws of thought, without which no use
of the understanding or of reason takes place at all, laws which conse-
quently are conditions under which the understanding can and ought to
agree with itself alone - the necessary laws and conditions of its correct
use — logic is, however, a canon. And as a canon of the understanding and
of reason it may not borrow any principles either from any science or from
any experience; it must contain nothing but laws a priori, which are neces-
sary and have to do with the understanding in general.
Some logicians, to be sure, do presuppose psychological principles in
logic. But to bring such principles into logic is just as absurd as to derive
morals from life. If we were to take principles from psychology, i.e., from
observations concerning our understanding, we would merely see how
thinking does take place and how it is under various subjective obstacles
and conditions; this would lead dien to cognition of merely contingent laws.
In logic, however, die question is not about contingent but about necessary
rules; not how we do think, but how we ought to think. The rules of logic
must thus be derived not from the contingent but from the necessary use of
die understanding, which one finds in oneself apart from all psychology.
In logic we do not want to know how die understanding is and does think
and how it has previously proceeded in diought, but radier how it ought to
proceed in diought. Logic is to teach us die correct use of die understand-
ing, i.e., dial in which it agrees with itself.
From die given explanation of logic, die remaining essential properties of
this science may now also be derived, namely, that it
4. is a science of reason not only as to its mere form but as to its matter^
since its rules are not derived from experience, and since at the same time
it has reason as its object. Logic is thus a self-cognition of the understand-
ing and of reason, not as to their faculties in regard to objects, however,
but merely as to form. In logic I will not ask what the understanding
cognizes and how much it can cognize or how far its cognition goes. For
diat would be self-cognition in regard to its material use and tiius belongs
to metaphysics. In logic the question is only, How will the understanding
cognize itself?

Finally, as a science that is rational as to matter and to form, logic is also
5. a doctrine or a demonstrated theory. For since it is occupied, not with
the common and as such merely empirical use of the understanding and
of reason, but rather merely with the universal and necessary laws of
thought in general, it rests on principles a priori, from which all its rules
can be derived and proved, as ones with which all cognition of reason has
to be in conformity.
By virtue of the fact that logic is to be taken as a science a priori, or as
a doctrine for a canon of the use of the understanding and of reason, it
is essentially distinct from aesthetics, which as mere critique of taste has no
canon (law) but only a norm (model or standard for passing judgment),
which consists in universal agreement. Aesthetics, that is, contains the
rules for the agreement of cognition with the laws of sensibility; logic, on
the other hand, contains the rules for the agreement of cognition with
the laws of the understanding and of reason. The former has only
empirical principles and thus can never be science or doctrine, provided
that one understands by doctrine a dogmatic instruction from principles
a priori, in which one has insight into everything through the understand-
ing without instruction from other quarters attained from experience,
and which gives us rules, by following which we procure the required
perfection.
Some, especially orators and poets, have tried to engage in reasoning*
concerning taste, but they have never been able to hand down a decisive
judgment concerning it. The philosopher Baumgarten in Frankfurt had a
plan for an aesthetic as a science. 8 But Home, 1 * more correctly, called
aesthetics critique, since it yields no rules a priori that determine judgment
sufficiently, as logic does, but instead derives its rules a posteriori, and
since it only makes more universal, through comparison, the empirical
laws according to which we cognize the more perfect (beautiful) and the
more imperfect.
Logic is thus more than mere critique; it is a canon that subsequently
serves for critique, i.e., as the principle for passing judgment on all use of
the understanding in general, although only on its correctness in regard to
mere form, since it is not an organon, any more than universal grammar
is.
Now as propaedeutic to all use of the understanding in general, univer-
sal logic is distinct also on another side from transcendental logic, in which
the object itself is represented as an object of the mere understanding;
universal logic, on the contrary, deals with all objects in general.
If we now join together all the essential marks that belong to a complete
'
"vernünfteln."
determination of the concept of logic, then we shall have to put forth the
following concept of it:
Logic is a science of reason, not as to mere form but also as to matter; a
science a priori of the necessary laws of thought, not in regard to particular
objects, however, but to all objects in general; — hence a science of the correct use
of the understanding and of reason in general, not subjectively,
however, i.e.,
not according to empirical (psychological) principles for how the understanding
does think, but objectively, i.e., according to principles a priori for how it ought
to think.
II.
Principal divisions of logic - Exposition — Use of this science — Sketch
of its history
Logic is divided
i.
into analytic and dialectic.
Analytic discovers through analysis all the actions of reason that we
perform in thinking. It is thus an analytic of the form of the understanding
and of reason and is rightly called the logic of truth, because it contains
the necessary rules of all (formal) truth, apart from which our cognition is
untrue in itself, regardless of its objects. Thus it is also nothing more than
a canon for adjudication (of the formal correctness of our cognition).
If one were to use this merely theoretical and universal doctrine as a
practical art, i.e., as an organon, then it would become dialectic. A logic of
illusion (ars sophistica, disputatoria'), which arises out of a mere misuse of
analytic, insofar as the illusion of a true cognition, the marks of which
have to be derived from agreement with objects and thus from content, is
fabricated according to mere logical form.
In earlier times dialectic was studied with great industry. This art ex-
pounded false principles under the illusion of truth and then sought, in
conformity with these, to maintain things in accordance with illusion.
Among the Greeks the dialecticians were the lawyers'' and orators, who
were able to lead the people wherever they wanted, because the people
allow themselves to be misled by illusion. At that time dialectic was thus
the art of illusion. For a long time, too, it was expounded in logic under
the name of the an of disputation, and as long as it was, all of logic and
philosophy were the cultivation of certain garrulous souls for fabricating
any illusion. Nothing can be less worthy of a philosopher, however, than
the cultivation of such an art. In this sense it must be done away with,
'
'
a sophistic art, an art of disputation.
"Sachwalter."
531
then, and instead of this a critique of this illusion must be introduced into
logic.
We would have two parts of logic, accordingly: analytic, which would
expound the formal criteria of truth, and dialectic, which would contain the
marks and rules in accordance with which we could recognize that some-
thing does not agree with the formal criteria of truth, although it seems to
agree with them. Dialectic in this sense would thus have its good use as
cathartic of the understanding.
It is customary to divide logic further into
2.
natural or popular logic and artificial or scientific logic (logica naturalis, log.
scholastica s. artificialis).
But this division is inadmissible. For natural logic, or logic of common
reason (sensus communis), is not really logic but an anthropological science
that has only empirical principles, in that it deals with the rules of the
natural use of the understanding and of reason, which are cognized only
in concrete, hence without consciousness of them in abstracto. - Only artifi-
cial or scientific logic deserves this name, then, as a science of the neces-
sary and universal rules of thought, which can and must be cognized a
priori, independently of the natural use of the understanding and of rea-
son in concreto, although these rules can first be found only through obser-
vation of that natural use.
3.
18
Another division of logic is that into theoretical and praaical logic. But this
division is also incorrect.
Universal logic, which as a mere canon abstracts from all objects, can-
not have a practical part. This would be a contmdictio in adjecto, because a
practical logic presupposes acquaintance with a certain kind of object, to
which it is applied. Thus we can call every science a practical logic; for in
each we must have a form of thought. Universal logic, considered as
practical, can thus be nothing more than a technique of learnedness in
general, an organon of scholastic method.
In consequence of this division logic would thus have a dogmatic and a
technical part. The first would be called the doctrine of elements, the second
the doctrine of method. The practical or technical part of logic would be a
logical art in regard to order and to logical terms of art and logical
distinctions, to make it easier for the understanding to act.
In both parts, both the technical and the dogmatic, it would be
impermissible to give the least consideration either to the objects or to the
subject of thought. In this latter relation logic could be divided
4.
into pure and applied logic.
In pure logic we separate the understanding from the other powers of
the mind and consider what it does by itself alone. Applied logic considers
the understanding insofar as it is mixed with the other powers of the mind,
which influence its actions and misdirect it, so that it does not proceed in
accordance with the laws which it quite well sees to be correct. Applied
logic really ought not to be called logic. It is a psychology in which we
consider how things customarily go on in our thought, not how they ought
to go on. In the end it admittedly says what one ought to do in order to
make correct use of the understanding under various subjective obstacles
and restrictions; and we can also learn from it what furthers the correct
use of the understanding, the means of aiding it, or the cures for logical
mistakes and errors. But propaedeutic it simply is not. For psychology,
from which everything in applied logic must be taken, is a part of the
philosophical sciences, to which logic ought to be the propaedeutic.
It is said, to be sure, that technique, or the way of building a science,
ought to be expounded in applied logic. But that is futile, indeed, even
harmful. One then begins to build before one has materials, and one gives
form, but content is lacking. Technique must be expounded within each
science.
Finally, as for what concerns
5.
the division of logic into the logic of the common and that of the speculative
understanding, we note that this science simply cannot be thus divided.
It cannot be a science of the speculative understanding. For as a logic of
speculative cognition or of the speculative use of reason it would be an
organon for other sciences and not a mere propaedeutic, which ought to
deal with all possible use of the understanding and of reason.
Just as little can logic be a product of the common understanding. The
common understanding is the faculty by which we have insight into the
rules of cognition in concreto. Logic, however, ought to be a science of the
rules of thought in abstraäo.
Nonetheless we can accept the universal human understanding as ob-
ject of logic, and to this extent it will abstract from the particular rules of
speculative reason and thus be distinct from the logic of the speculative
understanding.
As for what concerns the exposition of logic, it can be either scholastic or
popular.
It is scholastic insofar as it is adequate to the curiosity, the capabilities,
and the culture of those who want to treat the cognition of logical rules as
a science. Popular, however, if it condescends to the capabilities and needs
of those who do not study logic as science but only want to use it to
enlighten their understanding. - In the scholastic exposition the rules
must be presented in their universality or in abstracto; in the popular, on the
other hand, in the particular or in concreto. The scholastic exposition is the
533
20
foundation for the popular, for the only one who can expound something
in a popular way is one who could also expound it more thoroughly.
Here we distinguish exposition from method, by the way. By method is to
be understood, namely, the way to cognize completely a certain object, to
whose cognition the method is to be applied. It has to be derived from the
nature of the science itself and, as an order of thought that is determined
thereby and is necessary, it cannot be altered. Exposition means only the
manner of communicating one's thoughts in order to make a doctrine
understandable.
From what we have previously said concerning the nature' and the end of
logic, the worth of this science and the use of its study can be evaluated in
accordance with a correct and determinate standard.
Logic is thus not a universal art of discovery, to be sure, and not an
organon of truth - not an algebra, with whose help hidden truths can be
discovered.
It is useful and indispensable as a critique of cognition, however, or for
passing judgment on common as well as on speculative reason, not in
order to teach it, but only to make it correct and in agreement with itself.
For the logical principle of truth is agreement of the understanding with
its own universal laws.
As for what concerns the history of logic, finally, we want to cite only the
following:
Contemporary logic derives from Aristotle's Analytic. This philosopher
can be regarded as the father of logic. He expounded it as organon and
divided it into analytic and dialectic. His manner of teaching is very scholas-
tic and has to do with the development of the most universal concepts,
which lie at the basis of logic, but one has no use for it because almost
everything amounts to mere subtleties, except that one [has] drawn from
this the names for various acts of the understanding.
From Aristotle's time on, logic has not gained much in content, by the
way, nor can it by its nature do so. But it can surely gain in regard to
exactness, determinateness, and distinctness. There are only a few sciences
that can attain a permanent condition, where they are not altered any
more. These include logic and also metaphysics. Aristotle had not omitted
any moment of the understanding; we are only more exact, methodical,
and orderly in this.
It was believed of Lambert's Organon 10 that it would augment logic
considerably. But it contains nothing more except for subtler divisions,
which, like all correct subtleties, sharpen the understanding, of course,
but are of no essential use.
Among modern philosophers there are two who have set universal logic
in motion: Leibniz and Wolff.
Malebranche and Locke did not treat of real logic, since they also deal
with the content of cognition and with the origin of concepts.
The universal logic of Wolff" is the best we have. Some have combined
it with the Aristotelian logic, like Reusch, 1 * for example.
Baumgarten, a man who has much merit here, concentrated the Wolff-
ian logic,'3 und Meier then commented again on Baumgarten.
Crusius 1 * also belongs to the modern logicians, but he did not consider
how things stand with logic. For his logic contains metaphysical principles
and so to this extent oversteps the limits of this science; besides, it puts
forth a criterion of truth that cannot be a criterion, and hence to this
extent gives free reign to all sorts of fantastic notions.
In present times there has not been any famous logician, and we do not
need any new inventions for logic, either, because it contains merely the
form of thought.
III.
Concept of philosophy in general - Philosophy considered according to
the scholastic concept and according to the worldly concept - Essential
requirements and ends of philosophizing - The most universal and
highest tasks of this science
It is sometimes hard to explain what is understood by a science. But the
science gains in precision through establishment of its determinate con-
cept, and in this way many mistakes are avoided which otherwise creep in,
for certain reasons, if one cannot yet distinguish the science from sciences
related to it.
Before we try to give a definition of philosophy, however, we must first 22
investigate the character of various cognitions themselves, and since philo-
sophical cognitions belong to the cognitions of reason, we must explain in
particular what is to be understood by the latter.
Cognitions of reason are opposed to historical cognitions. The former
are cognitions from principles (ex principiis), the latter cognitions from data
(ex datis). — A cognition can have arisen from reason and in spite ofthat be
historical, however, as when a mere literator learns the products of some-
one else's reason his cognition of these products of reason is then merely
historical, for example.
One can distinguish cognitions, then,
1.
2.
according to their objective origin, i.e., according to the sources from which
alone a cognition is possible. In this respect all cognitions are either rational or
empirical;
according to their subjective origin, i.e., according to the way in which a
cognition can be acquired by men. Considered from this latter viewpoint,
cognitions are either rational or historical, however they may have arisen in
themselves. Hence something that is subjectively only historical can be objec-
tively a cognition of reason.
With some rational cognitions it is harmful to know them merely histori-
cally, while with others it makes no difference. Thus the sailor knows the
rules of navigation historically from his tables, for example, and that is
enough for him. But if the jurist knows jurisprudence merely historically,
then he is fully ruined as a genuine judge, and still more so as a legislator.
From the stated distinction between objectively and subjectively rational
cognitions it is also clear now that one can in a certain respect learn
philosophy without being able to philosophize. He who really wants to
become a philosopher must practice making a free use of his reason, then,
and not a merely imitative and, so to speak, mechanical use.
23
We have explained cognitions of reason as cognitions from principles, and
from this it follows that they must be a priori. But there are two kinds of
cognitions, which are both a priori, but which nevertheless have many
noteworthy differences, namely, mathematics and philosophy.
It is customary to maintain that mathematics and philosophy are distinct
from one another as to their object, in that the former deals with quantity,
the latter with quality. But^ this is wrong. The distinction between these
sciences cannot rest on the object, for philosophy deals with everything,
hence also with quanta, and mathematics does so in part too, insofar as
everything has a quantity. The specific difference between these two
sciences is constituted only by the different kind of cognition of reason, or of
the use of reason, in mathematics and philosophy. Philosophy is, namely,
cognition of reason from mere concepts, while mathematics is cognition of reason
from the construction of concepts.
We construct concepts when we exhibit them in intuition a priori without
experience, or when we exhibit in intuition the object that corresponds to
our concept of it. - The mathematician can never make use of his reason
in accordance with mere concepts, the philosopher never through con-
struction of concepts. In mathematics one uses reason in concreto, but the
intuition is not empirical; rather, here one makes something the object of
intuition for himself a priori.
f
Reading "Allein" for "Alles," in accordance with the published list of printer's errors
(KI, xli).
And as we see, mathematics has an advantage over philosophy here, in
that the cognitions of the former are intuitive cognitions while those of the
latter are only discursive. The cause of the fact that in mathematics we
consider quantities more lies in this, that quantities can be constructed a
priori in intuition, while qualities on the other hand cannot be exhibited in
intuition.
Philosophy is thus the system of philosophical cognitions or of cognitions
of reason from concepts. That is the scholastic concept g of this science.
According to the worldly concept 11 it is the science of the final ends of
human reason. This high concept gives philosophy dignity, i.e., an abso-
lute worth. And actually it is philosophy, too, which alone has only inner
worth, and which first gives a worth to all other cognitions.
Yet in the end people always ask what purpose is served by philosophiz-
ing and by its final end[,] philosophy itself considered as science in accor-
dance with the scholastic concept.
In this scholastic sense of the word, philosophy has to do only with skill,
but in relation to the worldly concept, on the other hand, with usefulness.
In the former respect it is thus a doctrine of skill; in the latter, a doctrine of
n>isdom[,] the legislator of reason[,] and the philosopher to this extent not
an artist of reason' but rather a legislator.
The artist of reason, or the philodox, as Socrates calls him, strives only
for speculative knowledge, without looking to see how much the knowl-
edge contributes to the final end of human reason; he gives rules for the
use of reason for any sort of end one wishes. The practical philosopher,
the teacher of wisdom through doctrine and example, is the real philoso-
pher. For philosophy is the idea of a perfect wisdom, which shows us the
final ends of human reason.
According to the scholastic concept, philosophy involves two things:
First, a sufficient supply of cognitions of reason, andrer the second thing,
a systematic connection of these cognitions, or a combination of them in
the idea of a whole.
Not only does philosophy allow such strictly systematic connection, it is
even the only science that has systematic connection in the most proper
sense, and it gives systematic unity to all other sciences.
As for what concerns philosophy according to the worldly concept (in
sensu cosmico), we can also call it a science of the highest maxim for the use of
our reason, insofar as we understand by a maxim the inner principle of
choice among various ends.
For philosophy in the latter sense is in fact the science of the relation of
all cognition and of all use of reason to the ultimate end of human reason,
to which, as the highest, all other ends are subordinated, and in which
they must all unite to form a unity.
The field of philosophy in this cosmopolitan sense can be brought
down to the following questions:
1.
2.
3.
4.
What
What
What
What
can I know?
ought I to do?
may I hope?
is man?
Metaphysics answers the first question, morals the second, religion the
third, and anthropology the fourth. Fundamentally, however, we could
reckon all of this as anthropology, because the first three questions relate
to the last one.
The philosopher must thus be able to determine
1.
2.
3.
26
the sources of human knowledge,
the extent of the possible and profitable use of all knowledge, and finally
the limits of reason.
The last is the most necessary but also the hardest, yet the philodox
does not bother himself about it.
To a philosopher two things chiefly pertain: i) Cultivation of talent and
of skill, in order to use them for all sorts of ends. 2) Accomplishment in
the use of all means toward any end desired. The two must be united; for
without cognitions-' one will never become a philosopher, but cognitions*
alone will never constitute the philosopher either, unless there is in addi-
tion a purposive combination of all cognitions and skills in a unity, and an
insight into their agreement with the highest ends of human reason.
No one at all can call himself a philosopher who cannot philosophize.
Philosophizing can be learned, however, only through practice and
through one's own use of reason.
How should it be possible to learn philosophy anyway? Every philosophi-
cal thinker builds his own work, so to speak, on someone else's ruins, but
no work has ever come to be that was to be lasting in all its parts. Hence
one cannot learn philosophy, then, just because it is not yet given. But even
granted that there were a philosophy actually at hand, no one who learned it
would be able to say that he was a philosopher, for subjectively his cogni-
tion' of it would always be only historical.
In mathematics things are different. To a certain extent one can proba-
become convinced of them; and on account of its evidence it can also, as it
were, be preserved as a certain and lasting doctrine.
He who wants to learn to philosophize, on the other hand, may regard
all systems of philosophy only as history of the use of reason and as objects
for the exercise of his philosophical talent.
Thus the true philosopher, as one who thinks for himself, must there-
fore make a free use of his reason on his own, not a slavishly imitative use.
But not a dialectical use, i.e., not one that aims only at giving cognitions the
illusion of truth and wisdom. This is the business of the mere sophist,
thoroughly incompatible with the dignity of the philosopher, as one who is
acquainted with and is a teacher of wisdom.
For science has an inner, true worth only as organ of wisdom. As such,
however, it is also indispensable for it, so that one may well maintain that
wisdom without science is a silhouette of a perfection to which we shall
never attain.
He who hates science but loves wisdom all the more is called a
misologist. Misology arises commonly out of an emptiness of scientific
cognitions'" and a certain vanity bound up with that. Sometimes, however,
people who had initially pursued sciences with great industry and fortune,
but who found in the end no satisfaction in the whole of their knowledge,
also fall into the mistake of misology.
Philosophy is the only science that knows how to provide for us this
inner satisfaction, for it closes, as it were, the scientific circle, and only
through it do the the sciences attain order and connection.
For the sake of practice in thinking for ourselves, or philosophizing, we
will have to look more to the method for the use of our understanding than
to the propositions themselves at which we have arrived through this
method.
IV.
Short sketch of a history of philosophy
There is some difficulty in determining the limits where the common use
of the understanding ends and the speculative begins, or where common
cognition of reason becomes philosophy.
Nevertheless there is a rather certain distinguishing mark here, namely,
the following:
Cognition of the universal in abstracto is speculative cognition, cognition
of the universal in concrete is common cognition. Philosophical cognition is
speculative cognition of reason, and thus it begins where the common use
of reason starts to make attempts at cognition of the universal in abstracto.
From this determination of the distinction between common and specu-
lative use of reason we can now pass judgment on the question, with
which people we must date the beginning of philosophizing. Among all
peoples, then, the Greeks first began to philosophize. For they first at-
tempted to cultivate cognitions of reason, not with images as the guiding
thread, but in abstracto, while other peoples always sought to make con-
cepts understandable only through images in concrete. Even today there are
peoples, like the Chinese and some Indians, who admittedly deal with
things that are derived merely from reason, like God, the immortality of
the soul, etc., but who nonetheless do not seek to investigate the nature of
these things in accordance with concepts and rules in abstracto. They
make no separation here between the use of the understanding in concrete
and that in abstracto. Among the Persians and the Arabs there is admittedly
some speculative use of reason, but the rules for this they borrowed from
Aristotle, hence from the Greeks. In Zoroaster's Zend-Avesta we find not the
slightest trace of philosophy. The same holds also for the prized Egyptian
wisdom, which in comparison with Greek philosophy was mere child's
Play-
As in philosophy, so too in regard to mathematics, the Greeks were the
first to cultivate this part of the cognition of reason in accordance with a
speculative, scientific method, by demonstrating every theorem from
elements.
When and where the philosophical spirit first arose among the Greeks,
however, one cannot really determine.
The first to introduce the speculative use of reason, and the one from
whom we derived the first steps of the human understanding toward
scientific culture, is Thales, the founder of the Ionian sect. He bore the
surname physicist, although he was also a mathematician, just as in general
mathematics has always preceded philosophy.
The first philosophers clothed everything in images, by the way. For
poetry, which is nothing other than a clothing of thoughts in images, is
older than prose. Thus in the beginning one had to make use of the
language of images and of poetic style even with things that are merely
objects of pure reason. Pherecydes is supposed to have been the first author
of prose.
The lonians were followed by the Eleatics. The principle of the Eleatic
philosophy and of its founder, Xenophanes, was: In the senses there is decep-
tion and illusion, the source of truth lies only in the understanding alone.
Among the philosophers of this school, Zeno distinguished himself as a
man of great understanding and acuity and as a subtle dialectician.
In the beginning dialectic meant the art of the pure use of the under-
standing in regard to abstract concepts separated from all sensibility.
Thus the many encomia of this art among the ancients. Subsequently
these philosophers, who completely rejected the testimony of the senses,
necessarily fell, given their claim, into many subtleties, and thus dialectic
degenerated into die art of maintaining and of disputing any proposition.
And so it became a mere exercise for the sophists, who wanted to engage in
reasoning" about everything, and who devoted diemselves to giving illu-
sion the veneer of trudi and to making black white. On account of this the
name sophist, by which one formerly meant a man who was able to speak
about all things rationally and with insight, became so hated and contempt-
ible, and the name philosopher was introduced instead.
Around the time of the Ionian school there appeared in Magna Graecia a
man of strange genius, who not only founded a school but also outlined
and brought into being a project, the like of which had never been before.
This man was Pythagoras, born on" Samos. He founded, namely, a society
of philosophers who were united with one another into a federation
through the law of silence. His divided his hearers into two classes: the
acusmatics (äxouonaÖixoi), who had simply to listen, and the acroamatics
(ccxooauctoixoi), who were permitted to ask too.
Among his doctrines there were some exoteric ones, which he ex-
pounded to die whole of the people; the remaining ones were secret and
esoteric, determined only for the members of his federation, some of whom
he took into his trusted friendship, separating diem wholly from the
others. He made physics and theology, hence the doctrines of the visible
and the invisible, the vehicle of his secret doctrines. He also had various
symbols, which presumably were nothing odier than certain signs that
allowed the Pytiiagoreans to communicate widi one another.
The end of his federation seems to have been none other than to purify
religion of the delusions of the people, to moderate tyranny, and to introduce more
lawfulness into states. This federation, which die tyrants began to fear, was
destroyed shortly before Pythagoras's death, however, and this philosophi-
cal society was broken up, partly by execution, partly by the flight and the
banning of the greatest part of its members. The few who remained were
novices. And since these knew little of die doctrines peculiar to Pythagoras,
we can say notiiing certain and determinate about them. Subsequendy
many doctrines dial were certainly only invented were attributed to
Pythagoras, who by die way was also a very madiematical mind.
The most important epoch of Greek philosophy starts finally with Socrates.
For it was he who gave to die philosophical spirit and to all speculative
minds a wholly new practical direction. Among all men, too, he was almost
the only one whose behavior comes closest to the idea of a wise man.
Among his disciples the most famous is Plato, who occupied himself
more with Socrates' practical doctrines, and among Plato's disciples Aris-
totle, who in turn raised speculative philosophy higher.
Plato and Aristotle were followed by the Epicureans and the Stoics, who
were openly declared mutual enemies. The former placed the highest good
in a joyful heart, which they called pleasure-/ the latter found it solely in
loftiness and strength of soul, whereby one can do without all the comforts of
life.
The Stoics, by the way, were dialectical in speculative philosophy, dog-
matic in moral philosophy, and in their practical principles, through which
they sowed the seed for the most sublime sentiments* that ever existed,
they showed uncommonly much dignity. The founder of the Stoic school
is Zeno of Citium. The most famous men from this school among the
Greek philosophers are Cleanthes and Chrysippus.
The Epicurean school was never able to achieve the same repute that
the Stoic did. Whatever one may say of the Epicureans, however, this
much is certain: they demonstrated the greatest moderation in enjoyment
and were the best natural philosophers among all the thinkers of Greece.
We note here further that the foremost Greek schools bore particular
names. Thus Plato's school was called the Academy, Aristotle's the Ly-
ceum, the Stoics' school porticus (oioot), a covered walkway, from which
the name Stoic is derived; Epicurus's school was called horti, because
Epicurus taught in gardens.
Plato's Academy was followed by three other Academies, which were
founded by his disciples. Speusippus founded the first, Arcesilaus the sec-
ond, Carneades the third.
These Academies inclined toward skepticism. Speusippus and Arcesilaus
both adjusted their mode of thought to skepticism, and in this Carneades
went still further. On this account the skeptics, these subtle, dialectical
thinkers, are also called Academics. Thus the Academics followed the first
great doubter, Pyrrho, and his successors. Their teacher Plato had himself
given them occasion for this by expounding many of his doctrines dialogi-
cally, so that the grounds pro and contra were put forth, without his decid-
ing about the matter himself, although he was otherwise very dogmatic.
If we begin the epoch of skepticism with Pyrrho, then we get a whole
school of skeptics, who are essentially distinct from the dogmatists in their
mode of thought and method of philosophizing, in that they made it the
first maxim for all philosophizing use of reason to withhold one's judgment
even when the semblance' of truth is greatest; and they advanced the principle
that philosophy consists in the equilibrium of judgment and teaches us to uncover
false semblance. 1 From these skeptics nothing has remained for us, however,
but the two works of Sextus Empiricus, in which he brought together all
doubts.
When philosophy subsequently passed from the Greeks to the Romans, it
was not extended; for the Romans always remained just disciples.
Cicero was a disciple of Plato in speculative philosophy, a Stoic in
morals. The Stoic sect included as die most famous Epictetus, Antonius the
Philosopher, and Seneca. There were no naturalists among the Romans
except for Pliny the Elder^ who left a natural history.
Finally culture disappeared among the Romans too, and barbarism
arose until the Arabs began in die 6th and yth centuries to apply them-
selves to the sciences and to revive Aristotle again. Then the sciences rose
in the Occident again, and in particular the regard for Aristotle, who was
followed, however, in a slavish way. In die nth and i2th centuries the
scholastics appeared; they elucidated Aristode and pursued his subtleties to
infinity. They occupied themselves with nothing but abstractions. This
scholastic mediod of pseudo-philosophizing was pushed aside at die time
of the Reformation, and now there were eclectics in philosophy, i.e., think-
ers who thought for themselves, who acknowledged no school, but who
instead sought the trudi and accepted it where they found it.
Philosophy owes its improvement in modern times partly to the greater
study of nature, partly to die combination of madiematics with natural
science. The order diat arose in thought through the study of these
sciences was also extended over die particular branches and parts of
philosophy proper. The first and greatest investigator of nature in modern
time was Bacon ofVerulam. In his investigations he followed the path of
experience and called attention to the importance and indispensability of
observations and experiments for the discovery of truth. It is hard to say, by
the way, from whence the improvement of speculative philosophy really
comes. Descartes rendered it no small service, in that he contributed much
to giving distinctness to thought by advancing his criterion of truth, which he
placed in the clarity and evidence of cognition.
Leibniz and Locke are to be reckoned among the greatest and most
meritorious reformers of philosophy in our times. The latter sought to
analyze die human understanding and to show which powers of die soul
and which of its operations belonged to this or that cognition. But he did
not complete the work of his investigation, and also his procedure is very
dogmatic, although we did gain from him, in that we began to study the
nature of the soul better and more thoroughly.
As for what concerns the special dogmatic method of philosophizing
peculiar to Leibniz and Wolff, it was quite mistaken. Also, there is so much
in it that is deceptive that it is in fact necessary to suspend the whole
procedure and instead to set in motion another, the method of critical
philosophizing, which consists in investigating the procedure of reason
itself, in analyzing the whole human faculty of cognition and examining
how far its limits may go.
In our age natural philosophy is in the most flourishing condition, and
among the investigators of nature there are great names, e.g., Newton.
Modern philosophers cannot now be called excellent and lasting, because
everything here goes forward, as it were, in flux. What one builds the
other tears down.
In moral philosophy we have not come further than the ancients. As for
what concerns metaphysics, however, it seems as if we had been stopped
short in the investigation of metaphysical truths. A kind of indifferentem
toward this science now appears, since it seems to be taken as an honor to
speak of metaphysical investigations contemptuously as mere cavilling.'
And yet metaphysics is the real, true philosophy!
Our age is the age of critique, and it has to be seen what will come of the
critical attempts of our time in respect to philosophy and in particular to
metaphysics.
V.
Cognition in general — Intuitive and discursive cognition; intuition
and concept and in particular their difference — Logical and aesthetic
perfection of cognition
All our cognition has a twofold relation, yzrtf a relation to the object, second a
relation to the subject. In the former respect it is related to representation,
in the latter to consciousness, the universal condition of all cognition in
general. - (Consciousness is really a representation that another represen-
tation is in me.)
In every cognition we must distinguish matter, i.e., the object, and form,
i.e., the way in which we cognize the object. If a savage sees a house from a
distance, for example, with whose use he is not acquainted, he admittedly
has before him in his representation the very same object as someone else
who is acquainted with it determinately as a dwelling established for men.
But as to form, this cognition of one and the same object is different in the
"Grübeleien."
two. With the one it is mere intuition, with the other it is intuition and
concept at the same time.
The difference in the form of the cognition rests on a condition that
accompanies all cognition, on consciousness. If I am conscious of the repre-
sentation, it is clear, if I am not conscious of it, obscure.
Since consciousness is the essential condition of all logical form of
cognitions, logic can and may occupy itself only with clear but not with
obscure representations. In logic we do not see how representations arise,
but merely how they agree with logical form. In general logic cannot deal
at all with mere representations and their possibility either. This it leaves
to metaphysics. Logic itself is occupied merely with the rules of thought in
concepts, judgments, and inferences, as that through which all thought
takes place. Something precedes, of course, before a representation be-
comes a concept. We will indicate that in its place, too. But we will not
investigate how representations arise. Logic deals with cognition too, to be
sure, because in cognition there is already thought. But representation is
not yet cognition, rather, cognition always presupposes representation.
And this latter cannot be explained at all. For we would always have to
explain what representation is by means of yet another representation.
All clear representations, to which alone logical rules can be applied,
can now be distinguished in regard to distinctness and indistinctness. If we
are conscious of the whole representation, but not of the manifold that is
contained in it, then the representation is indistinct. First, to elucidate
this, an example in intuition.
We glimpse a country house in the distance. If we are conscious that the
intuited object is a house, then we must necessarily have a representation of
the various parts of this house, the windows, doors, etc. For if we did not see
the parts, we would not see the house itself either. But we are not conscious
of this representation of the manifold of its parts, and our representation of
the object indicated is thus itself an indistinct representation.
If we want an example of indistinctness in concepts, furthermore, then
the concept of beauty may serve. Everyone has a clear concept of beauty.
But in this concept many different marks occur, among others that the
beautiful must be something that (i.) strikes the senses and (2.) pleases
universally. Now if we cannot explicate the manifold of these and other
marks of the beautiful, then our concept of it is still indistinct.
Wolff's disciples call the indistinct representation a confused one. But
this expression is not fitting, because the opposite of confusion is not
distinctness but order. Distinctness is an effect of order, to be sure, and
indistinctness an effect of confusion; and every confused cognition is thus
also an indistinct one. But the proposition does not hold conversely; not
every indistinct cognition is a confused one. For in the case of cognitions
in which there is no manifold at hand, there is no order, but also no
confusion.
This is the situation with all simple representations, which never be-
come distinct, not because there is confusion in them, but rather because
there is no manifold to be found in them. One must call them indistinct,
therefore, but not confused.
And even with compound representations, too, in which a manifold of
marks can be distinguished, indistinctness often derives not from confu-
sion but from weakness of consciousness. Thus something can be distinct as
to form, i.e., I can be conscious of the manifold in the representation, but
the distinctness can diminish as to matter if the degree of consciousness
becomes smaller, although all the order is there. This is the case with
abstract representations.
Distinctness itself can be of two sorts:
First, sensible. This consists in the consciousness of the manifold in
intuition. I see the Milky Way as a whitish streak, for example; the light
rays from the individual stars located in it must necessarily have entered
my eye. But the representation of this was merely clear, and it becomes
distinct only through the telescope, because then I glimpse the individual
stars contained in the Milky Way.
Secondly, intellectual; distinctness in concepts or distinctness of the understand-
ing. This rests on the analysis of the concept in regard to the manifold that
lies contained within it. Thus in the concept of virtue, for example, are
contained as marks (i.) the concept of freedom, (2.) the concept of adher-
ence to rules (to duty), (3.) the concept of overpowering the force of the
inclinations, in case they oppose those rules. Now if we break up the
concept of virtue into its individual constituent parts, we make it distinct
for ourselves through this analysis. By thus making it distinct, however, we
add nothing to a concept; we only explain it. With distinctness, therefore,
concepts are improved not as to matter but only as to form.
If we reflect on our cognitions in regard to the two essentially different
basic faculties, sensibility and the understanding, from which they arise,
then here we come upon the distinction between intuitions and concepts.
Considered in this respect, all our cognitions are, namely, either intuitions
or concepts. The former have their source in sensibility, the faculty of
intuitions, the latter in the understanding, the faculty of concepts. This is
the logical distinction between understanding and sensibility, according to
which the latter provides nothing but intuitions, the former on the other
hand nothing but concepts. The two basic faculties may of course be
considered from another side and defined in another way: sensibility,
namely, as a faculty of receptivity, the understanding as a faculty of spontane-
ity. But this mode of explanation is not logical but rather metaphysical. It is
also customary to call sensibility the lower faculty, the understanding on
the other hand the higher faculty, on the ground that sensibility gives the
mere material for thought, but the understanding rules over this material
and brings it under rules or concepts.
The difference between aesthetic and logical perfection of cognition is
grounded on the distinction stated here between intuitive and discursive
cognitions, or between intuitions and concepts.
A cognition can be perfect either according to laws of sensibility or
according to laws of the understanding; in the former case it is aesthetically
perfect, in the other logically perfect. The two, aesthetic and logical perfec-
tion, are thus of different kinds; the former relates to sensibility, the latter
to the understanding. The logical perfection of cognition rests on its
agreement with the object, hence on universally valid laws, and hence we
can pass judgment on it according to norms a priori. Aesthetic perfection
consists in the agreement of cognition with the subject and is grounded on
the particular sensibility of man. In the case of aesthetic perfection, there-
fore, there are no objectively and universally valid laws, in relation to
which we can pass judgment on it a priori in a way that is universally valid
for all thinking beings in general. Insofar as there are nonetheless univer-
sal laws of sensibility, which have validity subjectively for the whole of
humanity although not objectively and for all thinking beings in general,
we can think of an aesthetic perfection that contains the ground of a
subjectively universal pleasure. This is beauty, that which pleases the
senses in intuition and can be the object of a universal pleasure just
because the laws of intuition are universal laws of sensibility.
Through this agreement with the universal laws of sensibility the really,
independently beautiful, whose essence consists in mere form, is distin-
guished in kind from the pleasant," which pleases^ merely in sensation
through stimulation or excitement, and which on this account can only be
the ground of a merely private pleasure.*
It is this essential aesthetic perfection, too, which, among all [perfec-
tions], is compatible with logical perfection and may best be combined
with it.
Considered from this side, aesthetic perfection in regard to the essen-
tially beautiful can thus be advantageous to logical perfection. In another
respect it is also disadvantageous, however, insofar as we look, in the case
of aesthetic perfection, only to the non-essentially beautiful, the stimulating
or the exciting, which pleases the senses in mere sensation and does not
relate to mere form but rather to the matter of sensibility. For stimulation
and excitement, most of all, can spoil the logical perfection in our cogni-
tions and judgments.
In general, however, there always remains a kind of conflict between the
aesthetic and the logical perfection of our cognition, which cannot be fully
removed. The understanding wants to be instructed, sensibility enlivened;
the first desires insight, the second comprehensibility. If cognitions are to
instruct then they must to that extent be thorough; if they are to entertain
at the same time, then they have to be beautiful as well. If an exposition is
beautiful but shallow, then it can only please sensibility but not the under-
standing, but if it is thorough yet dry, only the understanding but not
sensibility as well.
Since the needs of human nature and the end of popularity in cognition
demand, however, that we seek to unite the two perfections with one
another, we must make it our task to provide aesthetic perfection for those
cognitions that are in general capable of it, and to make a scholastically
correct, logically perfect cognition popular through its aesthetic form. But
in this effort to combine aesthetic with logical perfection in our cognitions
we must not fail to attend to the following rules, namely: (i.) that logical
perfection is the basis of all other perfections and hence cannot be wholly
subordinated or sacrificed to any other; (2.) that one should look princi-
pally to formal aesthetic perfection, the agreement of a cognition with the
laws of intuition, because it is just in this that the essentially beautiful,
which may best be combined with logical perfection, consists; (3.) that one
must be very cautious with stimulation and excitement, whereby a cognition
affects sensation and acquires an interest for it, because attention can
thereby so easily be drawn from the object to the subject, whence a very
disadvantageous influence on the logical perfection of cognition must
evidently arise.
To acquaint us better with the essential differences that exist between the
logical and the aesthetic perfection of cognition, not merely in the univer-
sal but from several particular sides, we want to compare the two with one
another in respect to the four chief moments of quantity, quality, relation,
and modality, on which the passing of judgment as to the perfection of
cognition depends.
A cognition is perfect (i.) as to quantity if it is universal; (2.) as to quality
if it is distinct; (3.) as to relation if it is true; and finally (4.) as to modality if
it is certain.
Considered from the viewpoints indicated, a cognition will thus be
logically perfect as to quantity if it has objective universality (universality
of the concept or of the rule), as to quality if it has objective distinctness
(distinctness in the concept), as to relation if it has objective truth, and
finally as to modality if it has objective certainty.
To these logical perfections correspond now the following aesthetic
perfections in relation to those four principal moments, namely
aesthetic universality. This consists in the applicability of a cognition to a
multitude of objects that serve as examples, to which application of it can be
made, and whereby it becomes useful at the same time for the end of popular-
ity;
aesthetic distinctness. This is distinctness in intuition, in which a concept
thought abstractly is exhibited or elucidated in concrete through examples;
aesthetic truth. A merely subjective truth, which consists only in the agreement
of cognition with the subject and the laws of sensory illusion, and which is
consequently nothing more than a universal semblance/
aesthetic certainty. This rests on what is necessary in consequence of the
testimony of the senses, i.e., what is confirmed through sensation and experi-
ence.
39
With the perfections just mentioned two things are always to be found,
which in their harmonious union make up perfection in general, namely,
manifoldness and unity. Unity in the concept lies with the understanding,
unity of intuition with the senses.
Mere manifoldness without unity cannot satisfy us. And thus truth is
the principal perfection among them all, because it is the ground of unity
through the relation of our cognition to the object. Even in the case of
aesthetic perfection, truth always remains the conditio sine qua non, the
foremost negative condition, apart from which something cannot please
taste universally. Hence no one may hope to make progress in the belles
lettres^ if he has not made logical perfection the ground of his cognition. It
is in the greatest possible unification of logical with aesthetic perfection in
general, in respect to those cognitions that are both to instruct and to
entertain, that the character and the art of the genius actually shows itself.
VI.
PARTICULAR L O G I C A L P E R F E C T I O N S OF
COGNITION
A) Logical perfection of cognition as to quantity — Quantity —
Extensive and intensive quantity — Extensiveness and thoroughness or
importance andfruitfulness of cognition - Determination of the
horizon of our cognition
The quantity of cognition can be understood in two senses, either as
extensive or as intensive quantity. The former relates to the extension of
cognition and thus consists in its multitude and manifoldness; the latter
*
7
"ein allgemeiner Schein."
"in schönen Wissenschaften."
549
relates to its content,* which concerns the richness" or the logical impor-
tance and fruitfulness of a cognition, insofar as it is considered as ground
of many and great consequences («0« multa sed multum. b )
In expanding our cognitions or in perfecting them as to their extensive
quantity it is good to make an estimate as to how far a cognition agrees
with our ends and capabilities. This reflection concerns the determination
of the horizon of our cognitions, by which is to be understood the congru-
ence' of the quantity of all cognitions with the capabilities and ends of the subject.
The horizon can be determined
1.
2.
41
3.
logically, in accordance with the faculty or the powers of cognition in relation
to the interest of the understanding. Here we have to pass judgment on how far
we can go in our cognitions, how far we must go, and to what extent certain
cognitions serve, in a logical respect, as means to various principal cognitions
as our ends;
aesthetically, in accordance with taste in relation to the interest of feeling. He
who determines his horizon aesthetically seeks to arrange science according
to the taste of the public, i.e., to make it popular, or in general to attain only
such cognitions as may be universally communicated, and in which the class
of the unlearned, too, find pleasure and interest;
practically, in accordance with use in relation to the interest of the will. The
practical horizon, insofar as it is determined according to the influence which
a cognition has on our morality, '^.pragmatic and is of the greatest importance.
Thus the horizon concerns passing judgment on, and determining, what
man can know, what he is permitted to know, and what he ought to know.
Now as for what concerns the theoretically or logically determined hori-
zon in particular - and it is of this alone that we can speak here - we can
consider it either from the objective or from the subjective viewpoint.
In regard to objects, the horizon is either historical or rational. The
former is much broader than the other, indeed, it is immeasurably great,
for our historical cognition has no limits. The rational horizon, on the
other hand, may be fixed, e.g., it may be determined to what kind of
objects mathematical cognition cannot be extended. So too in respect of
philosophical cognition of reason, as to how far reason can go here a priori
without any experience.
In relation to the subject the horizon is either the universal and absolute,
or a particular and conditioned one (a private horizon).
By the absolute and universal horizon is to be understood the congru-
*
" "Gehalt."
" Vielgültigkeit.'"
b Not many but much.
'
"Angemessen/teil."
1
ence * of the limits of human cognitions with the limits of the whole of
human perfection in general. And here, then, the question is: In general,
what can man, as man, know?
The determination of the private horizon depends upon various empiri-
cal conditions' and special considerations, e.g., age, sex, station, mode of
life, etc. Every particular class of men has its particular horizon in relation
to its special powers of cognition, ends, and standpoints, every mind its
own horizon according to the standard of the individuality of its powers
and its standpoint. Finally, we can also think a horizon of healthy reason
and a horizon of science, which latter still requires principles, in accordance
with which to determine what we can and cannot know.
What we cannot know is beyond f our horizon, what we do not need to
know^ is outside* our horizon. This latter can hold only relatively, however,
in relation to various particular private ends, to whose accomplishment
certain cognitions not only do not contribute anything but could even be
an obstacle. For no cognition is, absolutely and for every purpose, useless
and unusable, although we may not always be able to have insight into its
use. Hence it is an objection as unwise as it is unjust that is made to great
men who labor in the sciences with painstaking industry when shallow
minds ask, What is the use ofthat?' We must simply never raise this ques-
tion if we want to occupy ourselves with the sciences. Even granted that a
science could give results only concerning some possible object, it would
still for that reason alone be useful enough. Every logically perfect cogni-
tion always has some possible use, which, although we are as yet unac-
quainted with it, will perhaps be found by posterity. If in the cultivation of
the sciences one had always looked only to material gain, their use, then
we would have no arithmetic or geometry. Besides, our understanding is
so arranged that it finds satisfaction in mere insight, even more than in the
use that arises therefrom. Plato noted this. Man feels in this his own
excellence, he senses what it means to have understanding. Men who do
not sense this must envy the animals. The inner worth that cognitions have
through logical perfection is not to be compared with the outer, their worth
in application.
Like that which lies outside our horizon, insofar as we, in accordance
with our purposes, do not need to know it, as dispensable for us, that which
d
"Congruenz."
' Ak, "empirischen und speciellen Rücksichten"; ist ed., "empirischen Bedingungen und
speciellen Rücksichten."
f
"über."
* "was wir nicht wissen dürfen oder nicht zu wissen brauchen."
* "außer."
'
Reading "wozu das nütze" for "wozu ist das nütze" in accordance with the published list of
printer's errors (KI, xli).
551
lies beneath' our horizon, insofar as we ought not to know it as harmful to
us, is to be understood in a relative sense but never in an absolute one.
With respect to the extension and the demarcation of our cognition, the
following rules are to be recommended:
One must
43
i.
2.
3.
4.
determine his horizon early, but of course only when one can determine it
oneself, which usually does not occur before the 2oth year;
not alter it lightly or often (not turn from one thing to another);
not measure the horizon of others by one's own, and not consider as useless
what is of no use to us; it would be presumptuous to want to determine others'
horizons, because one is not sufficiently acquainted, in part with their capabili-
ties, in part with their purposes;
neither extend it too far nor restrict it too much. For he who wants to know
too much ends by knowing nothing, and conversely, he who believes of some
things that they do not concern him, often deceives himself; as when, e.g., the
philosopher believes of history that it is dispensable for him[.]
One should also seek
5.
6.
7.
8.
44
to determine in advance the absolute horizon of the whole human race (as to
past and to future time), as well as also
to determine, in particular, the position that our science occupies in the whole
of cognition. The Universal Encyclopedia serves for this as a universal map
(mappe-monde k ) of the sciencesf.]
In determining his own particular horizon one should carefully consider for
which part of cognition one has the greatest capability and pleasure, what is
more or less necessary in regard to certain duties, what cannot coexist with
the necessary duties; and finally
one should of course always seek to expand his horizon rather than to narrow
it.
As for the extension of cognition, there need be no concern in general
about what concerned d'Alembert. 16 For the burden does not press us
down, but rather the volume of space for our cognitions constrains us.
Critique of reason, of history and historical writings, a universal spirit that
deals with human cognition en gros and not merely in detail^, will always
make the extension smaller, without diminishing anything in the content.
The metal merely separates from the slag, or the inferior vehicle, the
husk, which was necessary for so long. With the extension of natural
history, of mathematics, etc., new methods will be invented which will
shorten the old and make the multitude of books dispensable. It will be
because of the invention of such new methods and principles that we will
be able, with their help, to find anything we desire without burdening
>
*
"unter.'"
French: mappemonde, map of the world.
memory. Thus he who brings history under ideas that can always remain
renders it service as a genius.
Opposed to the logical perfection of cognition in regard to its extension
stands ignorance.' A negative imperfection, or imperfection of lack, which,
on account of the restrictions of the understanding, is inseparable from
our cognition.
We can consider ignorance from an objective or from a subjective view-
point.
1. Taken objectively, ignorance is either material or formal. The former
consists in a lack of historical cognitions, the other in a lack of rational
cognitions. One does not have to be completely ignorant in any field, but
one can well restrict historical knowledge in order to devote oneself more
to rational knowledge, or conversely.
2. In the subjective sense, ignorance is either learned, scientific, or is
common. He who has distinct insight into the restrictions of cognition,
hence into the field of ignorance from where it begins, e.g., the philoso-
pher who sees and proves how little one can know of gold in regard to its
structure due to a lack of the requisite data, is ignorant artfully™ or in a
learned way. He who is ignorant, on the other hand, without having
insight into the grounds of the limits of knowledge, and without concern-
ing himself with this, is so in a common, not a scientific way. Such a one
does not even know that he knows nothing. For one can never represent
his ignorance except through science, as a blind man cannot represent
darkness until he has become sighted.
Cognition" of one's ignorance presupposes science, then, and makes
one at the same time modest, while imagined knowledge puffs one up.
Hence Socrates' non-knowledge" was a laudable ignorance, really a knowl-
edge of non-knowledge, according to his own admission. It is precisely
those who possess very many cognitions/ then, and who for all that are
astounded at the multitude of what they do not know, who cannot be
reproached with their ignorance.
Ignorance in things whose cognition lies beyond our horizon is in
general irreproachable (inculpabilis), and in regard to the speculative use of
our faculty of cognition it can be allowed (although only in the relative
sense), insofar as the objects here lie not beyond our horizon but yet outside
it. It is shameful, however, in things that it is quite necessary and also easy
to know.
'
m
*
'
'
"Unwissenheit"
"kunstmäßig."
"Kenntniß."
"Nichtwissen."
"Kenntnisse."
553
46
There is a distinction between not knowing something and ignoring
something, i.e., taking no notice of it. It is good to ignore much that it is not
good for us to know. Abstracting is distinct from both of these. One ab-
stracts from a cognition when one ignores its application, whereby one
gets it in abstracto and can better consider it in the universal as a principle.
Such abstraction from what does not belong to our purpose in the cogni-
tion of a thing is useful and praiseworthy.
Scholars in matters of reason* are commonly ignorant historically.
Historical knowledge without determinate limits is poly history; this puffs
one up. Polymathy has to do with cognition of reason. Both historical and
rational knowledge, when extended without determinate limits, can be
called pansophy. Historical knowledge includes the science of die tools of
learnedness - philology, which comprises a critical acquaintance with
books and languages (literature and linguistics).
Mere polyhistory is cyclopic learnedness, which lacks one eye, the eye of
philosophy, and a cyclops among mathematicians, historians, natural histo-
rians, philologists, and linguists is a learned man who is great in all these
matters, but who for all that holds all philosophy to be dispensable.
One part of philology is constituted by the humaniora, by which is
understood acquaintance with the ancients, which furthers die unification
of science with taste, which rubs off coarseness and furthers die communica-
bility and urbanity in which humanity consists.
The humaniora, then, concern instruction in what serves the cultivation
of taste, in conformity with the models of the ancients. This includes, e.g.,
eloquence, poetry, wide reading in the classical authors, etc. All these
humanistic cognitions' can be reckoned in the practical part of philology,
which aims in the first instance at the cultivation of taste. If we separate
the mere philologist from the humanist, however, die two would differ
from one anotiier in dial die former seeks die tools of learnedness among
the ancients, the latter die tools for the cultivation of taste.
The belletrist, or bei esprit, 1 is a humanist according to contemporary
models in living languages. He is not learned, then, for only dead lan-
guages are now learned languages, but is radier a mere dilettante in
cognitions of taste' in accordance with fashion, with no need for the
ancients. We could call him one who apes the humanist. The polyhistor
must, as philologist, be a linguist and a literator, and as humanist a
classicist and expositor of the classics. As philologist he is cultivated, as
humanist civilized.
" Vernunftlehrer."
"Kenntnisse."
French: aesthete, belletrist.
" Geschmackskenntnisse."
In regard to the sciences, there are two degenerate forms of prevailing
taste: pedantry and gallantry." The one pursues the sciences only for the
school and thereby restricts them in respect of their use, the other pursues
them merely for intercourse or for the world and in this way restricts them
in respect of their content.
The pedant is either opposed, as learned man, to the man of the world,
and is to this extent the puffed up man of learning, unacquainted with the
world, i.e., with the ways of bringing his science to bear on men; or he is to
be considered, in general, as a man of skill, but only informalities, not as to
essence and as to his end. In the latter sense he is z fanatic for formalities?
restricted in regard to the core of things, he looks only to the clothing and
the shell. He is the unfortunate imitation or caricature of the methodical
mind. Thus pedantry can also be called cavilling fussiness 1 " and useless
exactitude (micrology) in formalities. And such formality of scholastic
method outside the schools is to be found not merely among learned
people and in learned things, but also in other classes and in other things. 47
What is the ceremonial at court and in intercourse but a pursuit of formalities"
and hair-splitting.' In the military this is not completely so, although it
seems so. But in conversation, in clothing, in diet, in religion, a good deal
of pedantry often prevails.
Thoroughness (scholarly, scholastic perfection 2 ) is a purposeful exacti-
tude in formalities. Pedantry is thus an affected thoroughness, just as
gallantry, as a mere courtesan seeking the approval of taste, is nothing but
an affected popularity. For gallantry only strives to gain the reader's affec-
tion and thus never to insult him with a hard word.
To avoid pedantry requires extensive cognitions" not only in the sci-
ences themselves but also in regard to their use. Only the true man of
learning, then, can free himself from pedantry, which is always the prop-
erty of a restricted mind.
In striving to procure for our cognition the perfection of scholastic
thoroughness and at the same time of popularity, without falling into the
indicated mistakes of affected thoroughness or of affected popularity, we
must look above all to the scholastic perfection of our cognition, the
scholastically correct form of thoroughness; and only then may we con-
cern ourselves about how we are to make our cognition, methodically
learned in school, truly popular, i.e., easily and universally communicable
to others, in such a way that thoroughness is not displaced by popularity.
"
°
"
*
*
z
"
"Galanterie."
"Formalienklauber."
"grüblerische Peinlichkeit."
"Formalienjagd."
"Klauberei."
"schulgerechte, scholastische Vollkommenheit."
"Kenntnisse."
48
For scholastic perfection, without which all science is nothing but tricks
and trifling,* must not be sacrificed for the sake of popular perfection, to
please the people.
To learn true popularity, however, one must read the ancients, e.g.,
Cicero's philosophical writings, the poets Horace, Virgil, etc., and among
the moderns Hume, Shaftesbury, et. al. Men who have all had a good deal
of intercourse with the refined world, without which one cannot be popu-
lar. For true popularity demands a good deal of practical acquaintance
with the world and with men, acquaintance with men's concepts, taste,
and inclinations, to which constant regard must be given in presentation
and even in the choice of expressions that are fitting and adequate to
popularity. This ability to descend' to the public's power of comprehen-
sion and to the customary expressions, in which scholastic perfection is
not slighted, but in which the clothing of thoughts is merely so arranged
that the framework, the scholastically correct and technical in that perfection,
may not be seen (just as one draws lines with a pencil, writes on them, and
subsequently erases them) - this truly popular perfection of cognition is
in fact a great and rare perfection, which shows much insight into the
science. It has this merit, too, in addition to many others, that it can
provide a proof of complete insight into a thing. For the merely scholastic
examination of a cognition leaves doubt as to whether that examination is
not one-sided and whether the cognition itself has a worth admitted by all
men. The school has its prejudices, just as does the common understand-
ing. One improves the other here. It is therefore important that a cogni-
tion be examined by men whose understanding does not depend on any
school.
This perfection of cognition, whereby it qualifies for easy and universal
communication, could also be called external extension/ or the extensive
quantity of a cognition, insofar as it is widespread externally among men.
Since cognitions are so many and manifold, one will do well to make
himself a plan, in accordance with which he orders the sciences in the way
that best agrees with, and contributes to the furtherance of, his ends. All
cognitions stand in a certain natural connection with one another. Now if,
in striving to expand his cognitions, one does not look to their connection,
then extensive knowledge' amounts to nothing more than a mere rhapsody.
If one makes one principal science his end, however, and considers all
other cognitions only as means for achieving it, then he brings a certain
systematic character into his knowledge. And in order to go to work on
*
' "Spielwerk und Tändelei."
"Eine solche Herablassung (Condescendenz)."
4 "äußere Extension."
'
" Vielwissen."
extending his cognitions according to such a well ordered and purposive
plan, one must seek, therefore, to become acquainted with this connection
of cognitions among themselves. For this, the sciences get guidance from
architectonic, which is a system in accordance with ideas, in which the sciences
are considered in regard to their kinship and systematic connection in a whole of
cognition that interests humanity.
49
Now as for what concerns the intensive quantity? of cognition — i.e., its
content, or its richness^ and importance, which is essentially distinct from
its extensive quantity, its mere extensiveness, as we noted above - we want
here to add only the following few remarks:
1.
2.
3.
A cognition that is concerned with what is great, i.e., with the whole in the use
of the understanding, is to be distinguished from subtlety in what is small
(micrology).
Every cognition that furthers logical perfection as to form is to be called
logically important, e.g., every mathematical proposition, every law of nature
into which we have distinct insight, every correct philosophical explanation.
Practical importance cannot be foreseen, one must simply wait and watch for it.
Importance must not be confused with difficulty. 11 A cognition can be difficult
without being important, and conversely. Difficulty, then, does not decide
either for or against the worth or the importance of a cognition. This rests on
the quantity or multiplicity of its consequences. A cognition is the more
important accordingly as it has more or greater consequences, as the use that
may be made of it is more. Cognition without important consequences is
called cavilling;' scholastic philosophy, e.g., was of this sort.
VII.
B) Logical perfection of cognition as to relation - Truth -Material
and formal, or logical, truth - Criteria of logical truth - Falsehood
and error — Illusion, as source of error — Means for avoiding errors
A principal perfection of cognition, indeed, the essential and inseparable
condition of all its perfection, is truth. Truth, it is said, consists in the
agreement of cognition with its object. In consequence of this mere nomi-
nal explanation, my cognition, to count as true, is supposed to agree with
its object. Now I can compare the object with my cognition, however, only
by cognizing it. Hence my cognition is supposed to confirm itself, which is
far short of being sufficient for truth. For since the object is outside me,
f Ak, "intensive Größe"; ist ed., "intensive Größe."
g "Vielgültigkeit."
*
' "Schwere."
"Grübelei."
557
the cognition in me, all I can ever pass judgment on is whether my
cognition of the object agrees with my cognition of the object. The an-
cients called such a circle in explanation a diallelon.' And actually the
logicians were always reproached with this mistake by the skeptics, who
observed that with this explanation of truth it is just as when someone
makes a statement before a court and in doing so appeals to a witness with
whom no one is acquainted, but who wants to establish his credibility by
maintaining that the one who called him as witness is an honest man. The
accusation was grounded, too. Only the solution of the indicated problem
is impossible without qualification and for every man.
The question here is, namely, whether and to what extent there is a
criterion of truth that is certain, universal, and useful in application. For
this is what the question, What is truth?, ought to mean.
To be able to decide this important question we must distinguish that
which belongs to the matter in our cognition and is related to the object
from that which concerns its mere form, as that condition without which a
cognition would in general never be a cognition. With respect to this
distinction between the objective, material relation in our cognition and the
subjective, formal relation, the question above thus breaks down into these
two particular ones:
1.
2.
51
Is there a universal material, and
Is there a universal formal criterion of truth?
A universal material criterion of truth is not possible; it is even self-
contradictory. For as a universal criterion, valid for all objects in general, it
would have to abstract fully from all difference among objects, and yet at
the same time, as a material criterion, it would have to deal with just this
difference, in order to be able to determine whether a cognition agrees
with just that object to which it is related and not just with any object in
general, in which case nothing would really be said. Material truth must
consist in this agreement of a cognition with just that determinate object
to which it is related, however. For a cognition that is true in regard to one
object can be false in relation to other objects. Hence it is absurd to
demand a universal material criterion of truth, which should abstract and
at the same time not abstract from all difference among objects.
If the question is about universal formal criteria of truth, however, then
here it is easy to decide that of course there can be such a thing. For formal
truth consists merely in the agreement of cognition with itself, in complete
abstraction from all objects whatsoever and from all difference among
them. And the universal formal criteria of truth are accordingly nothing
other than universal logical marks of the agreement of cognition with itself
Ak, "Diallele"; ist ed., "Diallele."
or - what is one and the same - with the universal laws of the understand-
ing and of reason.
These formal, universal criteria are of course not sufficient for objec-
tive truth, but they are nonetheless to be regarded as its conditio sine qua
non.
For the question of whether cognition agrees with its objects must be
preceded by the question of whether it agrees with itself (as to form). And
this is a matter for logic.
The formal criteria of truth in logic are
1.
2.
the principle of contradiction,
the principle of sufficient reason.
Through die former the logical possibility of a cognition is determined,
through the latter its logical actuality.
To the logical actuality of a cognition it pertains, namely:
First: that it be logically possible, i.e., not contradict itself. This character-
istic of internal logical truth is only negative, however; for a cognition that
contradicts itself is of course false, but if it does not contradict itself it is
not always true.
Second: that it be logically grounded, i.e., that it (a) have grounds and (b)
not have false consequences.
This second criterion of external logical truth or of accessibility to reason,*
which concerns the logical connection of a cognition with grounds and
consequences, is positive. And here the following rules are valid:
1.
From the truth of the consequence we may infer the truth of the cognition as
ground, but only negatively: if one false consequence flows from a cognition,
then the cognition itself is false. For if the ground were true, then the conse-
quence would also have to be true, because the consequence is determined by
the ground.
But one cannot infer conversely that if no false consequence flows from
a cognition, then it is true; for one can infer true consequences from a
false ground.
2.
If all the consequences of a cognition are true, then the cognition is true too. For if
there were something false in the cognition, then there would have to be a
false consequence too.
From the consequence, then, we may infer to a ground, but without
being able to determine this ground. Only from the complex of all conse-
quences can one infer to a determinate ground, infer that it is the true
ground.
The former mode of inference, according to which the consequence
"Rationabilitat."
559
can only be a negatively and indirectly sufficient criterion of the truth of a
cognition, is called in logic the apagogic mode (modus tollens).
This procedure, of which frequent use is made in geometry, has the
advantage that I may derive just one false consequence from a cognition in
order to prove its falsehood. To show, e.g., that the earth is not flat, I may
just infer apagogically and indirectly, without bringing forth positive and
direct grounds: If the earth were flat, then the pole star would always have
to be at the same height; but this is not die case, consequently it is not flat.
With the other, the positive and direct mode of inference (modus ponens)
the difficulty enters that the totality of the consequences cannot be
cognized apodeictically, and that one is therefore led by the indicated
mode of inference only to a probable and hypothetically true cognition
(hypotheses), in accordance witii the presupposition that where many
consequences are true, all the remaining ones' may be true too.
Thus we will be able to advance three principles here as universal,
merely formal or logical criteria of truth; these are
the principle of contradiction and of identity (principium contradictionis and identita-
tis), through which the internal possibility of a cognition is determined for
problematic judgments;
the principle of sufficient reason (principium rationis suffidentis), on which rests
the (logical) actuality of a cognition, the fact that it is grounded, as material for
assertoric judgments;
the principle of the excluded middle (principium exclusi medii inter duo contradic-
toria m ), on which the (logical) necessity of a cognition is grounded - that we
must necessarily judge thus and not otherwise, i.e., that the opposite is false -
for apodeictic judgments.
The opposite of truth is falsehood, which, insofar as it is taken for truth, is
called error. An erroneous judgment - for there is error as well as truth
only in judgment - is thus one that confuses the illusion of truth with
truth itself.
It is easy to have insight into how truth is possible, since here the under-
standing acts in accordance with its essential laws.
But it is hard to comprehend how error in the formal sense of the word, i.e.,
how the form of thought contrary to the understanding is possible, just as we
cannot in general comprehend how any power should deviate from its own
essential laws. We cannot seek the ground of errors in the understanding
itself and its essential laws, then, just as little as we can in the restrictions of
the understanding, in which lies the cause of ignorance, to be sure, but not
in any way the cause of error. Now if we had no other power of cognition
but the understanding, we would never err. But besides the understand-
ing, there lies in us another indispensable source of cognition. That is
sensibility, which gives us the material for thought, and in doing this works
according to other laws than those the understanding does. Error cannot
arise from sensibility in and by itself, however, because the senses simply
do not judge.
The ground for the origin of all error will therefore have to be sought
simply and solely in the unnoticed influence of sensibility upon the understand-
ing, or to speak more exactly, upon judgment. This influence, namely,
brings it about that in judgment we take merely subjective grounds to be
objective, and consequently confuse the mere illusion of truth with truth itself.
For it is just in this that the essence of illusion consists, which on this
account is to be regarded as a ground for holding a false cognition to be
true.
What makes error possible, then, is illusion, in accordance with which
the merely subjective is confused in judgment with the objective.
In a certain sense, however, one can make the understanding the author
of errors, namely, insofar as it allows itself, due to a lack of requisite
attention to that influence of sensibility, to be misled by the illusion arising
therefrom into holding merely subjective determining grounds of judg-
ment to be objective ones, or into letting that which is true only according
to the laws of sensibility hold as true in accordance with its own laws.
In the restrictions of the understanding, then, lies only the responsibil-
ity for ignorance; the responsibility for error we have to assign to our-
selves. Nature has denied us many cognitions, to be sure, it leaves us in
unavoidable ignorance concerning so much, but still it does not cause
error. We are misled into this by our own inclination to judge and to
decide even where, on account of our limitedness, we are not able to
judge and to decide.
Every error into which the human understanding can fall is only partial,
however, and in every erroneous judgment there must always lie some-
thing true. For a total error would be a complete opposition to the laws of
the understanding and of reason. But how could that, as such, in any way
come from the understanding and, insofar as it is still a judgment, be held
to be a product of the understanding.
In respect to the true and the erroneous in our cognition, we distinguish
an exact cognition from a rough one.
Cognition is exact when it is adequate to its object, or when there is not
the slightest error in regard to its object, and it is rough when there can be
errors in it yet without being a hindrance to its purpose.
This distinction concerns the broader or narrower determinateness of our
cognition (cognitio late vel stride determinata). Initially it is sometimes neces-
sary to determine a cognition in a broader extension (late determinare),
561
54
particularly in historical things. In cognitions of reason everything must be
determined exactly (striae), however. In the case of broad determination
one says that a cognition is determined praeter propter." Whether a cogni-
tion ought to be determined roughly or exactly always depends on its
purpose. Broad determination leaves a certain play for error, which still
can have its determinate limits, however. Error occurs particularly where a
broad determination is taken for a strict one, e.g., in matters of morality,"
where everything must be determined striae. Those who do not do this are
called by die English latitudinarians.
One can distinguish subtlety, as a subjective perfection of cognition, from
exactness, as an objective perfection - since here cognition is fully congru-
ent with its object.
A cognition is subtle when one discovers in it that which usually escapes
the attention of others. It requires a higher degree of attention, then, and
a greater application of power of the understanding.
Many reprove all subtlety because they cannot attain it. But in itself it
always brings honor to the understanding and is even laudable and neces-
sary, insofar as it is applied to an object worthy of observation. When one
could have attained the same end with less attention and effort of the
understanding, however, and yet one uses more, then one makes a useless
expense and falls into subtleties, which are difficult, to be sure, but do not
have any use (nugae difficiles p ).
As the rough is opposed to the exact, so is the crude to the subtle.
56
From the nature of error - whose concept, as we noted, contains as an
essential mark, besides falsehood, also the illusion of truth - we get the
following important rule for the truth of our cognition:
To avoid errors - and no error is unavoidable, at least not absolutely or
without qualification, although it can be unavoidable relatively, for the
cases where it is unavoidable for us to judge, even with the danger of
error - to avoid errors, then, one must seek to disclose and to explain
their source, illusion. Very few philosophers have done that, however.
They have only sought to refute the errors themselves, without indicating
the illusion from which they arise. This disclosure and breaking up of
illusion is a far greater service to truth, however, than the direct refutation
of errors, whereby one does not block their source and cannot guard
against the same illusion misleading one into errors again in other cases
because one is not acquainted with it. For even if we are convinced that we
have erred, then in case the illusion diat grounds our error has not been
"
°
f
approximately.
Ak, "Modalität"; ist ed., "Moralität."
difficult trivialities.
562THE J Ä S C H E LOGIC
removed we still have scruples, however little we can bring forth in justifica-
tion of them.
Through the explanation of illusion, furthermore, one grants to the one
who erred a kind of fairness. For no one will admit that he erred without
any illusion of truth, which might even have deceived someone more
acute, because here it is a matter of subjective grounds.
Where the illusion is evident even to the common understanding (sensus
communis), an error is called a stupidity or an absurdity. The charge of
absurdity is always a personal reproof, which one must avoid, particularly
in the refutation of errors.
For to him who maintains an absurdity, the very illusion that lies at the
ground of the evident falsehood is not evident. One must first make this
illusion evident to him. Then if he still persists, he is admittedly stupid;
but then nothing more can be undertaken with him either. He has thereby
made himself just as incapable of further correction and refutation as he is
unworthy of it. For one cannot really prove to anyone that he is absurd;
here all ratiocination would be vain. If one proves absurdity, then one is no
longer speaking with him who erred but with him who is rational. But then
the disclosure of the absurdity (deductio ad absurdum) is not necessary.
One can also call a stupid error one that nothing, not even illusion, serves
to excuse; just as a crude error is an error that proves ignorance in common
cognition or a slip in common attentiveness.
Error in principles is greater than in their application.
An external mark or an external touchstone of truth is the comparison of
our own judgments with those of others, because the subjective will not be
present in all others in the same way, so that illusion can thereby be
cleared up. The incompatibility of the judgments of others with our own is
thus an external mark of error and is to be regarded as a cue to investigate
our procedure in judgment, but not for that reason to reject it at once. For
one can perhaps be right about the thingbut not right in manner, i.e., in the
exposition.
The common human understanding (sensus communis) is also in itself a
touchstone for discovering the mistakes of the artificial use of the under-
standing. This is what it means to orient oneself in thought or in the
speculative use of reason by means of the common understanding, when
one uses the common understanding as a test for passing judgment on the
correctness of the speculative use.
Universal rules and conditions for avoiding error in general are: i) to
think for oneself, 2) to think oneself in the position of someone else, and
3) always to think in agreement with oneself. The maxim of thinking for
563
oneself can be called the enlightened mode of thought; the maxim of putting
oneself in the viewpoint of others in thought, the extended mode of thought;
and the maxim of always thinking in agreement with one self, the conse-
quent 9 or coherent' mode of thought.
58
VIII.
C) Logical perfection of cognition as to quality — Clarity — Concept of
a mark in general - Various kinds of marks — Determination of the
logical essence of a thing — Its distinction from the real essence —
Distinctness, a higher degree of clarity -Aesthetic and logical
distinäness — Distinction between analytic and synthetic distinctness
From the side of the understanding, human cognition is discursive, i.e., it
takes place through representations which take as the ground of cognition
that which is common to many things, hence through marks' as such. Thus
we cognize things through marks and that is called cognizing,' [the German
word for which] comes from [the German word for] being acquainted."
A mark is that in a thing which constitutes apart of the cognition of it, or -
what is the same - a partial representation, insofar as it is considered as
ground of cognition of the whole representation. All our concepts are marks,
accordingly, and all thought is nothing other than a representing through
marks.
Every mark may be considered from two sides:
First, as a representation in itself;
Second, as belonging, as a partial concept, to the whole representation of
a thing, and thereby as ground of cognition of this thing itself.
All marks, considered as grounds of cognition, have two uses, either an
internal or an external use. The internal use consists in derivation, in order
to cognize the thing itself through marks as its grounds of cognition. The
external use consists in comparison, insofar as we can compare one thing
with others through marks in accordance with the rules of identity or
diversity.
There are many specific differences among marks, on which the following
classification of them is grounded.
"consequente."
"Merkmale."
"Erkennen."
"Kennen."
564THE J Ä S C H E LOGIC
1. Analytic or synthetic marks. The forme r are partial concepts of my actual concept
(marks that I already think therein), while the latter are partial concepts of the
merely possible complete concept (which is supposed to come to be through a
synthesis of several parts). The former are all concepts of reason, the latter can
be concepts of experience.
2. Coordinate or subordinate. This division of marks concerns their connection
after" or under" one another.
59
Marks are coordinate insofar as each of them is represented as an immedi-
ate mark of the thing and are subordinate insofar as one mark is repre-
sented in the thing only by means of the other. The combination of
coordinate marks to form the whole of a concept is called an aggregate, the
combination of subordinate concepts a series. The former, the aggregation
of coordinate marks, constitutes the totality of the concept, which, in
regard to synthetic empirical concepts, can never be completed, but rather
resembles a straight line without limits.
The series of subordinate marks terminates a pane ante, or on the side
of the grounds, in concepts which cannot be broken up, which cannot be
further analyzed on account of their simplicity; a pane post, or in regard to
the consequences, it is infinite, because we have a highest genus but no
lowest species.
With the synthesis of every new concept in the aggregation of coordi-
nate marks, the extensive or extended distinctness grows, as intensive or deep
distinctness grows with the further analysis of the concept in the series of
subordinate marks. This latter kind of distinctness, since it necessarily
contributes to thoroughness and coherence of the cognition, is thus princi-
pally a matter of philosophy and is pursued to the highest degree in
metaphysical investigations in particular.
3. Affirmative or negative marks. Through the former we cognize what the thing
is, through the latter what it is not.
Negative marks serve to keep us from errors. Hence they are unneces-
sary where it is impossible to err, and are necessary and of importance only in
those cases where they keep us from an important error into which we can
easily fall. Thus in regard to the concept, e.g., of a being like God, negative
marks are quite necessary and important.
Through affirmative marks we seek to understand something, through
negative marks - into which all marks can be transformed - we only seek
not to misunderstand or not to err, even if we should not thereby become
acquainted with anything.
A mark is important and fruitful if it is a ground of cognition for great
and numerous consequences, partly in regard to its internal use, its use in
derivation, insofar as it is sufficient for cognizing thereby a great deal in
the thing itself, partly in respect to its external use, its use in comparison,
insofar as it thereby contributes to cognizing both the similarity of a thing
to many others and its difference from many others.
We have to distinguish logical importance and fruitfulness from practical,
from usefulness and utility, by the way.
Sufficient
and necessary or insufficient
and accidental marks.
A mark is sufficient insofar as it suffices always to distinguish the thing
from all others; otherwise it is insufficient, as the mark of barking is, for
example, for dogs. The sufficiency^ of marks, as well as their importance,
is to be determined only in a relative sense, in relation to ends that are
intended through a cognition.
Necessary marks, finally, are those that must always be there to be found
in the thing represented. Marks of this sort are also called essential and are
opposed to extra-essential and accidental marks, which can be separated
from the concept of the thing.
Among necessary marks there is another distinction, however.
Some of them belong to the thing as grounds of other marks of one and
the same thing, while others belong only as consequences of other marks.
The former are primitive and constitutive marks (constitutiva, essentialia in
sensu strictissimo z ), the others are called attributes (consectaria, rationata")
and belong admittedly to the essence of the thing, but only insofar as they
must first be derived from its essential points, as the three angles follow
from the three sides in the concept of the triangle, for example.
Extra-essential marks are again of two kinds; they concern either internal
determinations of a thing (moat) or its external relations (relationes). Thus
the mark of learnedness signifies an inner determination of a man, but being
a master or a servant only an external relation.
The complex of all the essential parts of a thing, or the sufficiency of its
marks as to coordination or subordination, is the essence (complexus notarum
primitivarum, interne conceptui dato sufficientium; s. complexus notarum, con-
ceptum aliquem primitive constituentium b ).
In this explanation, however, we must not think at all of the real or
*
"Zureichende."
"Hinlänglichkeit."
* things that are constitutive, things that are essential in the strictest sense.
" things that follow, things grounded.
* the complex of primitive marks internally sufficient for a given concept, or the complex of
marks that primitively constitute a certain concept.
y
natural essence of things, into which we are never able to have insight. For
since logic abstracts from all content of cognition, and consequently also
from the thing itself, in this science the talk can only be of the logical
essence of things. And into this we can easily have insight. For it includes
nothing further than the cognition of all the predicates in regard to which
an object is determined through its concept; whereas for the real essence of
the thing (esse ret) we require cognition of those predicates on which, as
grounds of cognition, everything that belongs to the existence of the thing
depends. If we wish to determine, e.g., the logical essence of body, then
we do not necessarily have to seek for the data for this in nature; we may
direct our reflection to the marks which, as essential points (constitutiva,
rationes) originally constitute the basic concept of the thing. For the logical
essence is nothing but the first basic concept of all the necessary marks of a
thing (esse conceptus).
The first stage of the perfection of our cognition as to quality is thus its
clarity. A second stage, or a higher degree of clarity, is distinctness. This
consists in clarity of marks.
First of all we must here distinguish logical distinctness in general from
aesthetic distinctness. Logical distinctness rests on objective clarity of
marks, aesthetic distinctness on subjective clarity. The former is a clarity
through concepts, the latter a clarity through intuition. The latter kind of
distinctness consists, then, in a mere liveliness and understandability, i.e., in
a mere clarity through examples in concreto (for much that is not distinct
can still be understandable, and conversely, much that is hard to under-
stand can still be distinct, because it goes back to remote marks, whose
connection with intuition is possible only through a long series).
Objective distinctness frequently causes subjective obscurity, and con-
versely. Hence logical distinctness is often possible only to the det-
riment of aesthetic distinctness, and conversely aesthetic distinctness
through examples and similarities which do not fit exactly but are only
taken according to an analogy often becomes harmful to logical distinct-
ness. Besides, examples are simply not marks and do not belong to the
concept as parts but, as intuitions, to the use of the concept. Dis-
tinctness through examples, mere understandability, is hence of a com-
pletely different kind than distinctness through concepts as marks.
Lucidity consists in the combination of both, of aesthetic or popular
distinctness and of scholastic or logical distinctness. For one thinks of a
lucid mind as the talent for a luminous presentation of abstract and
thorough cognitions that is congruent with the common understanding's
power of comprehension.
Next, as for what concerns logical distinctness in particular, it is to be
called complete distinctness insofar as all the marks which, taken together,
make up the whole concept have come to clarity. A completely' distinct
concept can be so, again, either in regard to the totality of its coordinate
marks or in respect to the totality of its subordinate marks. Extensively
complete or sufficient distinctness of a concept consists in the total clarity
of its coordinate marks, which is also called exhaustiveness. Total clarity of
subordinate marks constitutes intensively complete distinctness, profundity.
The former kind of logical distinctness can also be called the external
completeness (completudo externd) of the clarity of marks, the other the inter-
nal completeness (completudo internet). The latter can be attained only with
pure concepts of reason and with arbitrary concepts, but not with empiri-
cal concepts.
The extensive quantity of distinctness, insofar as it is not superfluous/ is
called precision/ Exhaustiveness^ (completudo) and precision (praecisio) to-
gether constitute adequacy 1 (cognitio, quae rem adaequat h ); and the completed
perfection of a cognition (consummata cognitionis perfectio) consists (as to qual-
ity) in intensively adequate cognition, profundity, combined with extensively
adequate cognition, exhaustiveness and precision.
Since, as we have noted, it is the business of logic to make clear concepts
distinct, the question now is in what way it makes them distinct.
Logicians of the Wolffian school place the act of making cognitions
distinct' entirely in mere analysis of them. But not all distinctness rests on
analysis of a given concept. It arises thereby only in regard to those marks
that we already thought in the concept, but not in respect to those marks
that are first added to the concept as parts of the whole possible concept.
The kind of distinctness that arises not through analysis but through
synthesis of marks is synthetic distinctness. And thus there is an essential
difference between the two propositions: to make a distinct concept and to
make a concept distinct.
For when I make a distinct concept, I begin with the parts and proceed
from these toward the whole. Here there are no marks as yet at hand; I
acquire them only through synthesis. From this synthetic procedure
emerges synthetic distinctness, then, which actually extends my concept
as to content through what is added as a mark beyond' the concept in (pure
or empirical) intuition. The mathematician and the natural philosopher
make use of this synthetic procedure in making distinctness in concepts.*
For all distinctness of properly mathematical cognition, as of all cognition
based on experience, rests on such an expansion of it through the synthe-
sis of marks.
When I make a concept distinct, however, my cognition does not grow
at all as to content through this mere analysis. The content remains the
same, only the form is altered, in that I learn to distinguish better, or to
cognize with clearer consciousness, what lay in the given concept already.
As nothing is added to a map through the mere illumination' of it, so a
given concept is not in the least increased through its mere illumination'"
by means of the analysis of its marks.
To synthesis pertains the making distinct of objects," to analysis the
making distinct of concepts." In the latter case the whole precedes the parts, in
the former the parts precede the whole. The philosopher only makes given
concepts distinct. Sometimes one proceeds synthetically even when the
concept that one wants to make distinct in this way is already given. This is
often the case with propositions based on experience, in case one is not
yet satisfied with the marks already thought in a given concept.
The analytic procedure for creating distinctness, with which alone logic
can occupy itself, is the first and principal requirement in making our
cognition distinct. For the more distinct our cognition of a thing is, the
stronger and more effective it can be too. But analysis must not go so far
that in the end the object itself disappears.
If we were conscious of all that we know, we would have to be aston-
ished at the great multitude of our cognitions.
In regard to the objective content of our cognition in general, we may
think the following degrees, in accordance with which cognition can, in this
respect, be graded:
The first degree of cognition is: to represent something;' 1
The second: to represent something with consciousness, or to perceive*
(percipere);
The third: to be acquainted' with something (noscere), or to represent
something in comparison with other things, both as to sameness 1 and as to
difference;'
The fourth: to be acquainted with something with consciousness, i.e., to
cognize" it (cognoscere). Animals are acquainted with objects too, but they do
not cognize them.
The fifth: to understand 11 something (intelligere), i.e., to cognize some-
thing through the understanding by means of concepts, or to conceive." One can
conceive* much, although one cannot comprehend-' it, e.g., a perpetuum
mobile, whose impossibility is shown in mechanics.
The sixth: to cognize something through reason, or to have insighf into
it (perspicere). With few things do we get this far, and our cognitions
become fewer and fewer in number the more that we seek to perfect them
as to content.
The seventh, finally: to comprehend" something (comprehendere), i.e., to
cognize something through reason or a priori to the degree that is suffi-
cient for our purpose. For all our comprehension is only relative, i.e.,
sufficient for a certain purpose; we do not comprehend anything without
qualification. Nothing can be comprehended more than what the mathe-
matician demonstrates, e.g., that all lines in the circle are proportional.
And yet he does not comprehend how it happens that such a simple
figure has these properties. The field of understanding or of the under-
standing is thus in general much greater than the field of comprehension
or of reason.
IX.
D) Logical perfection of cognition as to modality certainty — Concept of
holding-to-be-true in general- Modi of holding-to-be-true:
opining, believing and knowing - Conviction and persuasion -
Reservation and deferral of a judgment — Provisional judgments —
Prejudices, their sources and principal kinds
66
Truth is an objective property of cognition; the judgment through which
something is represented as true, the relation to an understanding and thus
to a particular subject, is, subjectively, holding-to-be-true. b
'
"Einerleiheit."
' "Verschiedenheit."
" "erkennen."
" "verstehen."
" "condpiren."
"
* "Concipiren."
"begreifen."
* "einsehen."
°
* "begreifen."
"Fiinvahrhalten."
Holding-to-be-true is in general of two kinds, certain or uncertain. Cer-
tain holding-to-be-true, or certainty, is combined with consciousness of
necessity, while uncertain holding-to-be-true, or uncertainty, is combined
with consciousness of the contingency or the possibility of the opposite.
The latter is again either subjectively as well as objectively insufficient, or
objectively insufficient but subjectively sufficient. The former is called opinion,'
the latter must be called belief. d
Accordingly, there are three kinds or modi of holding-to-be-true: opin-
ing' believing/ and knowing. 1 Opining is problematic judging, believing is
assertoric judging, and knowing is apodeictic judging. For what I merely
opine I hold in judging, with consciousness, only to be problematic; what I
believe I hold to be assertoric, but not as objectively necessary, only as
subjectively so (holding only for me); what I know, finally, I hold to be
apodeiäically certain, i.e., to be universally and objectively necessary (hold-
ing for all), even granted that the object to which this certain holding-to-
be-true relates should be a merely empirical truth. For this distinction in
holding-to-be-true according to the three modi just named concerns only
the power of judgment in regard to the subjective criteria for subsumption of
a judgment under objective rules.
Thus, for example, our holding-to-be-true of immortality would be
merely problematic in case we only act as if we were immortal, but it would
be assertoric in case we believe that we are immortal, and it would be apodeictic,
finally, in case we all knew that there is another life after this one.
There is an essential difference, then, between opining, believing, and
knowing, which we wish to expound more exactly and in more detail here.
i. Opining. Opining, or holding-to-be-true based on a ground of cogni-
tion that is neither subjectively nor objectively sufficient, can be regarded
as provisional judging (sub conditione suspensiva ad interim) that one cannot
easily dispense with. One must first opine before one accepts and main-
tains, but in doing so must guard oneself against holding an opinion to be
something more than mere opinion. For the most part, we begin with
opining in all our cognizing. Sometimes we have an obscure premonition
of truth, a thing seems to us to contain marks of truth; we suspect its truth
even before we cognize it with determinate certainty.
But now where does mere opining really occur? Not in any sciences that
contain cognitions a priori, hence neither in mathematics nor in metaphys-
ics nor in morals, but merely in empirical cognitions: in physics, psychol-
ogy, etc. For it is absurd to opine a priori. In fact, too, nothing could be
more ridiculous than, e.g., only to opine in mathematics. Here, as in
'
d
'
f
1
"Meinung."
"Glaube."
"Meinen."
"Glauben."
"Wissen."
571
metaphysics and in morals, the rule is either to know or not to know. Thus
matters of opinion can only be objects of a cognition by experience, a
cognition which is possible in itself but impossible for us in accordance
with the restrictions and conditions of our faculty of experience and the
attendant degree of this faculty that we possess. Thus, for example, the
ether of modern physicists is a mere matter of opinion. For with this as
with every opinion in general, whatever it may be, I see that the opposite
could perhaps yet be proved. Thus my holding-to-be-true is here both
objectively and subjectively insufficient, although it can become complete,
considered in itself.
2. Believing. Believing, or holding-to-be-true based on a ground that is
objectively insufficient but subjectively sufficient, relates to objects in
regard to which we not only cannot know anything but also cannot opine
anything, indeed, cannot even pretend there is probability, but can only be
certain that it is not contradictory to think of such objects as one does
think of them. What remains here is a free holding-to-be-true, which is
necessary only in a practical respect given a priori, hence a holding-to-be-
true of what I accept on moral grounds, and in such a way that I am certain
that the opposite can never be proved.*
68
* Believing is not a special source of cognition. It is a kind of incomplete holding-to-be-true
with consciousness, and if considered as restricted to a particular kind of object (which
pertains only to believing), it is distinguished from opining not by its degree but rather by the
relation that it has as cognition to action. Thus the businessman, for example, to strike a
deal, needs not just to opine that there will be something to be gained thereby, but to believe
it, i.e., to have his opinion be sufficient for an undertaking into the uncertain. Now we have
theoretical cognitions (of the sensible) in which we can come to certainty, and in regard to
everything that we can call human cognition this latter must be possible. We have just such
certain cognitions, and in fact completely a priori, in practical laws, but these are grounded
on a supersensible principle (of freedom) and in fact in us ourselves, as a principle of practical
reason. But this practical reason is a causality in regard to a likewise supersensible object, the
highest good, which is not possible through our faculty in the sensible world. Nature as object
of our theoretical reason must nonetheless agree with this, for the consequence or effect of this
idea is supposed to be met with in the world of the senses. Thus we ought to act so as to
make this end actual.
Now in the world of the senses we also find traces of an artistic misdom, h and we believe
that the cause of the world also works with moral wisdom toward the highest good. This is a
holding-to-be-true that is enough for action, i.e., a belief. Now we do not need this for action
in accordance with moral laws, for these are given through practical reason alone, but we
need to accept a highest wisdom as the object of our moral will, an object beyond the mere
legitimacy of our actions, toward which we cannot avoid directing our ends. Although
objectively this would not be a necessary relation of our faculty of choice, subjectively the
highest good is still necessarily the object of a good (even of a human) will, and hence belief
in its attainability is necessarily presupposed.
There is no mean between the acquisition of a cognition through experience (a posteriori)
and through reason (apriori). But there is a mean between the cognition of an object and the
mere presupposition of its possibility, namely, an empirical ground or a ground of reason for
'Kunstmeisheit."
Matters of belief are thus I) not objects of empirical cognition. Hence so-
called historical belief cannot really be called belief, either, and cannot be
opposed as such to knowledge, since it can itself be knowledge. Holding-
to-be-true based on testimony is not distinguished from holding-to-be-
true through one's own experience either as to degree or as to kind.
II) [N]or [are they] objects of cognition by reason (cognition a priori),
whether theoretical, e.g., in mathematics and metaphysics, or practical, in
morals.
One can believe mathematical truths of reason on testimony, to be sure,
partly because error here is not easily possible, partly, too, because it can
easily be discovered, but one cannot know them in this way, of course. But
accepting this possibility in relation to a necessary extension of the field of possible objects
beyond those whose cognition is possible for us. This necessity occurs only in regard to that
in which the object is cognized as practical and, through reason, as practically necessary, for
to accept something on behalf of the mere extension of theoretical cognition is always
contingent. This practically necessary presupposition of an object is the presupposition of the
possibility of the highest good as object of choice, hence also of the condition of this
possibility (God, freedom, and immortality). This is a subjective necessity to accept the
reality of the object for the sake of the necessary determination of the will. This is the casus
extravrdinarius, without which practical reason cannot maintain itself in regard to its neces-
sary end, and here a favor necessitate proves useful to it in its own judgment. It cannot acquire
an object logically, but can only oppose what hinders it in the use of this idea, which belongs
to it practically.
This belief is the necessity to accept the objective reality of a concept (of the highest
good), i.e., the possibility of its object, asapriori necessary object of choice. If we look merely
to actions, we do not need this belief. But if we wish to extend ourselves through actions to
possession of the end that is thereby possible, then we must accept that this end is com-
pletely possible. Hence I can only say that / see myself necessitated through my end, in
accordance with laws of freedom, to accept as possible a highest good in the world, but I
cannot necessitate anyone else through grounds (the belief is free).
A belief of reason can never aim at theoretical cognition, then, for there objectively
insufficient holding-to-be-true is merely opinion. It is merely a presupposition of reason for
a subjective though absolutely necessary practical purpose. The sentiment toward moral
laws leads to an object of choice, which [choice] is determinable through pure reason. The
acceptance of the feasibility of this object, and hence of the reality of its cause, is a moral
belief, or a free holding-to-be-true that is necessary for moral purposes for completion of
one's ends.
Fides is really good faith 1 in thepactum, or subjective trust-' in one another, that one will keep
his promise to the other, with full faith and credit.* The first when the pactum is made, the
second when it is to be concluded.
In accordance with the analogy, practical reason is, as it were, the promisor,' man the
promissee," the good expected from the deed the promised."
'
'
*
'
m
"
"Treue."
"subjectives Zutrauen."
"Treue und Glauben."
"der Prominent."
"der Promissarius."
"das Promissum."
573
68
70
71
philosophical truths of reason may not even be believed, they must simply
be known; for philosophy does not allow mere persuasion. And as for
what concerns in particular the objects of practical cognition by reason in
morals, rights and duties, there can just as little be mere belief in regard to
them. One must be fully certain whether something is right or wrong, in
accordance with duty or contrary to duty, allowed or not allowed. In moral
things one cannot risk anything on the uncertain, one cannot decide any-
thing on the danger of trespass against the law. Thus it is not enough for the
judge, for example, that he merely believe that someone accused of a crime
actually committed this crime. He must know it (juridically), or he acts
unconscientiously.
Ill) The only objects that are matters of belief are those in which
holding-to-be-true is necessarily free, i.e., is not determined through
objective grounds of truth that are independent of the nature and the
interest of the subject.
Thus also on account of its merely subjective grounds, believing yields
no conviction that can be communicated and that commands universal
agreement, like the conviction that comes from knowledge. Only / myself
can be certain of the validity and unalterability of my practical belief, and
my belief in the truth of a proposition or the actuality of a thing is what
takes the place of a cognition only in relation to me without itself being a
cognition.
He who does not accept what it is impossible to know but morally neces-
sary to presuppose is morally unbelieving. At the basis of this kind of
unbelief lies always a lack of moral interest. The greater a man's moral
sentiment," the firmer and more lively will be his belief in all that he feels
himself necessitated to accept and to presuppose out of moral interest, for
practically necessary purposes.
3. Knowing. Holding-to-be-true based on a ground of cognition that is
objectively as well as subjectively sufficient, or certainty, is either empirical
or rational, accordingly as it is grounded either on experience - one's own
as well as that communicated by others - or on reason. This distinction
relates, then, to the two sources from which the whole of our cognition is
drawn: experience and reason.
Rational certainty, again, is either mathematical or philosophical cer-
tainty. The former is intuitive/ the latter discursive.
Mathematical certainty is also called evidence," because an intuitive cog-
nition is clearer than a discursive one. Although the two, mathematical
and philosophical cognition of reason, are in themselves equally certain,
the certainty is different in kind in them.
' "moralische Gesinnung."
' "intuitiv."
"Evidenz."
f
Empirical certainty is original (originarie empirica) insofar as I become
certain of something from my own experience, and derived (derivative em-
pirica) insofar as I become certain through someone else's experience. The
latter is also usually called historical certainty.
Rational certainty is distinguished from empirical certainty by the con-
sciousness of necessity that is combined with it; hence it is apodeictic cer-
tainty, while empirical certainty is only assertoric. We are rationally certain
of that into which we would have had insight a priori even without any
experience. Hence our cognitions can concern objects of experience and
the certainty concerning them can still be both empirical and rational at
the same time, namely, insofar as we cognize an empirically certain propo-
sition from principles a priori.
We cannot have rational certainty of everything, but where we can have
it, we must put it before empirical certainty.
All certainty is either unmediated or mediated, i.e., it either requires a
proof, or it is not capable of and does not require any proof. Even if so
much in our cognition is certain only mediately, i.e., through a proof,
there must still be something indemonstrable or immediately certain, and the
whole of our cognition must proceed from immediately certain proposi-
tions.
The proofs on which any mediated or mediate certainty of a cog-
nition rests are either direct proofs or indirect, i.e., apagogical ones.
When I prove a truth from its grounds I provide a direct proof for it,
and when I infer the truth of a proposition from the falsehood of its
opposite I provide an indirect one. If this latter is to have validity,
however, the propositions must be opposed contradictorily or diametra-
liter. For two propositions opposed only as contraries (contrarie oppo-
sita) can both be false. A proof that is the ground of mathematical
certainty is called a demonstration, and that which is the ground of philo-
sophical certainty is called an acroamatic proof. The essential parts of
any proof in general are its matter and its form, or the ground of proof
and the consequentia.
From [the German word for] knowing' comes [the German word for]
science, 5 by which is to be understood the complex of a cognition as a
system. It is opposed to common cognition, i.e., to the complex of a cogni-
tion as mere aggregate. A system rests on an idea of the whole, which
precedes the parts, while with common cognition on the other hand, or a
mere aggregate of cognitions, the parts precede the whole. There are
historical sciences and sciences of reason.
In a science we often know only the cognitions but not the things repre-
" Wissen."
"Wissenschaft."
575
sented through them; hence there can be a science of that of which our
cognition is not knowledge.
From the foregoing observations concerning the nature and the kinds of
holding-to-be-true we can now draw the universal result that all our
conviction is thus either logical or practical. When we know, namely, that
we are free of all subjective grounds and yet the holding-to-be-true is
sufficient, then we are convinced, and in fact logically convinced, or con-
vinced on objective grounds (the object is certain).
Complete holding-to-be-true on subjective grounds, which in a practi-
cal relation hold just as much as objective grounds, is also conviction,
though not logical but rather praäical conviction (/ am certain). And this
practical conviction, or this moral belief of reason,' is often firmer than all
knowledge. With knowledge one still listens to opposed grounds, but not
with belief, because here it does not depend on objective grounds but on
the moral interest of the subject.*
73
Opposed to conviction stands persuasion, a holding-to-be-true on insuffi-
cient grounds, of which one does not know whether they are merely
subjective or also objective.
Persuasion often precedes conviction. We are conscious of many cogni-
tions only in such a way that we cannot judge whether the grounds of our
holding-to-be-true are objective or subjective. To be able to pass from
mere persuasion to conviction, then, we must first of all reflect, i.e., see to
which power of cognition a cognition belongs, and then investigate, i.e.,
test whether the grounds are sufficient or insufficient in regard to the
object. Many remain with persuasion. Some come to reflection, few to
investigation. He who knows what pertains to certainty will not easily mix
up persuasion and conviction, and hence will not let himself be easily
* This practical conviction is thus moral belief of reason," which alone can be called a belief in
the proper sense and be opposed as such to knowledge and to all theoretical or logical
conviction in general, because it can never elevate itself to knowledge. So-called historical
belief, on the other hand, as already observed, may not be distinguished from knowledge,
since as a kind of theoretical or logical holding-to-be-true it can itself be knowledge. We can
accept an empirical truth on the testimony of others with the same certainty as if we had
attained it through facta of our own experience. In the former kind of empirical knowledge
there is something deceptive, but also with the latter kind.
Historical or mediate empirical knowledge rests on die reliability of testimony. The
requirements of an irrefutable witness include authenticity* (competence") and integrity.'
' "moralische Vernunftglaube."
*
*' "der moralische Vemunfiglaube."
"Authentizität."
" "Tüchtigkeit."
* "Integrität."
persuaded, either. There is a ground of determination to approval, which
is composed of objective and subjective grounds, and most men do not
analyze^ this mixed effect.
Although all persuasion is false as to form (formaliter), namely, insofar
as an uncertain cognition appears here to be certain, it can nonetheless be
true as to matter (materialiter). And thus it is distinct from opinion, too,
which is an uncertain cognition, insofar as it is held to be uncertain.
The sufficiency of holding-to-be-true (in belief) can be put to the test
by betting and by taking oaths. For the first comparative sufficiency of objec-
tive grounds is necessary, for the second absolute sufficiency, instead of
which, if this is not available, a merely subjectively sufficient holding-to-
be-true nevertheless holds.
It is customary to use the expressions, to agree with someone's judgment, to
reserve, to defer, or give up one's judgment. These and similar expressions
seem to indicate that there is something arbitrary 2 in our judging, in that
we hold something to be true because we want to hold it to be true. The
question arises, accordingly, whether willing has an influence on our judg-
ments.
The will does not have any influence immediately on holding-to-be-
true; this would be quite absurd. When it is said that we gladly believe what
we wish, this means only our benign wishes, e.g., those of a father for his
children. If the will had an immediate influence on our conviction concern-
ing what we wish, we would constantly form for ourselves chimeras of a
happy condition, and always hold them to be true, too. But the will cannot
struggle against convincing proofs of truths that are contrary to its wishes
and inclinations.
Insofar as the will either impels the understanding toward inquiry into a
truth or holds it back therefrom, however, one must grant it an influence
on the use of the understanding, and hence mediately on conviction itself,
since this depends so much upon the use of the understanding.
As for what concerns in particular the deferral or reservation of our
judgment, however, this consists in the resolution not to let a merely
provisional judgment become determining. A provisional judgment is one in
which I represent that while there are more grounds for the truth of a
thing than against it, these grounds still do not suffice for a determining or
definitive judgment, through which I simply decide for the truth. Provi-
sional judging is thus merely problematic judging with consciousness.
Reservation of judgment can happen for two purposes: either in order to
seek for the grounds of the determining judgment, or in order never to
31
z
"setzen . . . nicht aus einander."
"etwas Willkürliches."
577
75
judge. In the former case the deferral of judgment is called critical
(suspensio judicii indagatoria), in the hiter skeptical (suspensiojudiciisceptica).
For the skeptic refrains from all judgment, while the true philosopher
merely suspends his judgment in case he does not yet have sufficient
grounds for holding something to be true.
To suspend one's judgment in accordance with maxims requires a prac-
ticed faculty of judgment, which is found only in advancing age. In gen-
eral, reservation of our approval is a very hard thing, partly because our
understanding is so desirous of expanding itself and enriching itself with
cognitions" by judging, partly because our inclination is always directed
more toward certain things than toward others. He who has often had to
retract his approval, however, and who has thereby become smart and
cautious, will not give it so quickly, out of fear of having subsequently to
retract his judgment again. This revocation is always mortifying and causes
one to mistrust all other cognitions.*
We observe here further that leaving one's judgment in dubio is some-
thing different from leaving it in suspense. In the latter case I always have
an interest in the thing, in the former it is not always in conformity with
my end and interest to decide whether the thing is true or not.
Provisional judgments are quite necessary, indeed, indispensable, for
the use of the understanding in all meditation and investigation. For they
serve to guide the understanding in its inquiries and to provide it with
various means thereto.
When we meditate concerning an object, we must always judge provi-
sionally and, as it were, get the scent of the cognition that is partly to come
to us through the meditation. And when we go after inventions or discover-
ies, we must always make a provisional plan, otherwise our thoughts go on
at random. We can think of provisional judgments, therefore, as maxims
for the investigation of a thing. We could also call them anticipations,
because we anticipate our judgment of a thing even before we have the
determining judgment. Judgments of this sort have their good use, then,
and rules can even be given for how we ought to judge provisionally
concerning an object.
Prejudices must be distinguished from provisional judgments.
Prejudices are provisional judgments insofar as they are accepted as princi-
ples. Every prejudice is to be regarded as a principle of erroneous judg-
ments, and from prejudices arise not prejudices, but rather erroneous
judgments. Hence one must distinguish the false cognition that arises
from prejudice from its source, the prejudice itself. Thus the interpreta-
"
*
"Kenntnissen."
"Kenntnisse."
tion of dreams, for example, is not in itself a prejudice, but rather an error,
which arises from the assumed universal rule that what happens a few
times happens always, or is always to be held to be true. And this princi-
ple, under which the interpretation of dreams belongs, is a prejudice.
Sometimes prejudices are true provisional judgments; what is wrong is
only that they hold for us as principles or as determining judgments. The
cause of this deception is to be sought in the fact that subjective grounds
are falsely held to be objective, due to a lack of reflection, which must
precede all judging. For even if we can accept some cognitions, e.g.,
immediately certain propositions, without investigating them, i.e., without
examining the conditions of their truth, we still cannot and may not judge
concerning anything without reflecting, i.e., without comparing a cognition
with the power of cognition from which it is supposed to arise (sensibility
or the understanding). If we accept judgments without this reflection,
which is necessary even where no investigation occurs, then from this
prejudices arise, or principles for judging based on subjective causes that
are falsely held to be objective grounds.
The principal sources of prejudices are: imitation, custom, and inclination.
Imitation has a universal influence on our judgments, for there is a
strong ground for holding to be true what others have put forth as true.
Hence the prejudice that what the whole world does is right. As for what
concerns prejudices that have arisen from custom, they can only be rooted
out in the course of time, as the understanding, having little by little been
held up and slowed down in judging by opposing grounds, is thereby
gradually brought to an opposite mode of thought. If a prejudice of cus-
tom has arisen at the same time from imitation, however, then the man
who possesses it is very hard to cure. The inclination toward passive use of
reason, or toward the mechanism of reason rather than toward its spontaneity
under laws, can also be called a prejudice of imitation.
Reason is an active principle, to be sure, which ought not to derive
anything from the mere authority of others, nor even, when its pure use is
concerned, from experience. But the indolence of many men is such that
they prefer to follow in the footsteps of others rather than strain their own
powers of understanding. Men of this sort can only be copies of others,
and if everyone were of this kind, the world would remain eternally in one
and the same place. Hence it is most necessary and important not to
confine youths to mere imitation, as customarily happens.
There are so many things that contribute to accustoming us to the
maxim of imitation, and thereby to making reason a fruitful ground of
prejudices. Such aids to imitation include:
i. Formulas. 0 These are rules whose expression serves as a model for imitation.
They are uncommonly useful, by the way, for making complicated proposi-
tions easier, and the most enlightened mind therefore seeks to discover such
things.
Sayings, whose expression has the great precision of pregnant meaning, so
that it seems one could not capture the sense with fewer words. Pronounce-
ments of this sort (dicta\), which must always be borrowed from others whom
one trusts to have a certain infallibility, serve, on account of this authority, as
rules and as laws. The pronouncements of the bible are called sayings xat'
Sentences/ i.e., propositions which recommend themselves and which,
through the force of the thoughts lying within them, often retain their prestige
through centuries as products of a mature power of judgment.
Canones. These are universal rules 1 that serve as foundations for die sciences
and indicate something sublime and thought through. One can express them
in a sententious way, too, so that they are the more pleasing.
Proverbs (proverbia). These are popular rules of the common understanding,
or expressions for signifying its popular judgments. Since provincial proposi-
tions of this sort serve only the common crowd as sentences and canons, tiiey
are not to be found among people of finer upbringing.
From the three universal sources of prejudices stated above, and from
imitation in particular, many particular prejudices arise, among which we
wish to touch here upon the following as the most common,
i. Prejudices of prestige. Among these are to be reckoned:
a) The prejudice of the prestige of the person. If, in things that rest on
experience and on testimony, we build our cognition on the prestige of
other persons, we are not thereby guilty of any prejudice; for in matters of
this kind, since we cannot experience everything ourselves and compre-
hend it with our own understanding, the prestige of the person must be
the foundation of our judgments. When we make the prestige of others
the ground of our holding-to-be-true in respect of cognitions of reason,
however, we accept these cognitions merely on the basis of prejudice. For
truths of reason hold anonymously; the question here is not, Who said it?
but rather, What did he say? It does not matter at all whether a cognition is
of noble descent; but the inclination toward the prestige of great men is
nonetheless very common, partly because of the restrictedness of our own
insight, partly due to a desire to imitate what is described to us as great.
Added to this is the fact that the prestige of the person serves to flatter our
vanity in an indirect way. Just as the subjects of a powerful despot are
proud of the fact that they are all just treated equally by him, since to this
extent the least can fancy himself the equal of the foremost, as they are
both nothing over against the unrestricted power of their ruler, so too do
the admirers of a great man judge themselves to be equal, insofar as the
superiorities that they may have compared to one another are to be re-
garded as insignificant when considered against the merit of the great
man. For more than one reason, therefore, highly prized great men con-
tribute not a little to the inclination toward the prejudice of the prestige of
the person.
b) The prejudice of the prestige of the multitude. It is principally the crowd
who are inclined to this prejudice. For since they are unable to pass
judgment on the merits, the capabilities, and the cognitions' of the person,
they hold rather to the judgment of the multitude, under the presupposi-
tion that what everyone says must surely be true. This prejudice among
the vulgar relates only to historical matters, however; in matters of reli-
gion, where they are themselves interested, they rely on the judgment of
the learned.
It is in general noteworthy that the ignorant man has a prejudice for
learnedness, while the learned man, on the other hand, has a prejudice for
the common understanding.
If, after a learned man has nearly run the course of the sciences, he
does not procure appropriate satisfaction from all his efforts, he finally
acquires a mistrust of learnedness, especially in regard to those specula-
tions where the concepts cannot be made sensible, and whose foundations
are unsettled, as, e.g., in metaphysics. Since he still believes, however, that
it must be possible to find the key to certainty concerning certain objects
somewhere, he seeks it now in the common understanding, after he had
sought it so long in vain on the path of scientific inquiry.
But this hope is quite deceptive, for if the cultivated faculty of reason
can accomplish nothing in respect to the cognition of certain things, the
uncultivated faculty will certainly do so just as little. In metaphysics the
appeal to pronouncements of the common understanding is completely
inadmissible, because here no case can be exhibited in concreto. With
morals, however, the situation is admittedly different. Not only can all
rules in morals be given in concreto, but practical reason even manifests
itself in general more clearly and more correctly through the organ of the
common use of the understanding than through that of its speculative use.
Hence the common understanding often judges more correctly concern-
ing matters of morality and duty than does the speculative.
c) The prejudice of the prestige of the age. Here the prejudice of antiquity is
one of the most significant. We do have reason to judge kindly of antiquity,
to be sure, but that is only a ground for moderate respect, whose limits we
all too often overstep by treating the ancients as treasurers of cognitions
1
"Kenntnisse."
581
80
81
and of sciences, elevating the relative worth of their writings to an absolute
one and trusting blindly to their guidance. To esteem the ancients so
excessively is to lead the understanding back into its childhood and to
neglect the use of one's own talent. We would also err greatly if we
believed that everyone in antiquity had written as classically as those
whose writings have come down to us. For since time sifts everything and
preserves only what has an inner worth, we may assume, not without
reason, that we only possess the best writings of the ancients.
There are several causes by which the prejudice of antiquity is created
and sustained.
If something exceeds expectation as a universal rule, one initially won-
ders at this, and this wonder then often turns to admiration. This is the
case with the ancients when one finds something in them that, in respect
of the circumstances of time in which they lived, one did not seek. An-
other cause lies in the circumstance that acquaintance with the ancients
and with antiquity proves learnedness and wide reading, which always
brings respect, however common and insignificant in themselves the
things may be that one has drawn from the study of the ancients. A third
cause is the gratitude we owe to the ancients for the fact that they blazed
the path toward many cognitions.-' It seems fair to show them special
esteem, whose measure we often overstep, however. A fourth cause, fi-
nally, is to be sought in a certain envy toward our contemporaries. He who
cannot contend with the moderns extols the ancients at their expense, so
that the moderns cannot raise themselves above him.
The opposite of this is the prejudice of modernity. Sometimes the pres-
tige of antiquity and the prejudice in its favor declined, particularly at the
beginning of this century, when the famous Fontenelle 1 ? took the side of
the moderns. In the case of cognitions that are capable of extension, it is
quite natural that we place more trust in the moderns than in the ancients.
But this judgment has ground only as a mere provisional judgment. If we
make it a determining one, it becomes a prejudice.
2. Prejudices based on self-love or logical egoism, in accordance with which
one holds the agreement of one's own judgment with the judgments of
others to be a dispensable criterion of truth. They are opposed to the
prejudices of prestige, since they express themselves in a certain prefer-
ence for that which is the product of one's own understanding, e.g., one's
own system.
Is it good and advisable to let prejudices stand or even to encourage them? It
is astonishing that in our age such questions can still be advanced, espe-
cially that concerning the encouragement of prejudices. Encouraging some-
'
"Kenntnissen."
one's prejudices amounts to deceiving someone with good intent. It would
be permissible to leave prejudices untouched, for who can occupy himself
with exposing and getting rid of every prejudice? But it is another question
whether it would not be advisable to work toward rooting them out with all
one's powers. Old and rooted prejudices are admittedly hard to battle,
because they justify themselves and are, as it were, their own judges. People
also seek to excuse letting prejudices stand on the ground that disadvan-
tages would arise from rooting them out. But let us always accept these
disadvantages; they will subsequently bring all the more good.
X.
Probability - Explanation of the probable — Distinction between
probability and plausibility - Mathematical and philosophical
probability — Doubt — Subjective and objective doubt - Skeptical,
dogmatic, and critical mode of thought or method of philosophizing -
Hypotheses
To the doctrine concerning the certainty of our cognition pertains also the
doctrine of the cognition of the probable, which is to be regarded as an
approximation to certainty.
By probability is to be understood a holding-to-be-true based on insuffi-
cient grounds which have, however, a greater relation to the sufficient
grounds than do the grounds of the opposite. By this explanation we
distinguish probability* (probabilitas) from mere plausibility' (verisimilitude),
a holding-to-be-true based on insufficient grounds insofar as these are
greater than the grounds of the opposite.
The ground of holding-to-be-true, that is, can be either objectively or
subjectively greater than that of the opposite. Which of the two it is one can
only discover by comparing the grounds of the holding-to-be-true with
the sufficient grounds; for then the grounds of the holding-to-be-true are
greater than the grounds of the opposite can be. With probability, then, the
ground of the holding-to-be-true is objeäively valid, while with mere plau-
sibility it is only subjectively valid. Plausibility is merely quantity of persua-
sion, probability is an approximation to certainty. With probability there
must always exist a standard in accordance with which I can estimate it.
This standard is certainty. For since I am supposed to compare the insuffi-
cient grounds with the sufficient ones, I must know how much pertains to
certainty. Such a standard is lacking, however, with mere plausibility,
since here I do not compare the insufficient grounds with the sufficient
ones, but only with the grounds of the opposite.
*
'
"Wahrscheinlichkeit."
"Scheinbarkeit."
583
The moments of probability can be either homogeneous or heterogeneous.
If diey are homogeneous, as in mathematical cognition, then they must be
enumerated; 1 " if they are heterogeneous, as in philosophical" cognition,
then they must be weighed," i.e., evaluated according to their effect, and
this according to the overpowering of obstacles in the mind. The latter
give no relation to certainty, but only a relation of one plausibility to
another. From this it follows that only the mathematician can determine
the relation of insufficient grounds to the sufficient ground; the philoso-
pher must content himself with plausibility, a holding-to-be-true that is
sufficient merely subjectively and practically. For in philosophical cogni-
tion probability cannot be estimated, on account of the heterogeneity of
the grounds; here the weights are not all stamped, so to speak. Hence it is
only of mathematical probability that one can really say that it is more than
half of certainty.
There has been much talk of a logic of probability (logica probabilium).
But this is not possible; for if the relation of the insufficient grounds to the
sufficient ground cannot be weighed mathematically, then rules do not
help at all. Also, one cannot give any universal rules of probability, except
that error will not occur on one side, but there must rather be a ground of
agreement in the object; likewise, that when there are as many errors, and
errors in equal degree, on two opposed sides, the truth is in the middle.
83
Doubt is an opposing ground or a mere obstacle to holding-to-be-true,
which can be considered either subjeäively or objectively. Doubt is some-
times taken subjeäively, namely, as a condition of an undecided mind, and
objectively as cognition of the insufficiency of the grounds for holding-to-
be-true. In the latter respect it is called an objection, that is, an objective
ground for holding to be false a cognition that is held to be true.
A scruple is a ground opposed to holding-to-be-true that is merely
subjectively valid. In the case of a scruple one does not know whether the
obstacle to holding-to-be-true is grounded objectively or only subjec-
tively, e.g., only in inclination, in custom, etc. One doubts without being
able to explain the ground of the doubt distinctly and determinately and
without being able to have insight into whether this ground lies in the
object itself or only in the subject. Now if it is to be possible to remove
such scruples, then they must be raised to the distinctness and determi-
nateness of an objection. For it is through objections that certainty is
brought to distinctness and completeness, and no one can be certain of a
thing unless opposing grounds have been stirred up, through which it can
"numerirt."
Ak, "philosophischen"; ist ed., "philosophischen."
"ponderirt."
be determined how far one still is from certainty or how close' one is to it.
Also, it is not enough merely to answer each doubt, one must also resolve
it, that is, make comprehensible how the scruple has arisen. If this does
not happen, the doubt is only turned back, but not removed, the seed of the
doubt still remains. In many cases, of course, we do not know whether the
obstacle to holding-to-be-true has only subjective grounds in us or objec-
tive ones, and hence we cannot remove the scruple by exposing the illu-
sion, since we cannot always compare our cognitions with the object but
often only with each other. It is therefore modesty to expound one's
objections only as doubts.
There is a principle of doubting which consists in the maxim that cogni-
tions are to be treated with the intention of making them uncertain and
showing the impossibility of attaining certainty. This method of philoso-
phizing is the skeptical mode of thought, or skepticism. It is opposed to the
dogmatic mode of thought, or dogmatism, which is a blind trust in the
faculty of reason to expand itself a priori through mere concepts, without
critique, merely on account of seeming success.
Both methods are mistaken if they become universal. For there are
many cognitions* in regard to which we cannot proceed dogmatically, and
on the other side skepticism, by renouncing all assertoric cognition/ ruins
all our efforts at attaining possession of a cognition of the certain.
As harmful as this skepticism is, though, the skeptical method is just as
useful and purposeful, provided one understands nothing more by this
than the way of treating something as uncertain and of bringing it to the
highest uncertainty, in the hope of getting on the trail of truth in this way.
This method is thus really a mere suspension of judging. It is quite useful
to the critical procedure, by which is to be understood that method of
philosophizing in accordance with which one investigates the sources of his
claims or objections and the grounds on which these rest, a method which
gives hope of attaining certainty.
In mathematics and physics skepticism does not occur. The only cogni-
tion that can occasion it is that which is neither mathematical nor empiri-
cal, purely philosophical cognition. Absolute skepticism pronounces every-
thing to be illusion. Hence it distinguishes illusion from truth and must
therefore have a mark of the distinction after all, and consequently must
presuppose a cognition of truth, whereby it contradicts itself.
p
*
'
Reading "wie nahe man" for "wie nahe man noch," with Hinske (KI, xlii).
"Kenntnisse."
"behauptende Erkenntniß."
585
84IMMANUEL KANT
85
86
Concerning probability, we observed above that it is merely an approxima-
tion to certainty. Now this is especially the case with hypotheses, through
which we can never attain apodeictic certainty in our cognition, but always
only a greater or lesser degree of probability.
A hypothesis is a holding-to-be-true of the judgment of the truth of a ground
for the sake of its sufficiency for given consequences,^ or more briefly, the
holding-to-be-true of a presupposition as a ground.
All holding-to-be-true in hypotheses is thus grounded on the fact that
the presupposition, as ground, is sufficient to explain other cognitions as
consequences. For we infer here from the truth of the consequence to the
truth of the ground. But since this mode of inference, as already observed
above, yields a sufficient criterion of truth and can lead to apodeictic
certainty only when all possible consequences of an assumed ground are
true, it is clear from this that since we can never determine all possible
consequences, hypotheses always remain hypotheses, that is, presupposi-
tions, whose complete certainty we can never attain. In spite of this, the
probability of a hypothesis can grow and rise to an analogue of certainty,
namely, when all the consequences that have as yet occurred to us can be
explained from the presupposed ground. For in such a case there is no
reason why we should not assume that we will be able to explain all
possible consequences thereby. Hence in this case we give ourselves over
to the hypothesis as if it were fully certain, although it is so only through
induction.
And in every hypothesis something must be apodeictically certain, too,
namely,
1. the possibility of the presupposition itself. If, for example, to explain
earthquakes and volcanoes we assume a subterranean fire, then such a
fire must be possible, if not as a flaming body, yet as a hot one. For the
sake of certain other appearances, however, to make the earth out to be an
animal, in which the circulation of the inner fluids produces warmth, is to
put forth a mere invention and not a hypothesis. For realities may be made
up, but not possibilities; these must be certain.
2. The consequentia. From the assumed ground the consequences
must flow correctly; otherwise the hypothesis becomes a mere chimera.
3. The unity. It is an essential requirement of a hypothesis that it be
only one and that it not need any subsidiary hypotheses for its support. If,
in the case of a hypothesis, we have to have several others to help, then it
thereby loses very much of its probability. For the more consequences that
may be derived from a hypothesis, the more probable it is, the fewer, the
more improbable. Thus Tycho Brake's hypothesis, for example, did not
suffice for the explanation of many appearances; hence he assumed sev-
eral new hypotheses to complete it.'9 Now here it is to be surmised that the
assumed hypothesis cannot be the real ground. The Copernican system,
on the other hand, is an hypothesis from which everything can be ex-
plained that ought to be explained therefrom, so far as it has yet occurred to
us. Here we do not need any subsidiary hypotheses (hypotheses subsidiarias).
There are sciences that do not allow any hypotheses, as, for example,
mathematics and metaphysics. But in the doctrine of nature they are
useful and indispensable.
APPENDIX
Of the distinction between theoretical and practical cognition
A cognition is called practical as opposed to theoretical, but also as opposed
to speculative cognition.
Practical cognitions are, namely, either
1.
2.
imperatives, and are to this extent opposed to theoretical cognitions; or they
contain
the grounds for possible imperatives and are to this extent opposed to speculative
cognitions.
By an imperative is to be understood in general every proposition that
expresses a possible free action, whereby a certain end is to be made real.
Every cognition that contains imperatives is practical, then, and is to be
called practical in opposition to theoretical cognition. For theoretical cogni-
tions are ones that express not what ought to be but rather what is, hence
they have as their object not an acting' but rather a being.'
On the other hand, if we oppose practical to speculative cognitions, then
they can also be theoretical, provided only that imperatives can be derived from
them. Considered in this respect they are then practical as to content (in
potentia) or objectively. By speculative cognitions we understand, namely,
ones from which no rules for proceeding can be derived, or which contain
no grounds for possible imperatives. There is a multitude of such specula-
tive propositions in theology, for example. Speculative cognitions of this
sort are always theoretical, then, but it is not the case, conversely, that
every theoretical cognition is speculative; it can also be at the same time
practical, considered in another respect.
In the end everything comes down to the practical, and the practical
worth of our cognition consists in this tendency of everything theoretical
and all speculation in regard to its use. This worth is unconditioned, how-
ever, only if the end toward which the practical use of the cognition is
directed is an unconditioned end. The sole, unconditioned, and final end
(ultimate end) to which all practical use of our cognition must finally
relate is morality, which on this account we may also call the practical
without qualification or the absolutely practical. And that part of philosophy
'
'
"'kein Handeln?
"ein Sein.'
587
87IMMANUEL KANT
which has morality as its object would accordingly have to be called
practical philosophy xat' e^o/riv;" although every other philosophical
science always has its practical part, i.e., can contain a directive from the
theories advanced for their practical use for the realization of certain ends.
"
par excellence.
588/.
Universal doctrine of elements
FIRST S E C T I O N
Of concepts
§i
The concept in general and its distinction from intuition
All cognitions, that is, all representations related with consciousness to an
object, are either intuitions or concepts. An intuition is a singular" representa-
tion (repraesentatio singularis), a concept a universal (repraesentatio per notas
communes) or reflected" representation (repraesentatio discursiva).
Cognition through concepts is called thought (cognitio discursiva).
Note i. A concept is opposed to intuition, for it is a universal representation, or a
representation of what is common to several objects, hence a representa-
tion insofar as it can be contained in various ones. 20
2. It is a mere tautology to speak of universal or common* concepts - a
mistake that is grounded in an incorrect division of concepts into univer-
sal, particular, and singular. Concepts themselves cannot be so divided,
but only their use.
§2
Matter and form of concepts
With every concept we are to distinguish matter and form. The matter of
concepts is the object, their form universality.
"einzelne."
"reflectirte."
"gemeinsamen."
589
92
§3
Empirical and pure concept
A concept is either an empirical or a pure concept (vel empiricus vel intellec-
tualis). A pure concept is one that is not abstracted- 1 ' from experience but
arises rather from the understanding even as to content.
An idea is a concept of reason whose object simply cannot be met with
in experience.
Note i. An empirical concept arises from the senses through comparison of
objects of experience and attains through the understanding merely the
form of universality. The reality of these concepts rests on actual experi-
ence, from which, as to their content, they are drawn/ But whether there
are pure concepts of the understanding (conceptus pun), which, as such, arise
merely from the understanding, independently of all experience, must be
investigated by metaphysics.
2. Concepts of reason, or ideas, simply cannot lead to actual objects, be-
cause these latter must all be contained in a possible experience. But they
serve to lead the understanding by means of reason in regard to experi-
ence and to the use of its rules in the greatest perfection, or also to show
that not all possible things are objects of experience, and that the princi-
ples of the possibility of the latter do not hold of things in themselves, nor
of objects of experience as things in themselves.
93
An idea contains the archetype for the use of the understanding, e.g., the
idea of the world whole, which idea must necessarily be, not as constitutive
principle for the empirical use of the understanding, but as regulative
principle for the sake of the thoroughgoing connection" of our empirical
use of the understanding. Thus it is to be regarded as a necessary basic
concept, either for objectively completing the understanding's actions of
subordination or for regarding them as unlimited. - The idea cannot be
attained by composition, either, for the whole is prior to the part. There are
ideas, however, to which an approximation occurs. This is the case with
mathematical ideas, or ideas of the mathematical production of a whole, which
differ essentially from dynamical ideas, which are completely heterogeneous
from all concrete concepts, because the whole is different from concrete
concepts not as to quantity (as with the mathematical ideas) but rather as
to kind.
One cannot provide objective reality for any theoretical idea, or prove it,
except for the idea of freedom, because this is the condition of the moral
*
*
*
"abgezogen."
"geschöpft."
"des durchgängigen Zusammenhanges."
law, whose reality is an axiom. The reality of the idea of God can only be
proved by means of this idea, and hence only with a practical purpose, i.e.,
to act as if there is a God, and hence only for this purpose.
In all sciences, above all in those of reason, the idea of the science is its
universal abstract or outline, hence the extension of all the cognitions* that
belong to it. Such an idea of the whole - the first thing one has to look to
in a science, and which one has to seek - is architeäonic, as, e.g., the idea
of jurisprudence/
Most men lack the idea of humanity, the idea of a perfect republic, of a
happy life, etc. Many men have no idea of what they want, hence they
proceed according to instinct and authority.
§4
Concepts that are given (a priori or a posteriori) and concepts
that are made
All concepts, as to matter, are either given (conceptus dati) or made (conceptus
factitii). The former are given either a priori or a posteriori.
All concepts that are given empirically or a posteriori are called concepts of
experience/ all that are given a priori are called notions.
Note.
The form of a concept, as that of a discursive representation, is always
made.
§5
Logical origin of concepts
The origin of concepts as to mere form rests on reflection and on abstrac-
tion from the difference among things that are signified by a certain
representation. And thus arises here the question: Which acts of the under-
standing constitute a concept? or what is the same, Which are involved in the
generation of a concept out of given representations?
Note i. Since universal logic abstracts from all content of cognition through
concepts, or from all matter of thought, it can consider a concept only in
respect of its form, i.e., only subjectively; not how it determines an object
through a mark, but only how it can be related to several objects. Hence
*
'
d
"Kenntnisse."
"Rechtswissenschaft."
"Erfahrungsbegriffe."
591
94IMMANUEL KANT
universal logic does not have to investigate the source of concepts, not how
concepts arise as representations, but merely horv given representations become
concepts in thought; these concepts, moreover, may contain something that
is derived from experience, or something invented, or borrowed from the
nature of the understanding. - This logical origin of concepts - the
origin as to their mere form - consists in reflection, whereby a represen-
tation common to several objects (conceptus communis) arises, as that form
which is required for the power of judgment. Thus in logic only the
difference in reflection in concepts is considered.
The origin of concepts in regard to their matter, according to which a
concept is either empirical or arbitrary or intellectual, is considered in
metaphysics.
§6
Logical Ac tu s of comparison, reflection, and abstraction
The logical actus of the understanding, through which concepts are gener-
ated as to their form, are:
1.
2.
3.
95
comparison 1 of representations among one another in relation to the unity of
consciousness;
reflection*as to how various representations can be conceived in one conscious-
ness; and finally
abstraction 1 of everything else in which the given representations differ.
Note i. To make concepts out of representations one must thus be able to
compare, to reflect, and to abstract, for these three logical operations of the
understanding are the essential and universal conditions for generation
of every concept whatsoever. I see, e.g., a spruce, a willow, and a linden.
By first comparing these objects with one another I note that they are
different from one another in regard to the trunk, the branches, the
leaves, etc.; but next I reflect on that which they have in common
among themselves, trunk, branches, and leaves themselves, and I ab-
stract from the quantity, the figure, etc., of these; thus I acquire a
concept of a tree.
2. The expression abstraction is not always used correctly in logic. We must
not speak of abstracting something (abstrahere aliquid), but rather of ab-
stracting from something (abstrahere ab aliquo). With a scarlet cloth, for
example, if I think only of the red color, then I abstract from the cloth; if I
abstract from this too and think the scarlet as a material stuff in general,
then I abstract from still more determinations, and my concept has in this
'
1
g
"Cotnparation, d.i., die Vergleichung."
"'Reflexion, d.i. die Überlegung."
"-Abstraction oder die Absonderung."
592THE J Ä S C H E L O G I C
way become still more abstract. For the more the differences among
things that are left out of a concept, or the more the determinations from
which we abstract in that concept, the more abstract the concept is.
Abstract concepts, therefore, should really be called abstracting concepts
(conceptus abstrahentes), i.e., ones in which several abstractions occur.
Thus the concept body is really not an abstract concept, for I cannot
abstract from body itself, else I would not have the concept of it. But I
must of course abstract from the size, the color, the hardness or fluidity,
in short, from all the special determinations of particular bodies. The
most abstract concept is the one that has nothing in common with any
distinct from itself. This is the concept of something, for that which is
different from it is nothing, and it thus has nothing in common with
something.
Abstraction is only the negative condition under which universal represen-
tations can be generated, the positive condition is comparison and reflec-
tion. For no concept comes to be through abstraction; abstraction only
perfects it and encloses it in its determinate limits.
§7
Content and extension of concepts
Every concept, as partial concept, is contained in the representation of
things; as ground of 'cognition, i.e., as mark, 1 " these things are contained under
it. In the former respect every concept has a content, in the other an
extension.
The content and extension of a concept stand in inverse relation to one
another. The more a concept contains under itself, namely, the less it
contains in itself, and conversely.
Note.
The universality or universal validity of a concept does not rest on the fact
that the concept is a partial concept, but rather on the fact that it is a ground
of cognition.
§8
Quantity of the extension of concepts
The more the things that stand under a concept and can be thought
through it, the greater is its extension or sphere.
Ak, "'d.i. alsMerkmar; ist ed., "d.i. als Merkmal."
593
96IMMANUEL KANT
Note.
As one says of a ground in general that it contains the consequence under
itself, so can one also say of the concept that as ground of cognition it
contains all those things under itself from which it has been abstracted,
e.g., the concept of metal contains under itself gold, silver, copper, etc. For
since every concept, as a universally valid representation, contains that
which is common to several representations of various things, all these
things, which are to this extent contained under it, can be represented
through it. And it is just this that constitutes the usefulness of a concept.
The more the things that can be represented through a concept, the
greater is its sphere. Thus the concept body, for example, has a greater
extension than the concept metal.
§9
Higher and lower concepts
Concepts are called higher (conceptus superiores) insofar as they have other
concepts under themselves, which, in relation to them, are called lower con-
cepts. A mark of a mark - a remote mark - is a higher concept, the concept
in relation to a remote mark is a lower one.
Note.
Since higher and lower concepts are so called only relatively (respective^,
one and the same concept can, in various relations, be simultaneously a
higher one and a lower one. Thus the concept man is a higher one in
relation to the concept Negro,' but a lower one in relation to the concept
animal.
§10
Genus and species
97
The higher concept, in respect to its lower one, is called genus/ the lower
concept in regard to its higher one species. k
Like higher and lower concepts, genus and species concepts are distin-
guished not as to their nature, then, but only in regard to their relation to
one another (termini a quo or ad quod 1 } in logical subordination.
Ak, "Neger"; ist ed., "Pferd."
"Gattung (genus). "
"An (species)."
terms from which [or] to which.
594THE J Ä S C H E LOGIC
Highest genus and lowest species
The highest genus is that which is not a species (genus summum non est
species), just as the lowest species is that which is not a genus (species, quae
non est genus, est infimd).
In consequence of the law of continuity, however, there cannot be
either a lowest or a next species.
Note.
If we think of a series of several concepts subordinated to one another,
e.g., iron, metal, body, substance, thing, then here we can attain ever
higher genera - for every species is always to be considered at the same
time as genus in regard to its lower concept, e.g., the concept learned man
in regard to the concept philosopher - until we finally come to a genus that
cannot in turn be a species. And we must finally be able to attain such a
one, because in the end there must be a highest concept (conceptus
summus), from which, as such, nothing further may be abstracted without
the whole concept disappearing. - But in the series of species and gen-
era there is no lowest concept (conceptus infimus) or lowest species, under
which no other would be contained, because such a one cannot possibly
be determined. For even if we have a concept that we apply immediately
to individuals, there can still be specific differences in regard to it, which
we either do not note, or which we disregard. Only comparatively for use
are there lowest concepts, which have attained this significance," 1 as it
were, through convention, insofar as one has agreed not to go deeper
here.
In respect to the determination of species and genus concepts, then, the
following universal law holds: There is a genus that cannot in turn be a species,
but there is no species that should not be able in turn to be a genus.
§12
Broader and narrower concept - Convertible concepts
The higher concept is also called a broader concept, the lower concept a
narrower one.
Concepts that have one and the same sphere are called convertible
concepts" (conceptus redproci).
"Bedeutung."
"Wechselbegriffe."
595
98IMMANUEL KANT
§13
Relation of the lower concept to the higher, of the broader
to the narrower
The lower concept is not contained in the higher, for it contains more in
itself than does the higher one; it is contained under it, however, because
the higher contains the ground of cognition of the lower.
Furthermore, one concept is not broader than another because it contains
more under itself - for one cannot know that - but rather insofar as it
contains under itself the other concept and besides this still more.
§i4
Universal rules in respeu of the subordination of concepts
In regard to the logical extension of concepts, the following universal rules
hold:
1.
2.
What belongs to or contradicts higher concepts also belongs to or contradicts
all lower concepts that are contained under those higher ones; and
conversely: What belongs to or contradicts all lower concepts also belongs to
or contradicts their higher concept.
Note.
99
Because that in which things agree flows from their universal properties,
and that in which they are different from one another flows from their
particular properties, one cannot infer that what belongs to or contradicts
one lower concept also belongs to or contradicts other lower concepts,
which belong with it to one higher concept. Thus one cannot infer, e.g.,
that what does not belong to man does not belong to angels either.
§15
Conditions for higher and lower concepts to arise: Logical abstraction
and logical determination
Through continued logical abstraction higher and higher concepts arise,
just as through continued logical determination, on the other hand, lower
and lower concepts arise. The greatest possible abstraction yields the
highest or most abstract concept - that from which no determination can
be further thought away. The highest, completed determination would
yield a thoroughly determinate concept (conceptus omnimode determinate),
i.e., one to which no further determination might be added in thought.
596THE J Ä S C H E LOGIC
Note.
Since only individual things," or individuals/ are thoroughly determinate,
there can be thoroughly determinate cognitions only as intuitions, but not
as concepts; in regard to the latter, logical determination can never be
regarded as completed (§ n, Note).
§16
Use of concepts in a b s t r a c t o and in concreto
Every concept can be used universally or particularly (in abstracto or in
concreto). The lower concept is used in abstracto in regard to its higher one,
the higher concept in concreto in regard to its lower one.
Note i. Thus the expressions abstract and concrete relate not to concepts in
themselves - for every concept is an abstract concept - but rather only to
dieir use. And this use can in turn have various degrees, accordingly as
one treats a concept more or less abstractly or concretely, i.e., as one
either leaves aside or adds more or fewer determinations. Through ab-
stract use a concept comes closer to the highest genus, through concrete
use, on die other hand, to the individual.
2. Which use of concepts, the abstract or die concrete, has an advantage
over the odier? Nodiing can be decided about this. The worth of the one
is not to be valued less than the worth of the other. Through very abstract
concepts we cognize little in many things, through very concrete concepts
we cognize much m few tilings; what we win on the one side, then, we lose
again on the other. A concept that has a large sphere is very useful insofar
as one can apply it to many things; but in return for that, there is that
much less contained in it. In the concept substance, for example, I do not
tiiink as much as in the concept chalk.
j. The art of popularity consists in finding die relation between representa-
tion in abstracto and in concreto in die same cognition, hence between
concepts and dieir exhibition, through which the maximum of cognition
is achieved, both as to extension and as to content.
SECOND SECTION
Of judgments
§17
Definition* of a judgment in general
A judgment is the representation of the unity of the consciousness of
various representations, or the representation of their relation insofar as
they constitute a concept.
°
p
*
"einzelne Dinge."
"Individuen."
"Erklärung."
597
IOO
101IMMANUEL KANT
§18
Matter and form of judgments
Matter and form belong to every judgment as essential constituents of it.
The matter of the judgment consists in the given representations that are
combined in the unity of consciousness in the judgment, the form in the
determination of the way that the various representations belong, as such,
to one consciousness.
§19
Object of logical reflection the mere form of judgments
Since logic abstracts from all real or objective difference of cognition, it
can occupy itself as litde with the matter of judgments as with the content
of concepts. Thus it has only the difference among judgments in regard to
their mere form to take into consideration.
102
§20
Logical forms of judgments: Quantity, quality, relation, and modality
The distinctions among judgments in respect of their form may be traced
back to the four principal moments of quantify, quality, relation, and modal-
ity, in regard to which just as many different kinds of judgments are
determined.
§21
Quantity of judgments: Universal, particular, singular
As to quantity, judgments are either universal or particular or singular,
accordingly as the subject is either wholly mcluded in or «-eluded from the
notion of the predicate or is only in part /«eluded in or excluded from it. In
the universal judgment, the sphere of one concept is wholly enclosed
within the sphere of another; in the particular, a part of the former is
enclosed under the sphere of the other; and in the singular judgment,
finally, a concept that has no sphere at all is enclosed, merely as part then,
under the sphere of another.
598THE J Ä S C H E LOGIC
Note i. As to logical form, singular judgments are to be assessed as like universal
ones in use, for in both the predicate holds of the subject without excep-
tion. In the singular proposition, Caius is mortal, for example, there can
just as little be an exception as in the universal one, All men are mortal.
For there is only one Caius.
2. In respect of the universality r of a cognition, there exists a real distinction
between general 1 and universal' propositions, which of course does not
have to do with logic, however. General propositions, namely, are those
that contain merely something of the universal regarding certain objects,
and consequently do not contain sufficient conditions of subsumption,
e.g., the proposition that one must make proofs thorough. Universal"
propositions are those that maintain something universally 1 ' of an object.
j. Universal rules are either analytically or synthetically universal. The former
abstract from differences, the latter attend to distinctions and conse-
quently determine in regard to them too. The more simply an object is
thought, the more possible is analytic universality in consequence of a
concept.
4. If we cannot have insight into universal propositions in their universality
without cognizing 01 them in concreto, then they cannot serve as a standard
and hence cannot hold heuristically in application, but are only assignments
to investigate the universal grounds for that with which we first became
acquainted in particular cases. For example, the proposition, He who has no
interest in lying and who knows the truth will speak the truth; we cannot have
insight into this proposition in its universality, because we are acquainted
with the restriction to the condition of the disinterested - namely, that
men can lie out of interest, which comes because they do not hold fast to
morality - only through experience. An observation that teaches us the
weakness of human nature.
5. Of particular judgments it is to be noted that if it is to be possible to have
insight into them through reason, and hence for them to have a rational,
not merely intellectual (abstracted) form, then the subject must be a
broader concept (conceptus latior) than the predicate. Let the predicate
always = O> the subject Q, then
is a particular judgment, for some of what belongs under a is b, some not
b - that follows from reason. But let it be
"Allgemeinheit."
"generalen."
"universalen."
"Universale."
"allgemein."
"kennen."
599
103IMMANUEL KANT
then at least all a can be contained under b, if it is smaller, but not if it is
greater, hence it is particular only by accident.
§22
Quality of judgments: Affirmative, negative, infinite
104
As to quality, judgments are either affirmative or negative or infinite. In the
affirmative judgment the subject is thought under the sphere of a predicate,
in the negative it is posited outside the sphere of the latter, and in the infinite
it is posited in the sphere of a concept that lies outside the sphere of
another.
Note i. The infinite judgment indicates not merely that a subject is not contained
under the sphere of a predicate, but that it lies somewhere in the infinite
sphere outside its sphere; consequently this judgment represents the
sphere of the predicate as restricted.
Everything possible is either^ or nonA. If I say, then, something is non
A, e.g., the human Soul is non-mortal, some men are non-learned, etc.,
then this is an infinite judgment. For it is not thereby determined, con-
cerning the finite sphere A, under which concept the object belongs, but
merely that it belongs in the sphere outside A, which is really no sphere at
all but only a sphere's sharing of a limit with the infinite* or the limiting
itself. y Now although exclusion is a negation, the restriction 2 of a concept
is still a positive act. Therefore limits" are positive concepts of restricted*
objects.
2. According to the principle of the excluded middle' (exclusi tertii), the
sphere of one concept relative to another is either exclusive or inclusive.
Now since logic has to do merely with the form of judgment, not with
concepts as to their content, the distinction of infinite from negative
judgments is not proper to this science.
j. In negative judgments the negation always affects the copula; in infinite
ones it is not the copula but rather the predicate that is affected, which
may best be expressed in Latin.
* "die Angrenzung einer Sphäre an das Unendliche."
y "die Begrenzung selbst."
'
*
* "Beschränkung."
"Grenzen."
"beschränkter."
' "Principium der Ausschließung jedes Dritten."
600THE J Ä S C H E LOGIC
Relation of judgments: Categorical, hypothetical, disjunctive
As to relation, judgments are either categorical or hypothetical or disjunctive.
The given representations in judgment are subordinated one to another
for the unity of consciousness, namely, either as predicate to subjert, or as
consequence to ground, or as member of the division to the divided concept.
Through the first relation categorical judgments are determined, through
the second hypothetical, and through the third disjunctive.
§24
Categorical judgments
In categorical judgments, subject and predicate constitute their matter;
the form, through which the relation (of agreement or of opposition)
between subject and predicate is determined and expressed, is called the
copula^.
Note.
Categorical judgments constitute the matter of the remaining judgments,
to be sure, but one must not on this account believe, as several logicians
do, that both hypothetical and disjunctive judgments are nothing more
than various clothings of categoricals and hence may be wholly traced back
to these latter. All three kinds of judgments rest on essentially different
logical functions of the understanding and must therefore be considered
according to their specific difference.
§25
Hypothetical judgments
The matter of hypothetical judgments consists of two judgments that are
connected with one another as ground and consequence. One of these
judgments, which contains the ground, is the antecedent 11 (antecedens, prius),
the other, which is related to it as consequence, is the consequent'
(consequens, posterius), and the representation of this kind of connection of
two judgments to one another for the unity of consciousness is called the
consequentia, which constitutes iheform of hypothetical judgments.
d
'
"Vordersatz."
"Nachsatz."
601
105IMMANUEL KANT
106
Note i. What the copula L is for categorical judgments, then, the consequentia is for
hypotheticals - their form.
2. Some believe it is easy to transform a hypothetical proposition into a
categorical. But this will not do, because the two are wholly different from
one another as to their nature. In categorical judgments nothing is problem-
atic, rather, everything is assertoric, but in hypotheticals only the conse-
quentia is assertoric. In the latter I can thus connect two false judgments
with one another, for here it is only a matter of the correctness of the
connection - the form of the consequentia, on which the logical truth of
these judgments rests. There is an essential difference between the two
propositions, All bodies are divisible, and, If all bodies are composite, then
they are divisible. In the former proposition I maintain the thing directly, in
the latter only under a condition expressed problematically.
§26
Modes of connection in hypothetical judgments: m o d u s p o n e n s
and m o d u s t o l l e n s
The form of the connection in hypothetical judgments is of two kinds: the
positing^ (modus ponens) and the denying^ (modus tollens).
1.
2.
If the ground (antecedens) is true, then the consequence (consequens) deter-
mined by it is true too; called modus ponens.
If the consequence (consequens) is false, then the ground (antecedens) is false
too; modus tollens.
§2?
Disjunaive judgments
A judgment is disjunctive if the parts of the sphere of a given concept
determine one another in the whole or toward a whole* as complements
(complementa).
§28
Matter and form of disjunctive judgments
The several given judgments of which the disjunctive judgment is com-
posed constitute its matter and are called the members of the disjunction or
opposition. The form of these judgments consists in the disjunction itself,
f
s
*
"setzende."
"aufhebende?
"in dem Ganzen oder zu einem Ganzen."
i.e., in the determination of the relation of the various judgments as
members of the whole sphere of the divided cognition which mutually
exclude one another and complement one another.
Note.
Thus all disjunctive judgments represent various judgments as in the com-
munity of a sphere and produce each judgment only through the restriction
of the others in regard to the whole sphere; they determine each judg-
ment's relation to the whole sphere, then, and thereby at the same time the
relation that these various members of the division (membra disjuncta) have
among themselves. Thus one member determines every other here only
insofar as they stand together in community as parts of a whole sphere of
cognition, outside of which, in a certain relation, nothing may be thought.
107
§29
Peculiar character of disjunctive judgments
The peculiar character of all disjunctive judgments, whereby their specific
difference from others, in particular from categorical judgments, is deter-
mined as to the moment of relation, consists in this: that the members of
the disjunction are all problematic judgments, of which nothing else is
thought except that, taken together as parts of the sphere of a cognition,
each the complement of the other toward the whole (complementum ad
totum), they are equal to the sphere of the first. And from this it follows
that in one of these problematic judgments the truth must be contained
or - what is the same - that one of them must hold assertorically, because
outside of them the sphere of the cognition includes nothing more under
the given conditions, and one is opposed to the other, consequently nei-
ther something outside them nor more than one among them can be true.
Note.
In a categorical judgment the thing whose representation is considered as a
part of the sphere of another, subordinated representation is considered as
contained under this, its higher concept; thus here, in the subordination of
the spheres, the part of the part is compared with the whole. But in
disjunctive judgments I go from the whole to all the parts taken together.
What is contained under the sphere of a concept is also contained under
one of the parts of this sphere. Accordingly, the sphere must first be
divided. If, for example, I make the disjunctive judgment, A learned man is
learned either historically or in matters of reason, I thereby determine that
these concepts are, as to the sphere, parts of the sphere of the learned man,
but not in any way parts of one another, and that taken together they are all
complete.
The following schema of comparison between categorical and disjunc-
603
108IMMANUEL KANT
live judgments may make it more intuitive that in disjunctive judgments the
sphere of the divided concept is not considered as contained in the sphere
of the divisions, but rather that which is contained under the divided
concept is considered as contained under one of the members of the
division.
In categorical judgments x, which is contained under b, is also under a:
In disjunctive ones x, which is contained under a, is contained either under
b or c, etc.:
Thus the division in disjunctive judgments indicates the coordination not
of the parts of the whole concept, but rather all the parts of its sphere.
Here I think many things through one concept, there one thing through' many
concepts, e.g., the definitum through all the marks of coordination.
§3°
Modality of judgments: Problematic, assertoric, apodeictic
As to modality, through which moment the relation of the whole judgment
to the faculty of cognition is determined, judgments are either problematic
or assertoric or apodeictic. The problematic ones are accompanied with the
consciousness of the mere possibility of the judging, the assertoric ones
with the consciousness of its actuality, the apodeictic ones, finally, with the
consciousness of its necessity.
109
Note i. This moment of modality indicates, then, only the way in which some-
thing is maintained or denied in judgment: whether one does not settle
anything concerning the truth or untruth of a judgment, as in the prob-
lematic judgment, The soul of man may be immortal; or whether one
determines something concerning it, as in the assertoric judgment, The
human soul is immortal; or, finally, whether one even expresses the truth
of a judgment with the dignity of necessity, as in the apodeictic judgment,
Ak, "durch"; ist ed., "durch."
604THE J Ä S C H E LOGIC
The soul of man must be immortal. This determination of merely possi-
ble or actual or necessary truth concerns only the judgment itself, then, not
in any way the thing about which we judge.
2. In problematic judgments, which one can also explain as ones whose
material is given with the possible relation between predicate and subject,
the subject must always have a smaller sphere than the predicate.
j. On the distinction between problematic and assertoric judgments rests
the true distinction between judgments and propositions, which is custom-
arily placed, wrongly, in the mere expression through words, without
which one simply could not judge at all. In judgment the relation of
various representations to the unity of consciousness is thought merely as
problematic, but in a proposition as assertoric. A problematic proposition
is a contradictio in adjecto. Before I have a proposition I must first judge;
and I judge about much that I cannot decide, which I must do, however,
as soon as I determine a judgment as a proposition. It is good, by the way,
first to judge problematically, before one accepts the judgment as asser-
toric, in order to examine it in this way. Also, it is not always necessary to
our purpose to have assertoric judgments.
§31
Exponible judgments
Judgments in which an affirmation and a negation are contained simulta-
neously, but in a covert way, so that the affirmation occurs distinctly but
the negation covertly, are exponible propositions.
Note.
In the exponible judgment, Few men are learned, for example lies (i.), but
in a covert way, the negative judgment, Many men are not learned, and (2.)
the affirmative one, Some men are learned. Since the nature of exponible
propositions depends merely on conditions of language, in accordance
with which one can express two judgments briefly at once, the observation
that in our language there can be judgments that must be expounded
belongs not to logic but to grammar.
§32
Theoretical and practical propositions
Those propositions that relate to the object and determine what belongs
or does not belong to it are called theoretical; practical propositions, on the
other hand, are those diat state the action whereby, as its necessary condi-
tion, an object becomes possible.
605
noIMMANUEL KANT
Note.
Logic has to deal only with propositions that are practical as to form, which
are to this extent opposed to theoretical ones. Propositions that are practical
as to content, and to this extent are distinct from speculative ones, belong to
morals.
§33
Indemonstrable and demonstrable propositions
Demonstrable propositions are those that are capable of a proof; those not
capable of a proof are called indemonstrable.
Immediately certain judgments are indemonstrable and thus are to be
regarded as elementary propositions.
§34
Principles
Immediately certain judgments a priori can be called principles/ insofar as
other judgments are proved from them, but they themselves cannot be
subordinated to any other. On this account they are also called principles' 1
(beginnings').
§35
Intuitive and discursive principles: Axioms and acroamata
Principles are either intuitive or discursive. The former can be exhibited in
intuition and are called axioms (axiomata), the latter may be expressed only
through concepts and can be called acroamata.
in
§36
Analytic and synthetic propositions
Propositions whose certainty rests on identity of concepts (of the predicate
with the notion of the subject) are called analytic propositions." 1 Proposi-
tions whose truth is not grounded on identity of concepts must be called
synthetic.
'
*
'
™
"Grundsätze."
"Prinzipien."
"Anfänge."
Ak, "Sätze"; ist ed., "Sätze."
606THE J Ä S C H E L O G I C
Note I.
2.
An example of an analytic proposition is, To everything x, to which the
concept of body (a + b) belongs, belongs also extension (b)
An example of a synthetic proposition is, To everything x, to which the
concept of body (a+b) belongs, belongs also attraction (c). Synthetic
propositions increase cognition materialiter, analytic ones merely for-
maliter. The former contain determinations (determinationes), the latter
only logical" predicates.
Analytic principles are not axioms, for they are discursive. And even
synthetic principles are axioms only if they are intuitive.
§37
Tautological propositions
The identity of the concepts in analytic judgments can be either explicit
(explicita) or non-explicit (implicita). In the first case the analytic proposi-
tions are tautological.
Note i.
2.
Tautological propositions are empty virtualiter, or empty of consequences, 0
for they are without value or use. The tautological proposition, Man is
man, is of this sort, for example. For if I do not know anything more to say
of man except that he is a man, then I know nothing more of him at all.
Propositions that are identical implicite, on the other hand, are not
empty of consequences or fruitless, for they make clear the predicate that
lay undeveloped (implicite) in the concept of the subject through develop-
ment 11 (explicatio).
Propositions that are empty of consequences must be distinguished from
ones that are empty of sense, 9 which are empty in meaning' because they
concern the determination of so-called hidden properties (qualitates oc-
cultae).
§ 3 8
Postulate and problem
A postulate is a practical, immediately certain proposition, or a principle
that determines a possible action, in the case of which it is presupposed
that the way of executing it is immediately certain.
Problems (problemata) are demonstrable propositions that require a direc-
" Ak, "logische"; ist ed., "logische."
° "folgeleer."
"Entwickelung."
p
* "sinnleeren."
' "leer an Verstand."
607
tive, or ones that express an action, the manner of whose execution is not
immediately certain.
Note i.
2.
There can also be theoretical postulates on behalf of practical reason.
These are theoretical hypotheses that are necessary for a practical pur-
pose of reason, like those of the existence of God, of freedom, and of
another world.
A problem involves (i.) the question, 1 which contains what is to be accom-
plished, (2.) the resolution, which contains the way in which what is to be
accomplished can be executed, and (3.) the demonstration that when I
have proceeded thus, what is required will occur.
§39
Theorems, corollaries, lemmas, and scholia
Theorems are theoretical propositions that are capable of and require
proof. Corollaries are immediate consequences from a preceding proposi-
tion. Propositions that are not indigenous to the science in which they are
presupposed as proved, but rather are borrowed from other sciences, are
called lemmas (lemmata). Scholia, finally, are mere elucidative propositions,'
which thus do not belong to the whole of the system as members.
Note.
Essential and universal moments of every theorem are the thesis and the
demonstration. Furthermore, one can place the distinction between theo-
rems and corollaries in the fact that the latter are inferred immediately, but
the former are drawn through a series of consequences from immediately
certain propositions.
114
§40
Judgments of perception and of experience
A judgment of perception is merely subjective, an objective judgment from
perceptions is a judgment of experience.
Note.
'
'
A judgment from mere perceptions is really not possible, except through
the fact that I express my representation as perception: I, who perceive a
"Quästion."
"Erläuterungssätze."
608THE J Ä S C H E LOGIC
tower, perceive in it the red color. But I cannot say: It is red. For this would
not be merely an empirical judgment, but a judgment of experience, i.e., an
empirical judgment through which I get a concept of the object. E.g., In
touching the stone I sense warmth, is a judgment of perception: but on the
other hand, The stone is warm, is a judgment of experience. It pertains to
the latter that I do not reckon to the object what is merely in my subject, for
a judgment of experience is perception from which a concept of the object
arises; e.g., whether points of light move on the moon or in the air or in my
eye.
THIRD SECTION
Of inferences
§4i
Inference in general
By inferring is to be understood that function of thought whereby one
judgment is derived from another. An inference is thus in general the
derivation of one judgment from the other.
§42
Immediate and mediate inferences
All inferences are either immediate or mediate.
An immediate inference (consequentia immediata) is the derivation (de-
ductio) of one judgment from the other without a mediating judgment
(judicium intermedium). An inference is mediate if, besides the concept that
a judgment contains in itself, one needs still others in order to derive a
cognition therefrom.
§43
Inferences of the understanding, inferences of reason, and inferences of
the power of judgment
Immediate inferences are also called inferences of the understanding; all
mediate inferences, on the other hand, are either inferences of reason or
inferences of the power of judgment. We deal here first with immediate infer-
ences, or inferences of the understanding.
609
115
I. INFERENCES OF THE UNDERSTANDING
§44
Peculiar nature of inferences of the understanding
The essential character of all immediate inferences and the principle of
their possibility consists simply in an alteration of the mere form of judg-
ments, while the matter of the judgments, the subject and predicate,
remains unaltered, the same.
Note i.
2.
By virtue of the fact that in immediate inferences only the form of judg-
ments is altered and not in any way the matter, these inferences differ
essentially from all mediate inferences, in which the judgments are distinct
as to matter too, since here a new concept must be added as mediating
judgment or as middle concept (terminus medius) in order to deduce the
one judgment from the other. If I infer, e.g., All men are mortal, hence
Caius is mortal too, this is not an immediate inference. For here I need for
the deduction the mediating judgment, Caius is a man; through this new
concept, however, the matter of the judgments is altered.
With inferences of the understanding ajudidum intermedium may also be
made, to be sure, but then this mediating judgment is merely tautological.
As, for example, in the immediate inference, All men are mortal, some
men are men, hence some men are mortal, the middle concept is a tauto-
logical proposition.
§45
M o d i of inferences of the understanding
Inferences of the understanding run through all the classes of the logical
functions of judgment and consequently are determined in their principal
kinds through the moments of quantity, quality, relation, and modality. On
this rests the following division of these inferences.
116
§46
i. Inferences of the understanding (in relation to the quantity of
judgments) per judicia subalternata
In inferences of the understanding per judicia subalternata the two judg-
ments are distinct as to quantity, and here the particular judgment is
derived from the universal in consequence of the principle: The inference
from the universal to the particular is valid (ab universali ad particulare valet
consequentia).
Note.
Ajudidum is called subaltematum insofar as it is contained under the other,
as, e.g., particular judgments are under universal ones.
§47
2. Inferences of the understanding (in relation to the quality of
judgments) per judicia opposita
In inferences of the understanding of this kind, the alteration concerns
the quality of the judgments, and this considered in relation to opposition.
Now since this opposition can be of three kinds, the following particular
division of immediate inference results: through contradictorily opposed,
through contrary, and subcontrary judgments."
Note.
Inferences of the understanding through equivalent" judgments (judicia
aequipollentia) cannot really be called inferences, for there is no conse-
quence here; they are rather to be regarded as a mere substitution of words
that signify one and the same concept, where the judgments themselves
also remain unaltered as to form. E.g., Not all men are virtuous, and, Some
men are not virtuous. The two judgments say one and the same thing.
§48
a. Inferences of the understanding per judicia contradictorie
opposita
In inferences of the understanding through judgments which are opposed
to one another contradictorily and which, as such, constitute genuine,
pure opposition, the truth of one of the contradictorily opposed judgments
is deduced from the falsehood of the other, and conversely. For genuine
opposition, which occurs here, contains no more and no less than what
belongs to opposition. In consequence of the principle of the excluded mid-
dle," the two contradicting judgments cannot both be true, and just as little
can they both be false. If the one is true, then the other is false, and
conversely.
" Ak, "durch conträre und durch subconträre Urtheile"; ist ed., "durch conträre und
subconträre Urtheile."
1
"gleichgeltende."
" "Princip des ausschließenden Dritten."
611
§49
b. Inferences of the understanding per judicia* contrarie opposita
Contrary or conflicting judgments (judicia contrarie opposita) are judg-
ments of which the one is universally affirmative, the other universally
negative. Now since one of them says more than the other, and since in
the excess - that it says more than the mere negation of the other -
there can lie falsehood, they cannot both be true, of course, but they can
both be false. In regard to these judgments, therefore, only the inference
from the truth of the one to the falsehood of the other holds, hut not conversely.
§50
c. Inferences of the understanding per judicia subcontrarie opposita
Subcontrary judgments are ones of which the one affirms or denies in
particular (particulariter) what the other denies or affirms in particular.
Since they can both be true but cannot both be false, only the following
inference holds in regard to them: If one of these propositions is false, the other
is true, but not conversely.
Note.
In the case of subcontrary judgments there is no pure, strict opposition, for
what is affirmed or denied in the one is not denied or affirmed of the same
objects in the other. In the inference, e.g., Some men are learned, hence
some men are not learned, what is denied in the second judgment is not
affirmed of the same men in the first.
118
§51
3. Inferences of the understanding (with respect to the relation of
judgments) per judicia conversa sive per conversionem
Immediate inferences through conversion concern the relation of judg-
ments and consist in the transposition of subject and predicate in the two
judgments, so that the subject of the one judgment is made the predicate
of the other judgment, and conversely.
Ak, "judica"; ist ed., "judicia."
612THE J Ä S C H E LOGIC
§52
Pure and altered conversion
In conversion the quantity of the judgments is either altered or it remains
unaltered. In the first case the converted (conversum) is distinct from what
converts (convertens) as to quantity, and the conversion is said to be altered 3 '
(conversio per accidens); in the latter case the conversion is called pure
(conversio simpliciter talis).
§53
Universal rules of conversion
With respect to inferences of the understanding through conversion, the
following rules hold:
1.
2.
3.
Universal affirmative 2 judgments may be converted only per accidens', for the
predicate in these judgments is a broader concept and thus only some of it is
contained in the concept of the subject.
But all universal negative" judgments may be converted simpliciter, for here the
subject is removed from the sphere of the predicate. Thus too, finally,
All particular affirmative propositions may be converted simpliciter; for in these
judgments a part of the sphere of the subject has been subsumed under the
predicate, hence a part of the sphere of the predicate may be subsumed under
the subject.
Note i.
2.
In universal affirmative judgments the subject is considered as a con-
tentum of the predicate, since it is contained under its sphere. Therefore I
may only infer, e.g., All men are mortal, hence some of those who are
contained under the concept mortal are men. The cause of the fact that
universal negative judgments may be converted simpliciter is that two
concepts that contradict one another universally contradict each other in
the same extension.
Some universal affirmative judgments may be converted simpliciter, to be
sure. But the ground of this lies not in their form but in the particular
character of their matter; as, e.g., the two judgments, Everything unalter-
able is necessary, and everything necessary is unalterable.
Ak, "veränderte"; ist ed., "veränderte."
Ak, "Allgemein bejahende"; ist ed., "Allgemein bejahende."
Ak, "allgemein verneinenden"; ist ed., "allgemein verneinenden."
613
119IMMANUEL KANT
§54
4. Inferences of the understanding (in relation to the modality of
judgments) per judicia contraposita
The immediate mode of inference through contraposition consists in that
transposition (metathesis^ of judgments in which merely the quantity re-
mains the same while the quality is altered. They concern only the modal-
ity of judgments, since they transform an assertoric into an apodeictic
judgment.
§55
Universal rule of contraposition
With respect to contraposition, the universal rule holds:
All universal affirmative judgments may be contraposed simpliciter. For if
the predicate, as that which contains the subject under itself, is denied,
and hence the whole sphere, then a part of it must also be denied, i.e., the
subject.
Note i. The metathesis^ of judgments through conversion and that through
contraposition are thus opposed to one another to the extent that the
former alters merely quantity, the latter merely quality.
2. The immediate modes of inference mentioned relate merely to categorical
judgments.
120
II. I N F E R E N C E S OF REASON
§56
Inference of reason in general
An inference of reason is the cognition of the necessity of a proposition
through the subsumption of its condition under a given universal rule.
§57
Universal principle of all inferences of reason
The universal principle on which the validity of all inference through
reason rests may be determinately expressed in the following formula:
What stands under the condition of a rule also stands under the rule
itself.
Note.
The inference of reason premises a universal rule and a subsumption under
its condition. Through this one cognizes the conclusion a priori, not in the
individual, but as contained in the universal and as necessary under a
certain condition. And this, that everything stands under the universal and
is determinable in universal rules, is just the principle of rationality or of
necessity (principium rationalitatis sive necessitatis).
§58
Essential components of the inference of reason
To every inference of reason belong the following essential three parts:
1.
2.
3.
a universal rule, which is called the major proposition 11 (propositio major),
the proposition which subsumes a cognition under the condition of the univer-
sal rule, and which is called the minor proposition' (propositio minor), and finally
the proposition which affirms or denies the rule's predicate of the subsumed
cognition: the conclusion 11 (conclusio).
The first two propositions, in their combination with one another, are
called premises. 1
Note.
A rule is an assertion under a universal condition. The relation of the
condition to the assertion, namely, how the latter stands under the former,
is the exponent^of the rule.
The cognition that the condition (somewhere) exists is the subsumption.
The combination of that which is subsumed under the condition with
the assertion of the rule is the inference.
"Obersatz."
"Untersatz."
"Schlußsatz."
"Vordersätze oder Prämissen."
"Exponent."
615
§59
Matter and form of inferences of reason
The matter of inferences of reason consists in the antecedent proposi-
tions or premises, the form in the conclusion insofar as it contains the
consequentia.
Note i.
2.
In every inference of reason, then, the truth of the premises is to be
examined first, and then the correctness of the consequentia. In rejecting
an inference of reason, one must never reject the conclusion first, but
rather one must first reject either the premises or the consequentia.
In every inference of reason the conclusion is given as soon as the
premises and the consequentia are given.
§60
Division of inferences of reason (as to relation) into categorical,
hypothetical, and disjunctive
122
All rules (judgments) contain objective unity of consciousness of the
manifold of cognition, hence a condition under which one cognition be-
longs with another to one consciousness. Now only three conditions of
this unity may be thought, however, namely: as subject of the inherence of
marks, or as ground of the dependence of one cognition on another, or,
finally, as combination of parts in a whole (logical division). Consequently
there can only be just as many kinds of universal rules (propositiones ma-
jores), through which the consequentia of one judgment from another is
mediated.
And on this is grounded the division of all inferences of reason into
categorical, hypothetical, and disjunäive.
Note i.
2.
Inferences of reason can be divided neither as to quantity, for every major L
is a rule, hence something universal; nor in regard to quality, for it is
equivalent whether the conclusion is affirmative or negative; nor, finally,
in respect of modality, for the conclusion is always accompanied with the
consciousness of necessity and consequently has the dignity of an
apodeictic proposition. Thus only relation remains as the sole possible
ground of division of inferences of reason.
Many logicians hold only categorical inferences of reason to be ordinary,
the others to be extraordinary. But this is groundless and false. For all
three of these kinds are products of equally correct, though equally
essentially different functions of reason.
§61
Peculiar difference between categorical, hypothetical and disjunctive
inferences of reason
The distinguishing feature among the three mentioned kinds of infer-
ences of reason lies in the major premise. In categorical inferences of reason
the major L is a categorical proposition, in hypothetical ones it is a hypotheti-
cal or problematic proposition, and in disjunctive ones it is a disjunctive
proposition.
§62
i. Categorical inferences of reason
In every categorical inference of reason there are three principal concepts 1
(termini), namely:
1.
2.
3.
the predicate in the conclusion, which concept is called the major concept 11
(terminus major), because it has a larger sphere than the subject,
the subject (in the conclusion), whose concept is called the minor concept'
(terminus minor), and
a mediating mark (nota intermedia), which is called the middle concept' (terminus
medius), because through it a cognition is subsumed under the condition of
the rule.
Note.
This distinction among the mentioned termini exists only in categorical
inferences of reason, because these alone infer through a terminus medius;
the others, on the other hand, infer only through the subsumption of a
proposition represented problematically in the m ajo r L and assertorically in
the minor L .
§6 3
Principle of categorical inferences of reason
The principle on which the possibility and validity of all categorical infer-
ences of reason rests is this:
What belongs to the mark of a thing belongs also to the thing itself; and what
*
*
'
>
"Hauptbegriffe."
"Oberbegriff."
"Unterbegriff."
"Mittelbegriff."
617
contradicts the mark of a thing contradicts also the thing itself (nota notae est nota
rei ipsius; repugnans notae, repugnat rei ipsi).
Note.
From the principle just set forth the so-called dictum de omni et nullo may
easily be deduced, and hence it can hold as the first principle neither for
inferences of reason in general nor for the categorical in particular.
Genus and species concepts are universal marks of all things that stand
under these concepts. Accordingly, the rule holds here: What belongs to or
contradicts the genus or species belongs to or contradicts all the objects that are
contained under that genus or species. And this rule is called just the dictum
de omni et nullo.
§64
Rules for categorical inferences of reason
From the nature and the principle of categorical inferences of reason flow
the following rules for them:
124
i.
2.
3.
4.
5.
6.
7.
In no categorical inference of reason can either more or fewer than three
principal concepts (termini) be contained; for here I am supposed to combine
two concepts (subject and predicate) through a mediating mark.
The antecedent propositions or premises may not be wholly negative (expuris
negatives nihil sequitur*); for the subsumption in the minor premise must be
affirmative, as that which states that a cognition stands under the condition of
the rule.
The premises may also not be wholly particular 1 propositions (ex puris particu-
laribus nihil sequitur™); for then there would be no rule, i.e., no universal
proposition, from which a particular cognition could be deduced.
The conclusion always follows the weaker pan of the inference, i.e., the negative and
the particular proposition in the premises, as that which is called the weaker
part of the categorical inference of reason (conclusio sequitur partem debilio-
rem"). If, therefore,
one of the premises is a negative proposition, then the conclusion must be
negative too, and
if one premise is a particular proposition, then the conclusion must be particu-
lar too.
In all categorical inferences of reason the majo r L must be a universal proposi-
tion (universalis), the minor L an affirmative one (affirmans), and from this it
follows, finally,
From purely negative [premises] nothing follows.
"besondere (particulare)."
From purely particular [premises] nothing follows.
The conclusion follows the weaker part.
618THE J Ä S C H E LOGIC
that the conclusion must follow the major premise in regard to quality, but the
minor premise as to quantity.
Note.
It is easy to see that the conclusion must in every case follow the negative
and the particular proposition in the premises.
If I make the minor premise only particular, and say, Some is contained
under the rule, then I can only say in the conclusion that the predicate of
the rule belongs to some, because I have not subsumed more than this under
the rule. And if I have a negative proposition as the rule (major premise),
then I must also make the conclusion negative. For if the major premise
says that this or that predicate must be denied of everything that stands
under the condition of the rule, then the conclusion must also deny the
predicate of that (subject) which has been subsumed under the condition
of the rule.
§65
Pure and mixed categorical inferences of reason
A categorical inference of reason is pure 0 (punts) if no immediate inference
is mixed in with it, nor is the legitimate order of the premises altered; in
the contrary case it is called impure or mixed (ratiocinium impurum or
hybridum).
§66
Mixed inferences of reason through conversion of
propositions - Figures
Those inferences which arise through the conversion of propositions and
in which the position of these propositions is thus not the legitimate one
are to be counted as mixed inferences. This case occurs in the three latter
so-called figures of the categorical inference of reason.
§6?
Four figures of inferences
By figures are to be understood those four modes of inferring whose
difference is determined through the particular position of the premises
and their concepts.
Ak, "rein"; ist ed., "rein."
619
§68
Ground of the determination of their distinction through the
position of the middle concept
126
different
The middle concept, namely, on whose position things really depend
here, can occupy either (i.) the place of the subject in the major premise
and the place of the predicate in the minor premise, or (2.) the place of the
predicate in both premises, or (3.) the place of the subject in both, or
finally (4.) the place of the predicate in the major premise and the place of
the subject in the minor premise. Through these four cases, the distinc-
tion among the four figures is determined. If we let 5 signify the subject of
the conclusion, P its predicate, and M the terminus medius, then the
schema for the four mentioned figures may be exhibited in the following
table:
M
S
S
P
M
P
P M
S M
S P
M P
M S
S P
P M
M S
S P
§69
Rules for the first figure, as the only legitimate one
The rule of \h& first figure is that the major^ be a universal proposition, the
minor L an affirmative. And since this must be the universal rule of all
categorical inferences of reason in general, it results from this that the
first figure is the only legitimate one, which lies at the basis of all the
others, and to which all the others must be traced back through conver-
sion of premises (metathesis praemissorum), insofar as they are to have
validity.
Note.
The first figure can have a conclusion of any quantity and quality. In the
other figures there are conclusions only of a certain kind; some modi of
them are excluded here. This already indicates that these figures are not
perfect, that there are certain restrictions in them which prevent the conclu-
sion from occurring in all modi, as in the first figure.
§70
Condition of the reduction of the three latter figures to the first
The condition of the validity of the three latter figures, under which a
correct modus of inference is possible in each of them, amounts to this:
that the medius terminus occupies a place such that from it, through
immediate inferences (consequentias immediatas), their position in accor-
dance with the rules of the first figure can arise. - From this the following
rules for the three latter figures emerge.
Rules of the second figure
In the second figure the minor L stands rightly, hence the major L must be
converted, and in such a way that it remains universal (universalis). This is
possible only if it is universally negative; but if it is affirmative, then it must
be contraposed. In both cases the conclusion becomes negative (sequitur
partem debiliorem).
Note.
The rule of the second figure is, What is contradicted by a mark of a thing
contradicts the thing itself. Now here I must first convert and say, What is
contradicted by a mark contradicts this mark; or I must convert the conclu-
sion, What is contradicted by a mark of a thing is contradicted by the thing
itself, consequently it contradicts the thing.
§72
Rule of the third figure
In the third figure the major L stands rightly, hence the minor L must be
converted, yet in such a way that an affirmative proposition arises there-
from. This is only possible, however, when the affirmative proposition is
particular, consequently the conclusion is particular.
Note.
The rule of the third figure is, What belongs to or contradicts a mark also
belongs to or contradicts some things under which this mark is contained.
Here I must first say, It belongs to or contradicts everything that is con-
tained under this mark.
621
128
§73
Rule of the fourth figure
If in the fourth figure the major L is universal negative, then it may be
converted purely (simplidter), just as the minor L may be as particular;
hence the conclusion is negative. If, on the other hand, the major L is
universal affirmative, then it may either be converted only per accidens or it
may be contraposed; hence the conclusion is either particular or negative.
If the conclusion is not to be converted (P S transformed into 5 P), a
transposition of the premises (metathesis praemissorum) or a conversion
(conversio) must occur.
Note.
In the fourth figure it is inferred that the predicate depends on the medius
terminus, the medius terminus on the subject (of the conclusion), conse-
quently the subject on the predicate; this simply does not follow, however, but
at most its converse. To make this possible, the major L must be made the
minor L and vice versa, and the conclusion must be converted, because in the
first alteration the terminus minor is transformed into the major L .
§74
Universal results concerning the three latter figures
From the rules stated for the three latter figures it is clear
1.
2.
3.
that in none of them is there a universal affirmative conclusion, but rather the
conclusion is always either negative or particular;
that in every one an immediate inference (consequentia immediatd) is mixed in,
which is not expressly signified, to be sure, but still must be silently included;
that on this account, too,
these three latter modi of inference must all be called impure inferences
(ratiocinia hybrida, impurd), not pure ones, since no pure inference can have
more than three principal propositions (termini).
129
§75
2. Hypothetical inferences of reason
A hypothetical inference is one that has a hypothetical proposition as
major L . Thus it consists of two propositions, (i.) an antecedent prop-
osition p (antecedent) and (2.) a consequent proposition 11 (consequens), and
f
q
"Vordersätze.™
"Nachsatze."
622THE J Ä S C H E LOGIC
here the deduction is either according to modus ponens or to modus
tollens.
Note i. Hypothetical inferences of reason have no medius terminus, then, but
instead the consequentia of one proposition from another is only indicated
in them. In their major L , namely, the consequentia of two propositions from
one another is expressed, of which the first is the premise, the second a
conclusion. The minor L is a transformation of the problematic condition
into a categorical proposition.
2. From the fact that the hypothetical inference consists only of two proposi-
tions, without having a middle concept, it may be seen that it is really not
an inference of reason, but rather only an immediate inference, to be
proved from an antecedent proposition and a consequent proposition, as
to matter or form (consequentia immediata demonstrabilis [ex antecedente et
consequente] vel quoad materiam vel quoad formam).
Every inference of reason is supposed to be a proof. But the hypotheti-
cal carries with it only the ground of proof. It is clear from this, conse-
quently, that it cannot be an inference of reason.
§76
Principle of hypothetical inferences
The principle of hypothetical inferences is the principle of the ground-/ a
ratione ad rationatum; a negatione rationati ad negationem rationis valet conse-
quential
§77
3. Disjunctive inferences of reason
In disjunctive inferences, the major L is a disjunctive proposition and as such
must therefore have members of division or disjunction.
Here we infer either (i.) from the truth of one member of the
disjunction to the falsehood of the others, or (2.) from the falsehood of all
members but one to the truth of this one. The former occurs through the
modus ponens (or ponendo tollens), the latter through the modus tollens
(tollendo ponens).
'
"Satz des Grundes."
'
The consequentia from the ground to the grounded, and from the negation of the
grounded to negation of the ground, is valid.
623
130IMMANUEL KANT
Note i. All the members of the disjunction but one, taken together, constitute the
contradictory opposite of this one. Here there is a dichotomy, then,
according to which, if one of the two is true, the other must be false, and
conversely.
2. All disjunctive inferences of reason of more than two members of
disjunction are thus really polysyllogistic. For every true disjunction can
only be bimembris, and logical division is also bimembris, but for brevity's
sake the membra subdividentia are posited among the membra dividentia.
§78
Principle of disjunctive inferences of reason
The principle of disjunctive inferences is the principle of the excluded mid-
dle:'
A contradictorie oppositorum negatione unius ad affirmationem
positione unius ad negationem alterius valet consequential
alterius, a
§79
Dilemma
A dilemma is a hypothetical-disjunctive inference of reason, or a hypo-
thetical inference, whose consequens is a disjunctive judgment. The hypo-
thetical proposition whose consequens is disjunctive is the major proposi-
tion; the minor proposition affirms that the consequens (per omnia membra)
is false, and the conclusion affirms that the antecedent is false. (A remotione
consequents ad negationem antecedents valet consequential)
Note.
131
'
The ancients made a great deal of the dilemma and named this inference
comutus. They knew how to drive an opponent into a corner by rehearsing
everything to which he could turn and then contradicting everything, too.
They showed him many difficulties with any opinion he accepted. But it is
a sophistical trick, not to refute propositions directly but rather only to
show difficulties, which is feasible with many, indeed, most things.
Now if we wish to declare as false everything with which we find difficul-
ties, then it is an easy game to reject everything. It is good, to be sure, to
show the impossibility of the opposite, only there still lies something decep-
"'Grundatz des ausschließenden Dritten.™
" From the negation of one contradictory opposite to the affirmation of the other, and from
the positing of one to the negation of the other, the consequentia is valid.
1
From the denial of the consequence to the negation of the antecedent the consequence is
valid.
tive in this in case one takes the incomprehensibility of the opposite for its
impossibility. Hence dilemmata have much that is captious in them, even
though they infer correctly. They can be used to defend true propositions,
but also to attack true propositions by means of difficulties that one throws
up against them.
§80
Formal and coven inferences of reason (ratiocinia formalia and
cryptica)
A formal inference of reason is one that not only contains everything
required as to matter but also is expressed correctly and completely as to
form. Opposed to formal inferences of reason are coven ones (cryptica), to
which all those can be reckoned, in which either the premises are trans-
posed, or one of the premises is left out, or, finally, the middle concept
alone is combined with the conclusion. A covert inference of reason of the
second kind, in which one premise is not expressed but only thought, is
called a truncated" one or an enthymeme. Those of the third kind are called
contracted" inferences.
I I I . I N F E R E N C E S OF THE POWER OF J U D G M E N T
§81
Determinative and reflective power of judgment
The power of judgment is of two kinds: the determinative 11 or the reflective 7 '
power of judgment. The former goes from the universal to the particular, the
second from the particular to the universal. The latter has only subjective
validity, for the universal to which it proceeds from the particular is only
empirical universality - a mere analogue of the logical.
§82
Inferences of the (reflective) power of judgment
Inferences of the power of judgment are certain modes of inference for
coming from particular concepts to universal ones. They are not functions
"verstümmelter."
"zusammengezogene."
"bestimmende."
"reflectirende."
625
of the determinative power of judgment, then, but rather of the reflective;
hence they also do not determine the object, but only the mode of reflection
concerning it, in order to attain its cognition.
§83
Principle of these inferences
The principle that lies at the basis of these inferences of the power of
judgment is this: that the many will not agree in one without a common ground,
but rather that which belongs to the many in this way will be necessary due to a
common ground.
Note.
Since such a principle lies at the basis of the inferences of the power of
judgment, they cannot, on that account, be held to be immediate inferences.
§84
Induction and analogy — The two modes of inference of the power of
judgment
The power of judgment, by proceeding from the particular to the univer-
sal in order to draw from experience (empirically) universal - hence not a
priori — judgments, infers either from many to all things of a kind, or from
many determinations and properties, in which things of one kind agree, to
the remaining ones, insofar as they belong to the same principle. The former
mode of inference is called inference through induction, the other infer-
ence according to analogy.
133
Note i. Induction infers, then, from the particular to the universal (aparticular? ad
universale) according to the principle of universalization:* What belongs to
many things of a genus belongs to the remaining ones too. Analogy infers from
particular to total similarity of two things, according to the principle of
specification: Things of one genus, which we know to agree in much/ also
agree in what remains, with which we are familiar in some things of this
genus but which we do not perceive in others. Induction extends the
empirically given from the particular to the universal in regard to many
objects, while analogy extends the given properties of one thing to several
[other properties] of the very same thing[.] - One in many, hence in all:
Induction; many in one (which are also in others), hence also what remains
* "Princip der Allgemeinmachung."
* "von denen man vieles Uebereinstimmende kennt."
626THE J Ä S C H E LOGIC
in the same thing: Analogy. Thus the ground of proof for immortality
from the complete .development of natural dispositions of each creature
is, for example, an inference according to analogy.
In the inference according to analogy, however, identity of the ground
(par ratio) is not required. In accordance with analogy we infer only
rational inhabitants of the moon, not men. Also, one cannot infer accord-
ing to analogy beyond the tertium comparationis.
Every inference of reason must yield necessity. Induction and analogy are
therefore not inferences of reason, but only logical presumptions, or even
empirical inferences; and through induction one does get general proposi-
tions, but not universal ones.
The mentioned inferences of the power of judgment are useful and
indispensable for the sake of the extending of our cognition by experi-
ence/ But since they give only empirical certainty, we must make use of
them with caution and care.
§85
Simple and composite inferences of reason
An inference of reason is called simple if it consists of one inference of
reason, composite if of several.
§86
Ratiocinatio polysyllogistica
A composite inference, in which the several inferences of reason are
combined with one another not through mere coordination but through
subordination, i.e., as grounds and consequences, is called a chain of
inferences of reason (ratiodnatio polysyllogistica).
§8?
Prosyllogisms and episyllogisms
In the series of composite inferences one can infer in two ways, either
from the grounds down to the consequences, or from the consequences
up to the grounds. The first occurs through episyllogisms, the other
through prosyllogisms.
An episyllogism is that inference, namely, in the series of inferences,
whose premise becomes the conclusion of a prosyllogism, hence of an
inference that has the premises of the former as conclusion.
'
"Erfahrungserkenntnisses."
627
§88
Sorites or chain inference
An inference consisting of several inferences that are shortened and com-
bined with one another for one conclusion is called a sorites or a chain
inference, which can be either progressive or regressive, accordingly as one
climbs from the nearer grounds up to the more distant ones, or from the
more distant grounds down to the nearer ones.
§8 9
Categorical and hypothetical sorites
Both progressive and regressive chain inferences can in turn be either
categorical or hypothetical. The former consists of categorical propositions, as a
series of predicates, the latter of hypothetical ones, as a series of conse-
quences.
§90
Fallacy — Paralogism - Sophism
135
An inference of reason that is wrong as to form, although it has for itself
the illusion of a correct inference, is called a fallacy {fallacia). Such an
inference is a paralogism insofar as one deceives oneself through it, a
sophism insofar as one intentionally seeks to deceive others through it.
Note.
The ancients occupied themselves very much with the art of making such
sophisms. Therefore many of this kind have emerged, e.g., the sophisma
figurae dictionis, in which the medius terminus is taken in different
meanings - fallacia a dicta secundum quid ad dictum simpliciter, sophisma
heterozeteseos, elenchi d ignorationis,' etc.
§9!
Leap in inference
A leap (saltus) in inference or proof is the combination of a premise with the
conclusion so that the other premise is left out. Such a leap is legitimate
d
Reading "elenchi" for "elenchi,".
' the fallacy [of inferring] from what is said with qualification to what is said simpliciter, the
sophism of misdirection, the ignoratio elenchis, etc.
(legitimus) if everyone can easily add the missing premise in thought, but
illegitimate (illegitimus) if the subsumption is not clear. Here a distant mark
is connected with a thing without an intermediate mark (nota intermedia).
§92
Petitio principii - Circulus in p r o b a n d o
By a petitio principii is understood the acceptance of a proposition as
ground of proof as an immediately certain proposition, although it still
requires a proof. And one commits a circle in proof if one lays at the basis of
its own proof the very proposition that one wanted to prove.
Note.
The circle in proof is often hard to discover, and this mistake is usually
most frequently committed where the proofs are hard.
§93
Probatio plus and minus probans
A proof can prove too much, but also too little. In the latter case it proves
only a part of what is to be proved, in the former it also goes on to that
which is false.
Note.
A proof that proves too little can be true and thus is not to be rejected. But
if it proves too much, then it proves more than is true, and that is then
false. Thus, for example, the proof against suicide, that he who has not
given life cannot take it either, proves too much; for on this ground we
would not be permitted to kill animals. Thus it is false.
629
136//.
Universal doctrine of method
§94
Manner and method
All cognition, and a whole of cognition, must be in conformity with a rule.
(Absence of rules is at the same time unreason.) But this rule is either that
of manner (free) or that of method (compulsion).
§95
Form of science — Method
Cognition, as science, must be arranged in accordance with a method. For
science is a whole of cognition as a system, and not merely as an aggre-
gate. It therefore requires a systematic cognition, hence one composed in
accordance with rules on which we have reflected.
§96
Doctrine of method — Its object and end
As the doctrine of elements in logic has for its content the elements and
conditions of the perfection of a cognition, so the universal doctrine of
method, as the other part of logic, has to deal with the form of a science in
general, or with the ways of acting so as to connect the manifold of
cognition in a science.
§97
Means for furthering the logical perfection of cognition
140
The doctrine of method is supposed to expound the way for us to attain
the perfection of cognition. Now one of the most essential logical perfec-
tions of cognition consists in its distinctness, thoroughness, and system-
atic ordering into the whole of a science. Accordingly, the doctrine of
method will have principally to provide the means through which these
perfections of cognition are furthered.
§98
Conditions of the distinctness of cognition
The distinctness of cognitions and their combination in a systematic
whole depends on the distinctness of concepts both in regard to what is
contained in them and in respect of what is contained under them.
The distinct consciousness of the content of concepts is furthered by
exposition and definition of them, while the distinct consciousness of their
extension, on the other hand, is furthered through logical division of them.
— First of all, then, of the means for furthering the distinctness of con-
cepts in regard to their content.
I. F U R T H E R I N G L O G I C A L P E R F E C T I O N OF
COGNITION THROUGH DEFINITION,
E X P O S I T I O N , AND D E S C R I P T I O N OF C O N C E P T S
§99
Definition
A definition is a sufficiently distinct and precise concept (conceptus rei
adaequatus in minimis terminis, complete determinatus f ).
Note.
The definition alone is to be regarded as a logically perfect concept, for in
it are united the two essential perfections of a concept: distinctness, and
completeness and precision in distinctness (quantity of distinctness).
§100
Analytic and synthetic definition
All definitions are either analytic or synthetic. The former are definitions
of a concept that is given, the latter of one that is made.
A concept adequate to the thing, in minimal terms, completely determined.
631
§IOI
Concepts that are given and made a priori and a posteriori
The given concepts of an analytic definition are given either a priori or a
posteriori, just as the concepts of a synthetic definition, which are made,
are made either a priori or a posteriori.
§102
Synthetic definitions through exposition or construction
The synthesis of concepts that are made, out of which synthetic defini-
tions arise, is either that of e3Cposition s (of appearances) or that of construc-
tion. The latter is the synthesis of concepts that are made arbitrarily, the
former the synthesis of concepts that are made empirically, i.e., from given
appearances as their matter (conceptus factitii vel a priori vel per synthesin
empiricam 1 '). Concepts that are made arbitrarily are the mathematical ones.
Note.
All definitions of mathematical concepts and - provided that definitions
could exist in the case of empirical concepts - of concepts of experience
must be made synthetically, then. For even in the case of concepts of the
latter kind, e.g., the empirical concepts water, fire, air, etc., I ought not to
analyze what lies in them, but to become acquainted, through experience,
with what belongs to them. All empirical concepts must thus be regarded as
concepts that are made, whose synthesis is not arbitrary, however, but
empirical
§103
Impossibility of empirically synthetic definitions
142
Since the synthesis of empirical concepts is not arbitrary but rather is
empirical and as such can never be complete (because one can always
discover in experience more marks of the concept), empirical concepts
cannot be defined, either.
Note.
*
h
Thus only arbitrary concepts may be defined synthetically. Such defini-
tions of arbitrary concepts, which are not only always possible but also
necessary, and which must precede all that is said by means of an arbitrary
"Exposition."
concepts that are made either a priori or through empirical synthesis.
632THE J Ä S C H E LOGIC
concept, could also be called declarations, insofar as through them one
declares his thoughts or gives account of what one understands by a word.
This is the case among mathematicians.
§104
Analytic definitions through analysis of concepts given a p r i o r i or a
posteriori
All given concepts, be they given a priori or a posteriori, can be defined only
through analysis. For one can make given concepts distinct only insofar as
one successively makes their marks clear. If all the marks of a given
concept are made clear, then the concept becomes completely distinct; if it
does not contain too many marks, then it is at the same time precise, and
from this there arises a definition of the concept.
Note.
Since one cannot become certain through any test whether one has ex-
hausted all the marks of a given concept through a complete analysis, all
analytic definitions are to be held to be uncertain.
§105
Expositions and descriptions
Not all concepts can be defined, and not all need to be.
There are approximations to the definition of certain concepts; these
are partly expositions' (expositiones), partly descriptions (descriptions).
The expounding' of a concept consists in the connected (successive)
representation* of its marks, insofar as these are found through analysis.
Description is the exposition of a concept, insofar as it is not precise.
Note i. We can expound either a concept or experience. The first occurs through
analysis, the second through synthesis.
2. Exposition occurs only with given concepts, then, which are thereby made
distinct; it is thereby distinct from declaration, which is a distinct represen-
tation of concepts that are made.
Since it is not always possible to make analysis complete, and since in
general an analysis must first be incomplete before it becomes complete,
an incomplete exposition, as part of a definition, is also a true and useful
exhibition of a concept. Definition always remains here only the idea of a
logical perfection that we must seek to attain.
j. Description can occur only with empirically given concepts. It has no
determinate rules and contains only the materials for definition.
§106
Nominal and real definitions
By mere definitions of names, 1 or nominal definitions, are to be understood
those that contain the meaning that one wanted arbitrarily to give to a
certain name, and which therefore signify only the logical essence of their
object, or which serve merely for distinguishing it from other objects.
Definitions of things," 1 or real definitions, on the other hand, are ones that
suffice for cognition of the object according to its inner determinations,
since they present the possibility of the object from inner marks.
144
Note i. If a concept is internally sufficent for distinguishing the thing then it
certainly is externally sufficient too, but if it is not internally sufficient
then it can be externally sufficient merely in a certain relation, namely, in
the comparison of the definitum with other things. But unrestricted exter-
nal sufficiency is not possible without internal sufficiency.
2. Objects of experience allow only nominal explanations. Logical nominal
definitions of given concepts of the understanding are derived from an
attribute, real definitions, on the other hand, from the essence of the
thing, the first ground of possibility. Thus the latter contain what always
belongs to die thing - its real essence. Merely negative definitions cannot
be called real definitions either, because negative marks can serve just as
well as affirmative ones for distinguishing one thing from others, but not
for cognition of the thing according to its inner possibility.
In matters of morals real definitions must always be sought; all our
striving must be directed toward this. There are real definitions in mathe-
matics, for the definition of an arbitrary concept is always real.
j. A definition is genetic if it yields a concept dirough which the object can be
exhibited a priori in concreto; all mathematical definitions are of this sort.
§107
Principal requirements of definition
The essential and universal requirements that pertain to the completeness
of a definition in general may be considered under the four principal
moments of quantity, quality, relation, and modality:
' "Namen-Erklärungen."
m "Sach-Erklarungen."
as to quantity - what concerns the sphere of the definition - the definition
and the definitum must be convertible concepts" (conceptus redprocf), and hence
the definition must be neither broader nor narrower than its definitum;
as to quality, the definition must be a detailed and at the same time precise
concept,
as to relation, it must not be tautological, i.e., the marks of the definitum must,
as grounds of its cognition, be different from it itself, and finally
as to modality, the marks must be necessary, and hence not such as are added
through experience.
Note.
The condition that the genus concept and the concept of the specific
difference (genus and differentia specified) ought to constitute the definition
holds only in regard to nominal definitions in comparison, but not for real
definitions in derivation.
§108
Rules for testing definitions
In the testing of definitions four acts are to be performed; it is to be
investigated, namely, whether the definition
1.
2.
3.
4.
considered as a proposition is true, whether
as a concept it is distinct,
whether as a distinct concept it is also detailed, and finally whether
as a detailed concept it is at the same time determinate, i.e., adequate to the
thing itself.
§109
Rules for preparation of definitions
The very same acts that belong to the testing of definition are also to be
performed in the preparation of them. Toward this end, then, seek (i.)
true propositions, (2.) whose predicate does not presuppose the concept
of the diing; (3.) collect several of them and compare them with the
concept of the thing itself to see if they are adequate; and finally (4.) see
whether one mark does not lie in another or is not subordinated to it.
Note i.
"
These rules hold, as is surely understood without any reminder, only of
analytic definitions. But now since one can never be certain here whether
the analysis has been complete, one may put forth the definition only as
" Wechselbegriffe."
635
145IMMANUEL KANT
an experiment and make use of it only as if it were not a definition. Under
this restriction one can still use it as a distinct and true concept and draw
corollaries from the concept's marks. I will be able to say, namely, that to
that to which the concept of the definitum belongs, the definition belongs
too, but of course not conversely, since the definition does not exhaust
the whole definitum.
To make use of the concept of the definitum in the explanation, or to
make the definitum the basis of the definition, is called explaining through
a circle (drculus in definiendo).
146
II. FURTHERING THE PERFECTION
OF COGNITION THROUGH LOGICAL DIVISION
OF CONCEPTS
§110
Concept of logical division
Every concept contains a manifold under itself insofar as the manifold
agrees, but also insofar as it is different. The determination of a concept
in regard to everything possible that is contained under it, insofar as
things are opposed to one another, i.e., are distinct from one another, is
called the logical division of the concept. The higher concept is called the
divided concept (divisus), the lower concepts the members of the division
(membra dividentia).
Note i. To take apart' a concept and to divide? it are thus quite different things. In
taking a concept apart I see what is contained in it (through analysis), in
dividing it I consider what is contained under it. Here I divide the sphere
of the concept, not the concept itself. Thus it is a great mistake to
suppose that division is the taking apart of the concept; rather, the mem-
bers of the division contain more in themselves than does the divided
concept.
2. We go up from lower to higher concepts, and afterward we can go down
from these to the lower ones - through division.
§in
Universal rules of logical division
In every division of a concept we must see to it:
'
f
"theilen."
"eintheilen."
1.
2.
3.
that the members of the division exclude or are opposed to one another, that
furthermore they
belong under one higher concept (conceptus communis), and finally that
taken together they constitute the sphere of the divided concept or are equal
to it.
Note.
The members of the division must be separated from one another through
contradictory opposition, not through mere contrariety* (contrarium).
§112
Codivision and subdivision
Various divisions of a concept, which are made in various respects, are
called codivisions, and division of the members of division is called a
subdivision (mbdivisio).
Note i.
2.
Subdivision can be carried on to infinity, but comparatively it can be
finite. Codivision, especially with concepts of experience, always goes to
infinity; for who can exhaust all the relations of concepts?
Codivision can also be called a division according to the variety of con-
cepts of the same object (viewpoints), as subdivision can be called a
division of the viewpoints themselves.
§"3
Dichotomy and polytomy
A division into two members is called dichotomy; but if it has more than two
members, it is called polytomy.
Note i. All polytomy is empirical; dichotomy is the only division from principles a
priori, hence the only primitive division. For the members of a division are
supposed to be opposed to one another, but for each A the opposite is
nothing more than nonA.
2. Polytomy cannot be taught in logic, for it involves cognition of the object.
Dichotomy requires only the principle of contradiction, however, without
being acquainted, as to content, with the concept one wants to divide.
Polytomy requires intuition, either a priori, as in mathematics (e.g., in the
division of conic sections), or empirical intuition, as in description of
nature. However, division based on the principle of synthesis' a priori has
* "Widerspiel."
' "Princip der Synthesis."
637
147IMMANUEL KANT
148
trichotomy, namely: (i.) the concept as condition, (2.) the conditioned,
and (3.) the derivation of the latter from the former.
§"4
Various divisions of method
Now as for what concerns in particular method itself in working up and
treating scientific cognitions, there are various principal kinds of it, which
we can present in accordance with the following division.
i. Scientific or popular method
1
Scientific or scholastic method differs from popular method through the
fact that the former proceeds from basic and elementary propositions,
but the latter from the customary and the interesting. The former aims
for thoroughness and thus removes everything foreign, the latter aims at
entertainment.
Note.
These two methods differ as to kind, then, and not merely as to exposition,
and popularity in method is hence something other than popularity in
exposition.
§n6
2. Systematic or fragmentary
method
Systematic method is opposed to fragmentary or rhapsodic method. If one
has thought in accordance with a method and then also expressed this
method in the exposition, and if the transition from one proposition to
another is distinctly presented, then one has treated a cognition systemati-
cally. If, on the other hand, one has thought according to a method but has
not arranged the exposition methodically, such a method is to be called
rhapsodic.
Note.
149
Systematic exposition is opposed to the fragmentary as methodical exposition
is to the tumultuous. He who thinks methodically can expound systemati-
"scientifische."
cally or fragmentarily, that is. Exposition that is outwardly fragmentary but
in itself methodical is aphoristic.
§"7
3. Analytic or synthetic method
Analytic is opposed to synthetic method. The former begins with the condi-
tioned and grounded and proceeds to principles (a prindpiatis ad prin-
dpia), while the latter goes from principles to consequences or from the
simple to the composite. The former could also be called regressive, as the
latter could progressive.
Note. Analytic method is also called the method of invention. Analytic method is
more appropriate for the end of popularity, synthetic method for the end of
scientific and systematic preparation of cognition.
§118
4. Syllogistic [or] tabular method
Syllogistic method is that according to which a science is expounded in a
chain of inferences.
That method in accordance with which a finished system is exhibited in
its complete connection is called tabular.
§"9
Sf.Acroamatic or erotematic method
Method is acroamatic insofar as someone only teaches, erotematic insofar as
he asks as well. The latter method can be divided in turn into dialogic or
Socratic method and catechistic method, accordingly as the questions are
directed either to the understanding or merely to memory.
Note.
One cannot teach erotematically except through the Socratic dialogue, in
which both must ask and also answer one another in turn, so that it seems
as if the pupil is also himself teacher. The Socratic dialogue teaches, that
is, through questions, by acquainting the learner with his own principles of
reason and sharpening his attention to them. Through common catechism,
however, one cannot teach but can only elicit what one has taught
acroamatically. Catechistic method also holds, then, only for empirical and
historical cognitions, while the dialogic holds for rational cognitions.
639
§120
Meditation
By meditation is to be understood reflection, or methodical thought. Medi-
tation must accompany all reading and learning, and for this it is requisite
that one first undertake provisional investigations and then put his
thoughts in order, or connect them in accordance with a method.